% ==============================================================================
% CEC 320 FINAL EXAM CHEATSHEET --- COMBINED PREAMBLE
% Auto-generated header for single-file 14-page output.
% Concatenate with all partials in part1/ then part2/ then _footer.tex.
% Built by build.sh.
% ==============================================================================
\documentclass[8pt,landscape]{extarticle}
\usepackage[letterpaper,landscape,margin=0.18in,top=0.18in,bottom=0.18in]{geometry}
\usepackage{multicol}
\usepackage{amsmath,amssymb}
\usepackage{array}
\usepackage{booktabs}
\usepackage{xcolor}
\usepackage{listings}
\usepackage{enumitem}
\usepackage{titlesec}
\usepackage[T1]{fontenc}
\usepackage{lmodern}
\usepackage{microtype}

\AtBeginDocument{\small}

% ---------- Spacing / layout ----------
\setlength{\columnsep}{0.3em}
\setlength{\columnseprule}{0.15pt}
\setlength{\parindent}{0pt}
\setlength{\parskip}{1pt}
\setlist{nosep,leftmargin=0.9em,itemsep=0pt,topsep=0.5pt,parsep=0pt,partopsep=0pt}
\setlength{\tabcolsep}{2pt}
\renewcommand{\arraystretch}{1.0}
\pagestyle{empty}
\raggedcolumns

% ---------- Section styling (tight) ----------
\titleformat{\section}{\normalfont\bfseries\footnotesize\color{blue!60!black}}
  {\thesection.}{0.25em}{}
\titlespacing*{\section}{0pt}{3pt}{0.5pt}
\titleformat{\subsection}{\normalfont\bfseries\scriptsize}{}{0em}{}
\titlespacing*{\subsection}{0pt}{2pt}{0.2pt}

% ---------- ARM listings ----------
\lstdefinelanguage{arm}{
  morekeywords={ldr,ldrb,ldrh,ldrsb,ldrsh,ldrd,str,strb,strh,strd,mov,movw,movt,%
    movs,add,adds,adc,sub,subs,sbc,cmp,cmn,tst,teq,bl,bx,b,beq,bne,bgt,bge,blt,%
    ble,bhi,bhs,blo,bls,bcs,bcc,push,pop,cbz,cbnz,it,itt,ite,itte,itet,itee,%
    ittt,ittee,iteee,ittte,itttt,tbb,tbh,adr,lsl,lsls,lsr,lsrs,asr,asrs,ror,%
    rors,rrx,rrxs,and,ands,orr,orrs,eor,eors,bic,bics,orn,mvn,mvns,bfc,bfi,%
    mul,muls,mla,mls,umull,umlal,smull,smlal,udiv,sdiv,nop,wfi,wfe,svc,rsb,%
    sxtb,sxth,uxtb,uxth,clz,rbit,rev,%
    lslge,lslgt,lsllt,lslle,asrge,asrgt,asrlt,asrle,ldrle,ldrgt,ldrlt,ldrge,%
    ldrshgt,ldrshlt,ldrhge,ldrhlt,lsrlt,lsrge},
  morecomment=[l]{@},
  sensitive=false
}
\lstset{
  language=arm,
  basicstyle=\ttfamily\scriptsize,
  keywordstyle=\color{blue!70!black}\bfseries,
  commentstyle=\color{green!40!black}\itshape,
  breaklines=false,
  frame=none,
  aboveskip=0.5pt,
  belowskip=0.5pt,
  columns=fullflexible,
  keepspaces=true,
  xleftmargin=0.25em,
  showstringspaces=false
}
\newcommand{\asm}[1]{\texttt{\scriptsize #1}}
\newcommand{\cee}[1]{\texttt{\scriptsize #1}}

% Tag environments
\newenvironment{ex}
  {\par\smallskip\noindent{\scriptsize\color{red!55!black}\textbf{Ex.}}\,\scriptsize\ignorespaces}
  {\par}
\newenvironment{qp}
  {\par\smallskip\noindent{\scriptsize\color{purple!70!black}\textbf{Qz.}}\,\scriptsize\ignorespaces}
  {\par}
\newenvironment{hwp}
  {\par\smallskip\noindent{\scriptsize\color{orange!70!black}\textbf{HW}}\,\scriptsize\ignorespaces}
  {\par}
\newenvironment{reg}
  {\par\smallskip\noindent{\scriptsize\color{teal!75!black}\textbf{Reg}}\,\scriptsize\ignorespaces}
  {\par}

\begin{document}
\begin{center}
  \textbf{\footnotesize CEC 320 --- FINAL EXAM Cheatsheet} \;
  \scriptsize HW 1--12 \textbullet{} Lctr 2--9, 13--27 \textbullet{} Quiz 1--4, 6a \textbullet{} STM32 Peripheral Reference
\end{center}
\vspace{-8pt}

\begin{multicols*}{4}
\scriptsize

% ----- part1/01_lectures_14_27.tex -----
\section{Lctr 14--27 Recall Bullets}

\subsection{L14 --- Mixed C/ASM (EABI)}
\begin{itemize}
\item EABI: \asm{R0--R3} pass args + return; caller-saved (NOT preserved). \asm{R4--R11} callee-saved $\Rightarrow$ push/pop if used.
\item 64-bit arg/return: low word in \asm{R0}, high word in \asm{R1} (little-endian).
\item Stack = Full Descending: SP points at last pushed item; \asm{push} decrements SP by 4/word.
\item Multi-reg push: low-numbered regs end up at LOW addresses regardless of brace order in \asm{\{...\}}.
\item Return: \asm{bx lr} if LR not stacked; \asm{pop \{pc\}} if LR was pushed ($\equiv$\asm{pop \{lr\}}+\asm{bx lr}).
\item Leaf fn $\Rightarrow$ no LR push needed; stem fn (calls another) MUST push LR before BL.
\end{itemize}

\subsection{L15 --- Arithmetic}
\begin{itemize}
\item \asm{ADD/SUB/RSB} (no S) leave flags alone; \asm{ADDS/SUBS/RSBS} update NZCV. \asm{RSB} exists because Rn must be a reg.
\item \asm{ADC Rn,Op2}$=$Rn$+$Op2$+$C; \asm{SBC Rn,Op2}$=$Rn$-$Op2$+$C$-$1 (C$=$1 means no-borrow).
\item 64-bit add idiom: \asm{adds r0,r2} / \asm{adc r1,r3} (low sets C, high consumes it).
\item \asm{MUL/MLA/MLS} 32-bit signed$\equiv$unsigned in low 32; \asm{UMULL/SMULL/UMLAL/SMLAL} for 64-bit.
\item \asm{UDIV} vs \asm{SDIV} differ; no flag-setting form, non-constant cycles.
\item Modulo idiom: \asm{udiv r2,r0,r1} / \asm{mls r0,r1,r2,r0} ($r0=r0-r1\cdot r2$).
\end{itemize}

\subsection{L16 --- Shifts \& Op2}
\begin{itemize}
\item \asm{LSL} both signed/unsigned; \asm{LSR} unsigned-only (0-fill); \asm{ASR} signed-only (sign-bit replicate). Last bit out $\to$ C.
\item \asm{RRX} rotates right by 1 through C: C $\to$ MSB, LSB $\to$ C.
\item Op2 forms: \asm{\#imm}, \asm{Rm}, or \asm{Rm,LSL/LSR/ASR \#n} (barrel shifter, free combined w/ arithmetic).
\item Shifted-reg Op2: destination canNOT be omitted even if same as Rn.
\item Q31 accurate mul (3 instr): \asm{smull r1,r0,r0,r1} / \asm{lsl r0,\#1} / \asm{add r0,r0,r1,lsr \#31}.
\item Fast scaling by $(2^n\pm1)/2^n$: \asm{add r0,r0,r0,asr \#2}$=5/4$; \asm{sub r0,r0,r0,asr \#3}$=7/8$.
\end{itemize}

\subsection{L17 --- Bitwise Logic}
\begin{itemize}
\item \asm{MVN}$=\sim$, \asm{AND}$=\&$, \asm{ORR}$=|$, \asm{EOR}$=\hat{}\,$. No-C-equivalents: \asm{BIC Rn,Op2}$=$Rn\,\&\,$\sim$Op2; \asm{ORN}$=$Rn\,$|$\,$\sim$Op2.
\item \asm{BFC Rd,\#s,\#w} clears $w$ bits at pos $s$; \asm{BFI Rd,Rn,\#s,\#w} inserts low-$w$ bits of Rn at pos $s$.
\item Clear-bit-N: \asm{mov mask,\#1}/\asm{lsl mask,\#N}/\asm{bic r0,mask} (BIC$=$AND-NOT).
\item Read APSR: \asm{MRS r0,apsr}; write: \asm{MSR apsr\_nzcvq,Rn}.
\item Bit-field assign $v$: build mask $((1\!\ll\!w)-1)\!\ll\!s$, \asm{bic} target, shift $v$ to pos, \asm{and} mask, \asm{orr} in.
\end{itemize}

\subsection{L18 --- NZCV Flags}
\begin{itemize}
\item \textbf{N}$=$MSB(result); \textbf{Z}$=$(result$=$0); \textbf{C}$=$unsigned carry/no-borrow; \textbf{V}$=$signed overflow.
\item HW always updates BOTH C and V; programmer chooses interpretation.
\item Unsigned uses \{C,N,Z\}; signed uses \{V,N,Z\}.
\item C on SUB: C$=$1 if no borrow (Rn$\ge$Op2), C$=$0 if borrow (counterintuitive naming).
\item \asm{CMP}$=$SUB w/o writeback (sets flags). \asm{CMN}$=$add-form. \asm{TST}$=$AND-only-flags. \asm{TEQ}$=$EOR-only-flags.
\item N-bit verify trick: \asm{LSL \#(32-N)} operands, \asm{ADDS}, \asm{ASR \#(32-N)} back, read APSR via \asm{MRS}.
\end{itemize}

\subsection{L19 --- CCS \& Conditional Branch}
\begin{itemize}
\item Signed CCS: \asm{EQ/NE/GT/GE/LT/LE}. Unsigned: \asm{EQ/NE/HI/HS(CS)/LO(CC)/LS}.
\item \asm{HS}$=$\asm{CS} (C set $\Rightarrow$ unsigned $\ge$); \asm{LO}$=$\asm{CC} (C clear $\Rightarrow$ unsigned $<$).
\item Mnemonic ``HiLow unSun'': \textbf{Hi/Low} family $=$ unSigned; \textbf{G/L} ($>$/$<$) $=$ Signed.
\item Consistency rule: signed operands $\Rightarrow$ signed CCS + \asm{ASR} for divide; unsigned $\Rightarrow$ unsigned CCS + \asm{LSR}.
\item \asm{cmp r0,\#-5} assembled as \asm{CMN r0,\#5}; signed CCSs still behave correctly.
\end{itemize}

\subsection{L20 --- LDR/STR Immediate}
\begin{itemize}
\item Three forms: \asm{[Ra,\#os]} basic (no update); \asm{[Ra,\#os]!} pre (update BEFORE, ``urgent bang''); \asm{[Ra],\#os} post (update AFTER).
\item LOAD size: none(word), \asm{B}(u8), \asm{SB}(i8), \asm{H}(u16), \asm{SH}(i16). STORE: NO sign variant (truncate only).
\item C$\to$ASM: \asm{*p++}$=$post-index; \asm{*++p}$=$pre-index; \asm{*(p+k)}$=$basic offset, $k\cdot$sizeof(elem).
\item \asm{STRB} writes only low byte of Rv; upper bits ignored (no signed variant needed).
\item Valid offset range: word LDR/STR typically $\pm$4095; byte/halfword $\pm$255 on certain forms.
\end{itemize}

\subsection{L21 --- LDR/STR Register + MOV}
\begin{itemize}
\item Reg-offset: \asm{LDR Rv,[Ra,Ri,LSL \#n]}, $n=\log_2$(elem-size): word$=$2, half$=$1, byte$=$0, dword$=$3.
\item NO writeback form for reg-offset (only immediate-offset has \asm{!} or trailing \asm{,\#os}).
\item Iter idioms: post-index advances Ra; reg-offset keeps Ra constant, increments Ri.
\item \asm{MOV} $=$ small imm only; \asm{MOVW} loads 16b low; \asm{MOVT} loads 16b high (upper halfword).
\item Pseudo \asm{LDR Rd,=expr} $\to$ assembler picks \asm{MOV}/\asm{MOVW}+\asm{MOVT}/PC-relative literal-pool load.
\end{itemize}

\subsection{L22 --- Increment Ops (vars + ptrs)}
\begin{itemize}
\item C precedence: postfix \asm{++} level 1 (LTR, highest); prefix \asm{++},\asm{*},\asm{\&}, cast all level 2 (RTL).
\item \asm{*p++} $=$ deref OLD p, then p advances. \asm{*++p} $=$ p advances, then deref NEW.
\item \asm{++*p} $=$ load val, +1, store back (modifies \emph{value}, not pointer).
\item \asm{++*p++}: post-inc p (so p now points past), but modify mem at OLD addr via \asm{[r0,\#-1]}.
\item Ptr arith scales by sizeof(elem): \asm{p8++}$+1$, \asm{p16++}$+2$, \asm{p32++}$+4$. Diff of typed ptrs $\to$ elements.
\item \asm{B[i]=A[i++]++;} is UB (unsequenced i mods); split or use comma operator.
\end{itemize}

\subsection{L23 --- ASM Calls ASM (LDRD/STRD)}
\begin{itemize}
\item \textbf{Leaf} fn (never BLs): \asm{bx lr} suffices, no LR push.
\item \textbf{Stem} fn (calls another): MUST \asm{push \{lr\}} (BL clobbers LR), \asm{pop \{pc\}} to return.
\item \asm{LDRD Rv\_L,Rv\_H,[Ra,\#os]}: 8 bytes, low word at low addr, high word at high addr (little-endian).
\item Rt2$=$Rt$+1$ (consecutive regs). NO register-offset form for \asm{LDRD/STRD} (basic/pre/post only).
\item Arg-prep by C type: int32$\to$\asm{LDR}; int16$\to$\asm{LDRSH}; uint16$\to$\asm{LDRH}; int8$\to$\asm{LDRSB}; uint8$\to$\asm{LDRB}; int64$\to$\asm{LDRD}.
\end{itemize}

\subsection{L24 --- Combo Branch / CEX (CBZ/IT)}
\begin{itemize}
\item \asm{CBZ Rn,label} / \asm{CBNZ Rn,label}: FORWARD only ($\sim$126 byte range), Rn $\in$ R0--R7, does NOT read/set flags.
\item Perfect after \asm{ldrb} for null-terminator test (loaded byte is flag-neutral).
\item IT block $=$ up to 4 conditional instr sharing one CCS. Forms: \asm{IT, ITT, ITE, ITTT, ITTE, ITEE, ITTEE, ITTTT}.
\item First IT slot is ALWAYS T (matches CCS); later slots T (same) or E (inverse). Each instr inside still needs explicit CCS suffix.
\item HARD bans: \asm{CBZ}/\asm{CBNZ} forbidden inside IT; conditional branch only allowed in LAST IT slot.
\end{itemize}

\subsection{L25 --- If / If-Else (PL/NL/CEX)}
\begin{itemize}
\item \textbf{PL} (positive logic): same-CCS branch to \asm{\_then}; else-block laid FIRST. Reads ``backward'' from C.
\item \textbf{NL} (negative logic): inverse-CCS branch to \asm{\_else}; then-block FIRST. Closer to C source order.
\item if-else-if: chain \asm{cmp}+NL-branch; \asm{b end} at each arm tail to skip the rest.
\item Compound \textbf{OR} (short-circuit): two CMPs each branch to shared \asm{\_then}; fall-through $=$ else.
\item Compound \textbf{AND} (De Morgan): both CMPs branch to \asm{\_else} on FAILURE; fall-through $=$ then.
\item CEX style: small bodies inside IT block instead of branches (saves a label).
\end{itemize}

\subsection{L26 --- Loops}
\begin{itemize}
\item Labels: \asm{<fn>\_lp} (top), \asm{<fn>\_end} (exit).
\item Top-tested: \asm{for}, \asm{while} (0+ iters). Bottom-tested: \asm{do-while} ($\ge$1 iter).
\item \asm{while} top-test form: NL-branch to \asm{\_end} at top, body, update, \asm{b lp}.
\item \asm{while} late-check: jump-over-body on entry, test at bottom $\Rightarrow$ saves one branch per iteration.
\item \asm{while} CEX form: \asm{cmp}+\asm{it cond}+conditional B as last slot in IT (only legal cond-branch-in-IT).
\item Priming-read pattern: \asm{ldrb} before loop + \asm{ldrb} inside loop (eliminated by \asm{while(1)+break}).
\end{itemize}

\subsection{L27 --- break / continue / switch}
\begin{itemize}
\item \asm{break} $=$ forward \asm{b <loop\_end>} (or \asm{CBZ}/\asm{CBNZ} to end on a value test).
\item \asm{continue} in while/do-while $\to$ branch to CONDITION test; in \asm{for} $\to$ branch to UPDATE expr (so \asm{i++} still runs!).
\item \asm{while(1) \{ ...; if(c) break; \}} removes priming-read duplication $\Rightarrow$ single \asm{ldrb} inside loop.
\item Switch naive: \asm{CMP}/\asm{BEQ} ladder; stacked case labels w/ no code between $\Rightarrow$ fall-through.
\item \asm{ADR Rd,label} $=$ PC-relative addr load (position-independent).
\item \asm{TBB [Rn,Rm]}: \asm{PC = PC+4+2*table[Rm]}; table bytes $=$ \asm{(case\_X-first\_case)/2} (Thumb-2 halfword aligned).
\end{itemize}

% ----- part1/02_hw05_qformat.tex -----
\section{HW5: Q-Format Fixed-Point Real Numbers}

\subsection{Format anatomy}
\textbf{Qm.n} (signed): 1 sign + $m$ int + $n$ frac $=$ $N=1+m+n$ bits, two's-comp.
\textbf{UQm.n} (unsigned): $m$ int + $n$ frac $=$ $N=m+n$ bits.
Scale factor $=2^{-n}$. LSB resolution $=2^{-n}$.

\begin{tabular}{@{}lllll@{}}
\toprule
Fmt & $N$ & $n$ & range (signed) & res \\
\midrule
Q3.4   & 8  & 4  & $-8.0\ldots+7.9375$ & $1/16$ \\
UQ3.5  & 8  & 5  & $0.0\ldots 7.96875$ & $1/32$ \\
Q2.5   & 8  & 5  & $-4.0\ldots+3.96875$ & $1/32$ \\
Q15 ($=$Q0.15) & 16 & 15 & $-1.0\ldots+0.99997$ & $2^{-15}$ \\
Q15.16 & 32 & 16 & $\pm 32{\rm k}$ & $2^{-16}$ \\
\bottomrule
\end{tabular}

\subsection{Encode rule (real $\to$ hex code)}
\begin{enumerate}
\item Multiply: $I = f\cdot 2^{n}$
\item Round to nearest integer.
\item If $I<0$ (signed Qm.n): $I \leftarrow I + 2^{N}$ (two's-comp wrap).
\item Express $I$ in hex (width $=\lceil N/4\rceil$ digits).
\end{enumerate}
Out-of-range $\Rightarrow$ saturate or flag overflow.

\subsection{Decode rule (hex code $\to$ real)}
\begin{enumerate}
\item Read $U =$ unsigned int value of code.
\item If signed Qm.n and MSB$=$1: $I = U - 2^{N}$; else $I=U$.
\item $f = I \cdot 2^{-n} = I/2^{n}$.
\end{enumerate}

\subsection{Multiplication / MAC rule}
$\text{Q}m.n \times \text{Q}p.q = \text{Q}(m{+}p{+}1).(n{+}q)$ in $2N$-bit product.
Renormalize back to Q$m.n$ by arithmetic shift right $n$ (\asm{ASR \#n}).
\textbf{Q15 MAC:} \asm{smull} gives Q30 in 64b across (HW:LW); shift the
\emph{low word} right by 15 to recover Q15 product. Sum of Q15 products is
plain integer add (no scaling).

\subsection{Q15.16 result extraction (32-bit Q15.16 from \texttt{smull})}
\asm{smull lo,hi,a,b} gives $a\cdot b$ as Q30.32 in (\asm{hi}:\asm{lo}).
To get Q15.16 (32b): \asm{lsl hi,\#16} (keep low 16 bits of hi as int part),
\asm{lsr lo,\#16} (keep high 16 bits of lo as frac part), then
\asm{orr/add} them: result $= (\text{hi}\!\ll\!16)\,|\,(\text{lo}\!\gg\!16)$.

\begin{lstlisting}
@ Q15 multiply: r0=A (Q15), r1=B (Q15) -> r0=A*B (Q15)
    smull r2, r3, r0, r1   @ r3:r2 = A*B (Q30, 64b)
    lsr   r0, r2, #15      @ shift Q30 -> Q15 (low 32 ok if no overflow)
    bx    lr
\end{lstlisting}

\subsection{HW5 P1 --- Encode 2.75 in UQ3.5 (unsigned, 8b, $n{=}5$)}
\begin{hwp}
$2.75\cdot 2^{5}=2.75\cdot 32 = 88$. Hex: $88=$\textbf{0x58}.
\end{hwp}

\subsection{HW5 P2 --- Decode Q3.4 ($N{=}8$, $n{=}4$)}
\begin{hwp}
\textbf{(a)} 0x77: $U{=}119$, MSB$=$0 $\Rightarrow I{=}119$;
$f=119/16=$\textbf{7.4375}.\\
\textbf{(b)} 0xAA: $U{=}170$, MSB$=$1 $\Rightarrow I=170-256=-86$;
$f=-86/16=$\textbf{$-5.375$}.
\end{hwp}

\subsection{HW5 P3 --- Decode Q2.5 (\texttt{0b1001\_1001})}
\begin{hwp}
$U{=}153$, MSB$=$1 $\Rightarrow I=153-256=-103$;
$f=-103/32=$\textbf{$-3.21875$}.
\end{hwp}

\subsection{HW5 P4 --- Encode in Q3.4 (range $[-8.0,+7.9375]$)}
\begin{hwp}
\textbf{(a)} $0.5\cdot 16=8 \Rightarrow$ \textbf{0x08}.\\
\textbf{(b)} $6.73\cdot 16 = 107.68 \to 108 \Rightarrow$ \textbf{0x6C}
(decode: $108/16=6.75$, err $0.02$).\\
\textbf{(c)} $-2.25\cdot 16=-36$, $+256=220 \Rightarrow$ \textbf{0xDC}.\\
\textbf{(d)} $-4.5\cdot 16=-72$, $+256=184 \Rightarrow$ \textbf{0xB8}.
\end{hwp}

\subsection{HW5 P5 --- Q15 MAC: $f_1f_2+f_3f_4$}
\begin{hwp}
$f_1{=}0.5,\,f_2{=}0.25,\,f_3{=}0.125,\,f_4{=}0.0625$.
\textbf{Encode} ($I{=}f\cdot 2^{15}$): \;
$I_1{=}$0x4000, $I_2{=}$0x2000, $I_3{=}$0x1000, $I_4{=}$0x0800.\\
\textbf{Mult+renorm} ($I_P=(I_AI_B)\!\gg\!15$):
$I_{P1}=2^{14}\!\cdot\!2^{13}\!\gg\!15=2^{12}=$0x1000;
$I_{P2}=2^{12}\!\cdot\!2^{11}\!\gg\!15=2^{8}=$0x0100.\\
\textbf{Add:} $I_A=$0x1100$=4352$.
\textbf{Decode:} $4352/32768=$\textbf{0.1328125} $\checkmark$
($0.5\!\cdot\!0.25+0.125\!\cdot\!0.0625$).
\end{hwp}

\subsection{Worked encode/decode shortcuts}
\begin{ex}
Powers of two: $f=2^{-k}\Rightarrow I=2^{n-k}$ (single-bit code).
Q15 of $0.5=2^{14}=$0x4000; Q15 of $0.0625=2^{11}=$0x0800.
\end{ex}

\begin{ex}
Q15.16 of $3.25$: $I=3.25\cdot 2^{16}=3\cdot 0x10000+0x4000=$\textbf{0x00034000}.
Returned in R0 (32-bit EABI return).
\end{ex}

\subsection{Common pitfalls}
\begin{itemize}
\item Forgetting the $+2^{N}$ wrap on negative encodes (yields wrong sign).
\item Using LSR (logical) instead of ASR for signed renormalize after multiply.
\item Adding numbers of \emph{different} Q-formats without aligning $n$
(shift smaller-$n$ left, or shift larger-$n$ right, before add).
\item Q15$\times$Q15 product is Q30 not Q15 --- must \asm{ASR \#15}.
\item Q15 cannot store $+1.0$ exactly (max $\approx 0.99997$).
\end{itemize}

\subsection{Cookbook table}
\begin{tabular}{@{}lll@{}}
\toprule
op & rule & note \\
\midrule
encode & $I=\mathrm{round}(f\cdot 2^{n})$ & $+2^{N}$ if $I<0$ \\
decode & $f=I\cdot 2^{-n}$ & $I=U-2^{N}$ if MSB \\
add Qm.n & $I_S=I_A+I_B$ & same format only \\
mul Qm.n & $I_P=(I_AI_B)\!\gg\!n$ & ASR for signed \\
Q15 mul & \asm{smull;lsr \#15} & 16$\to$32$\to$16 \\
Q15.16  & 32b, $n{=}16$ & R0 return \\
\bottomrule
\end{tabular}

% ----- part1/03_hw06_arith_eabi.tex -----
\section{HW6: Arithmetic + EABI Calling}

\subsection{Arithmetic instructions (Cortex-M4)}
\begin{tabular}{@{}lll@{}}
\toprule
Instr & Form & Flags \\
\midrule
\asm{ADD/SUB/RSB} & \asm{Rd,Rn,op2} & none \\
\asm{ADDS/SUBS/RSBS} & \asm{Rd,Rn,op2} & N,Z,C,V \\
\asm{ADC/SBC/RSC} & \asm{Rd,Rn,op2} & uses C in (S sets flags) \\
\asm{MUL/MLA/MLS} & 32$\times$32$\to$lo32 & none (\asm{MULS}: only N,Z) \\
\asm{UMULL/SMULL} & 32$\times$32$\to$64 & none \\
\asm{UMLAL/SMLAL} & RdLo:RdHi += A$\cdot$B & none \\
\asm{UDIV/SDIV} & \asm{Rd,Rn,Rm} & \textbf{never} set flags \\
\bottomrule
\end{tabular}\\
\textbf{RSB:} \asm{rsb Rd,Rn,op2} $=$ op2$-$Rn (reverse). \textbf{MLA} \asm{Rd,Rn,Rm,Ra}: Rd=Ra+Rn$\cdot$Rm. \textbf{MLS} \asm{Rd,Rn,Rm,Ra}: Rd=Ra$-$Rn$\cdot$Rm.\\
\textbf{Flag suffix rule:} append \asm{S} to update NZCV (e.g.\ \asm{ADD}$\to$\asm{ADDS}). Plain \asm{ADD} leaves flags untouched. \asm{UDIV/SDIV} have no S form. Div-by-zero $\to$ Rd=0 (no fault when CFG.DIV\_0\_TRP=0).

\subsection{ARM C / V flag conventions (subtract!)}
\textbf{ADD: } C=1 iff unsigned overflow out of bit 31. V=1 iff signed overflow.\\
\textbf{SUB/CMP: } \emph{C = NOT borrow}. C=1 means \emph{no} borrow (Rn$\ge$op for unsigned); C=0 means borrow occurred. V=1 iff signed overflow.\\
\asm{SBC Rd,Rn,Rm}: Rd $=$ Rn $-$ Rm $-$ \asm{!C} $=$ Rn$-$Rm$-$(1$-$C). So a low-word borrow (C=0) propagates an extra $-1$ into the high word.

\subsection{64-bit ADD / SUB pattern}
\begin{lstlisting}
@ A=R1:R0 (hi:lo), B=R3:R2; result back in R1:R0
mp_add_64bit_s:
    ADDS r0, r2     @ low,  C = unsigned ovf
    ADC  r1, r3     @ high + C
    BX   lr
mp_sub_64bit_s:
    SUBS r0, r2     @ low,  C = !borrow
    SBC  r1, r3     @ high - !C
    BX   lr
\end{lstlisting}
Convention: low word at lower-numbered reg (R0,R2), high at next (R1,R3). The S on the low op is mandatory; the high op must be ADC/SBC, never ADD/SUB.

\subsection{Modulo via UDIV+MUL+SUB (no MOD instr)}
Identity \asm{A\%B = A - (A/B)*B}. Leaf, R0--R2 caller-saved $\Rightarrow$ no PUSH/POP.
\begin{lstlisting}
umod32bit_s:                @ uint32 mod: R0=A,R1=B
    UDIV r2, r0, r1         @ q = A/B
    MUL  r2, r2, r1         @ q*B
    SUB  r0, r0, r2         @ A - q*B
    BX   lr
@ One-instruction form (M4 has MLS):
    UDIV r2, r0, r1
    MLS  r0, r2, r1, r0     @ r0 = r0 - r2*r1
    BX   lr
\end{lstlisting}
For signed modulo use \asm{SDIV} (sign of result follows dividend).

\subsection{EABI / AAPCS register usage}
\begin{tabular}{@{}llp{4.4cm}@{}}
\toprule
Reg & Role & Notes \\
\midrule
R0--R3 & arg / ret, scratch & \emph{caller-saved}; R0(:R1) hold return \\
R4--R11 & local & \emph{callee-saved}; PUSH on entry, POP on exit \\
R12 (IP) & scratch & may be clobbered by veneers \\
R13 (SP) & stack ptr & 8-byte aligned at public boundaries \\
R14 (LR) & link reg & save if you call (BL) anything \\
R15 (PC) & program ctr & POP \{...,pc\} returns \\
\bottomrule
\end{tabular}\\
\textbf{Args:} first 4 word-sized go in R0--R3, rest spill to stack (lowest addr = 5th arg). 64-bit args use an aligned even/odd pair (R0:R1 or R2:R3); cannot split across R3/stack--R3 is left unused if needed for alignment.\\
\textbf{Return:} 32-bit in \textbf{R0}. 64-bit in \textbf{R0:R1} (R0=lo, R1=hi). \texttt{void} returns nothing.\\
\textbf{Leaf rule:} if you don't BL anything and only touch R0--R3,R12, no PUSH/POP needed.

\subsection{Sub-word arg promotion at the CALL SITE}
Caller widens \asm{int8\_t/int16\_t/uint8\_t/uint16\_t} to 32-bit before placing in R0--R3:
\begin{itemize}
\item \textbf{Unsigned} src $\Rightarrow$ \emph{zero-extend} (UXTB/UXTH): upper bits = 0.
\item \textbf{Signed} src $\Rightarrow$ \emph{sign-extend} (SXTB/SXTH): upper bits = MSB of src.
\end{itemize}
Callee assumes the value is already 32-bit clean--reading the low byte without re-extending is wrong if caller failed to extend.

\subsection{Stack \& PUSH/POP rules (full descending)}
SP points at last-pushed word; PUSH pre-decrements, POP post-increments. \textbf{Lowest reg\# stored at lowest address, regardless of brace order.}\\
\asm{PUSH \{r7,r5,r8,r6\}} with SP=0x2000\_1020:\\
\begin{tabular}{@{}ll@{}}
0x2000\_101C & R8 (highest \#) \\
0x2000\_1018 & R7 \\
0x2000\_1014 & R6 \\
0x2000\_1010 & R5 (lowest \#) $\leftarrow$ new SP \\
\end{tabular}\\
SP after $=$ 0x2000\_1020$-$4$\cdot$4 $=$ 0x2000\_1010. POP reverses: lowest addr$\to$lowest reg.\\
\asm{POP \{...,pc\}} replaces the usual final \asm{BX lr}.

\subsection{HW6 P1: PUSH ordering walkthrough}
\begin{hwp}\textbf{P1.1} 4 regs $\cdot$ 4B $=$ 16 $\Rightarrow$ SP$=$0x20001020$-$0x10$=$\textbf{0x20001010}.\\
\textbf{P1.2} R5$=$(1$\ll$5)$+$(1$\ll$10)$=$32$+$1024$=$\textbf{0x0000\_0420}.\\
\textbf{P1.3} Sorted: R5@1010, R6@1014, R7@1018, R8@101C. R6$=$(1$\ll$6)$+$(1$\ll$12)$=$64$+$4096$=$\textbf{0x0000\_1040}.\\
\textbf{P1.4} Little-endian byte@0x20001014 $=$ LSB of 0x1040 $=$ \textbf{0x40}.
\end{hwp}

\subsection{HW6 P2--P3: arg promotion}
\begin{hwp}\textbf{P2 (unsigned, zero-ext):} \cee{fn1((uint8\_t)0xF,(uint8\_t)0xFF,0xFFF,0xFFFF)} into \cee{uint32\_t}\\
R0$=$0x0000\_000F, R1$=$0x0000\_00FF, R2$=$0x0000\_0FFF, R3$=$0x0000\_FFFF.\end{hwp}
\begin{hwp}\textbf{P3 (signed, sign-ext):} \cee{fn2((int8\_t)0xF,(int8\_t)0xFF)} into \cee{int32\_t}\\
0xF: MSB$=$0 $\Rightarrow$ R0$=$\textbf{0x0000\_000F}.\\
0xFF: MSB$=$1 (it's $-1$) $\Rightarrow$ R1$=$\textbf{0xFFFF\_FFFF}. Value $-1$ preserved.\end{hwp}

\subsection{HW6 P4: return-reg + Q15.16 encode}
\begin{hwp}Encode 3.25 in Q15.16: I$=$f$\cdot$2\textsuperscript{16}. $3\!\ll\!16=$0x0003\_0000; $0.25\!\cdot\!65536=$16384$=$0x4000.\\
$I=$\textbf{0x0003\_4000}; 32-bit return $\Rightarrow$ goes in \textbf{R0}.\end{hwp}

\subsection{HW6 P5: 64-bit ADD flag analysis}
\begin{hwp}A$=$R1:R0, B$=$R3:R2.\\
\textbf{T1:} A1$=$0x388$\!\ll\!$24, B1$=$0x289$\!\ll\!$24.\\
lo: 0x88000000$+$0x89000000$=$0x1\_11000000 $\Rightarrow$ R0$=$0x11000000, \textbf{C$=$1}.\\
hi: 3$+$2$+$1$=$6 $\Rightarrow$ R1$=$0x6. Result $=$ \cee{(int64\_t)0x611<<24}.\\
\textbf{T2:} A2$=$0x378$\!\ll\!$24, B2$=$0x279$\!\ll\!$24.\\
lo: 0x78000000$+$0x79000000$=$0xF1000000, fits $\Rightarrow$ \textbf{C$=$0}.\\
hi: 3$+$2$+$0$=$5. Result $=$ \cee{(int64\_t)0x5F1<<24}.\end{hwp}

\subsection{HW6 P6: 64-bit SUB flag analysis}
\begin{hwp}\textbf{T1:} A1$=$0x387$\!\ll\!$28, B1$=$0x28$\!\ll\!$28.\\
lo: 0x70000000$-$0x80000000 $\Rightarrow$ borrow! lo$=$0xF000\_0000, \textbf{C$=$0}.\\
hi: 0x38$-$0x02$-$1$=$0x35. Result $=$ \cee{(int64\_t)0x35F<<28}.\\
\textbf{T2:} A2$=$0x388$\!\ll\!$28, B2$=$0x288$\!\ll\!$28.\\
lo: 0x80000000$-$0x80000000$=$0, no borrow $\Rightarrow$ \textbf{C$=$1}.\\
hi: 0x38$-$0x28$-$0$=$0x10. Result $=$ \cee{(int64\_t)0x100<<28}.\\
\emph{Mantra:} subtract sets C$=$\textbf{!borrow}; SBC subtracts the inverse.\end{hwp}

\subsection{HW6 P7: modulo example}
\begin{hwp}A$=$17, B$=$5: \asm{UDIV} $\to$ q$=$3; \asm{MUL} $\to$ 15; \asm{SUB} $\to$ 17$-$15$=$\textbf{2}. Leaf, no stack frame, R2 is scratch.\end{hwp}

\subsection{Quick gotchas}
\begin{itemize}
\item \asm{MUL} (no S) never sets flags; \asm{MULS} sets only N,Z (C,V undefined). \asm{UMULL/SMULL/UMLAL/SMLAL} never set flags.
\item Reading flags after \asm{UDIV/SDIV} reads stale flags from the previous flag-setting instr.
\item \asm{MLS Rd,Rn,Rm,Ra}: Rd$=$Ra$-$Rn$\cdot$Rm--note the operand order.
\item Signed $\to$ unsigned promotion of a negative byte yields a huge positive uint32 (0xFF\_int8 $\to$ 0xFFFF\_FFFF as int32 $\to$ 4294967295 if reinterpreted unsigned).
\item PUSH brace order is cosmetic; assembler sorts. Don't rely on listed order for memory layout.
\item Full-descending: SP points at \emph{valid} top word; word at [SP] is the most recently pushed.
\end{itemize}

% ----- part1/04_hw07_shifts_logic.tex -----
\section{HW7: Shifts, Op2, Bitwise Logic}

\subsection{Shift instructions (4 forms + RRX)}
\begin{tabular}{@{}l l l@{}}
\textbf{LSL} & logical shift L & signed/unsigned, 0-fill at LSB \\
\textbf{LSR} & logical shift R & unsigned only, 0-fill at MSB \\
\textbf{ASR} & arith shift R & signed, sign-bit replicated at MSB \\
\textbf{ROR} & rotate right & bits wrap LSB$\to$MSB \\
\textbf{RRX} & rot R thru C & 33-bit rotate by \textbf{1} (C $\to$ MSB, LSB $\to$ C) \\
\end{tabular}

\textbf{S-suffix} (\asm{lsls/lsrs/asrs/rors}): updates flags. \textbf{C = last bit shifted out}; N/Z from result. For \asm{lsls Rd,Rs,\#n} $\Rightarrow$ C = bit$(32{-}n)$ of Rs. For \asm{lsrs Rd,Rs,\#n} $\Rightarrow$ C = bit$(n{-}1)$ of Rs.

\subsection{Operand-2 (flexible second operand)}
Three legal forms in any data-processing instr:
\begin{itemize}
\item \asm{\#imm} -- 8-bit rotated immediate (must encode)
\item \asm{Rm} -- bare register
\item \asm{Rm, \{LSL|LSR|ASR|ROR\} \#n} -- register w/ inline shift (\textbf{free}, no extra cycle)
\end{itemize}
\textbf{Shift count must be immediate} when used as Op2. Standalone \asm{lsl/lsr/asr/ror Rd,Rm,Rs} can take a register count -- only because the shift IS the instruction.

\textbf{GNU quirk:} \asm{add r0, r1, lsl \#3} is illegal even though it looks like \asm{r0 += r0<<3}; you must specify dest+src+Op2 explicitly: \asm{add r0, r0, r1, lsl \#3}. Never omit the Rn slot when Op2 carries a shift.

\subsection{One-instr scaling idioms (decompose $k = 2^n \pm 1$)}
\textbf{Multiply by integer $k$:}
\begin{lstlisting}
lsl r0, r1, #4         @ A = 16*B  (2^n)
add r0, r1, r1, lsl #1 @ A =  3*B
add r0, r1, r1, lsl #2 @ A =  5*B
add r0, r1, r1, lsl #3 @ A =  9*B
add r0, r1, r1, lsl #4 @ A = 17*B  (1+2^n)
rsb r0, r1, r1, lsl #3 @ A =  7*B  (2^n - 1)
rsb r0, r1, r1, lsl #4 @ A = 15*B
\end{lstlisting}
\textbf{Divide / fractional ($k/2^n$, signed):} use ASR.
\begin{lstlisting}
asr r0, r1, #4         @ A = B/16
add r0, r1, r1, asr #4 @ A = 17B/16  = B + B/16
sub r0, r1, r1, asr #4 @ A = 15B/16  = B - B/16
sub r0, r0, r0, asr #2 @ A = 0.75*A  (3/4)
\end{lstlisting}
Recall \asm{rsb Rd,Rn,Op2} $\Rightarrow$ Rd = Op2$-$Rn (reverse). LSR on a signed negative gives garbage; \textbf{always ASR for signed division}.

\subsection{Q31 / Q15.16 renormalization (post-SMULL)}
After \asm{smull Rlo,Rhi,Rn,Rm}, the 64-bit signed product spans $b_{63}\ldots b_0$. Extract Q$m$.$n$ output (bits $b_{m+n+15}\ldots b_n$) with a fused 3-instr renorm:
\begin{lstlisting}
@ Q15.16 result := bits b47..b16 of {Rhi:Rlo}
smull r1, r0, r0, r1   @ Rlo=R1, Rhi=R0
lsr   r0, r1, #16      @ low half: top 16 of LW
orr   r0, r0, r0, lsl #16  @ ... typical pair:
@ canonical:
lsr   r0, r1, #16          @ R0 = (LW >> 16)
orr   r0, r0, r3, lsl #16  @ OR (HW << 16)
\end{lstlisting}
\textbf{Rule:} for Q$m$.$n$, LW $\gg n$, HW $\ll (32{-}n)$, OR. \quad Q31: $\ll 1 / \gg 31$. \quad Q15.16: $\ll 16/\gg 16$. \quad Q16.15: $\ll 17/\gg 15$.

\subsection{Bitwise logic family}
\begin{tabular}{@{}l l l@{}}
\asm{AND Rd,Rn,Op2} & Rd = Rn \& Op2 & mask in \\
\asm{ORR Rd,Rn,Op2} & Rd = Rn $|$ Op2 & set bits \\
\asm{EOR Rd,Rn,Op2} & Rd = Rn $\oplus$ Op2 & toggle bits \\
\asm{BIC Rd,Rn,Op2} & Rd = Rn \& $\sim$Op2 & \textbf{clear} bits set in Op2 \\
\asm{ORN Rd,Rn,Op2} & Rd = Rn $|$ $\sim$Op2 & (no C operator) \\
\asm{MVN Rd,Op2} & Rd = $\sim$Op2 & bitwise NOT \\
\asm{BFC Rd,\#lsb,\#w} & clear $w$ bits at lsb & field clear \\
\asm{BFI Rd,Rn,\#lsb,\#w} & insert low $w$ of Rn at lsb & field insert \\
\end{tabular}

\textbf{C-equivalent gap:} BIC has no direct C operator -- write \cee{a \&= \~{}mask;} (or \cee{a \& \~{}m}). Same for ORN: \cee{a | \~{}b}.

\subsection{APSR access}
\asm{MRS Rd, APSR} -- read flags into Rd. \asm{MSR APSR\_nzcvq, Rs} -- write flags. Used for save/restore around critical sections; rare on exam.

\subsection{Bit-field idioms (memorize)}
\textbf{Replace field of width $w$ at position $s$ with value $v$:}
\begin{lstlisting}
@ a &= ~(MASK << s);   a |= (v << s);
ldr  r1, =((1<<w)-1)   @ width mask (e.g. 0xF for nibble)
bic  r0, r0, r1, lsl #s  @ clear field
ldr  r2, =v
orr  r0, r0, r2, lsl #s  @ insert value
\end{lstlisting}
\textbf{Insert digit at LSB end} (vacated nibble already 0):
\begin{lstlisting}
lsl r0, r0, #4    @ H3H2H1H0 -> H2H1H0 0
orr r0, r0, #3    @ -> H2H1H0 3
\end{lstlisting}
\textbf{Insert digit at MSB end:}
\begin{lstlisting}
lsr r0, r0, #4         @ -> 0 H3H2H1
ldr r1, =0x7000        @ 7000 not 8-bit-rot encodable -> use ldr=
orr r0, r0, r1         @ -> 7 H3H2H1
\end{lstlisting}
\textbf{Set/Clear/Toggle triad:} ORR sets, BIC clears, EOR toggles. MVN inverts.

\subsection{HW7 P1 -- Q15.16 from SMULL}
\begin{hwp}\textbf{P1.} Given \asm{smull r2,r3,r0,r1} produces \{R3:R2\} 64-bit, extract Q15.16 ($b_{47}\ldots b_{16}$):
\begin{lstlisting}
lsr r0, r2, #16        @ low 16 of result = top 16 of LW
orr r0, r0, r3, lsl #16  @ high 16 of result = low 16 of HW
\end{lstlisting}
\end{hwp}

\subsection{HW7 P2 -- LSLS/LSRS \& C flag}
\begin{hwp}\textbf{P2.} R0 = \asm{0xDD} = \asm{0b1101\_1101}.
\begin{lstlisting}
lsls r1, r0, #1   @ R1 = 0x1BA; bit shifted out = R0[31] = 0 -> C=0
lsrs r2, r0, #3   @ R2 = 0x1B;  last bit out = R0[2]  = 1 -> C=1
\end{lstlisting}
After both, C = \textbf{1} (from the last S-instr). Rule: LSLS\#$n$ $\Rightarrow$ C=src$[32{-}n]$; LSRS\#$n$ $\Rightarrow$ C=src$[n{-}1]$.
\end{hwp}

\subsection{HW7 P3 -- Single-instr scaling table}
\begin{hwp}\textbf{P3.} R1=B, write A=R0 in one instruction.
\begin{tabular}{@{}l l@{}}
$A = 16B$  & \asm{lsl r0, r1, \#4} \\
$A = 17B$  & \asm{add r0, r1, r1, lsl \#4} \\
$A = 15B$  & \asm{rsb r0, r1, r1, lsl \#4} \\
$A = B/16$ & \asm{asr r0, r1, \#4} \\
$A = 17B/16$ & \asm{add r0, r1, r1, asr \#4} \\
$A = 15B/16$ & \asm{sub r0, r1, r1, asr \#4} \\
\end{tabular}
\end{hwp}

\subsection{HW7 P4 -- BIC+ORR replace nibble}
\begin{hwp}\textbf{P4.} \cee{uint16\_t a = 0xH3H2H1H0} $\to$ \asm{0xH3H2 A H0} (set bits[7:4]=A).
\begin{lstlisting}
ldr r1, =15
bic r0, r0, r1, lsl #4   @ clear bits 7:4
ldr r2, =10
orr r0, r0, r2, lsl #4   @ OR in 0xA<<4
\end{lstlisting}
C: \cee{(a \& \textasciitilde(0xF<<4)) | (0xA<<4);} or \cee{(a \& 0xFF0F) | 0x00A0;}.
\end{hwp}

\subsection{HW7 P5 -- BIC/ORR/EOR hand-trace}
\begin{hwp}\textbf{P5.} R0=\asm{0xAD}=\asm{1010\_1101}, R1=\asm{0x03}=\asm{0000\_0011}.
\begin{tabular}{@{}l l l@{}}
\asm{lsr r2,r1,\#2} & R2 = 0 & \asm{0x00} \\
\asm{lsl r2,\#2}    & R2 = 0 & \asm{0x00} \\
\asm{bic r3,r0,r1}  & \asm{0xAD \& \~{}0x03 = 0xAD \& 0xFC} & \asm{0xAC} \\
\asm{orr r4,r3,r1}  & \asm{0xAC | 0x03} & \asm{0xAF} \\
\asm{eor r5,r0,r1}  & \asm{0xAD $\oplus$ 0x03} & \asm{0xAE} \\
\end{tabular}
\end{hwp}

\subsection{HW7 P6 -- Digit insert/replace patterns}
\begin{hwp}\textbf{P6(a).} \asm{H3H2H1H0} $\to$ \asm{H3H2 5 H0} (replace H1 with 5):
\begin{lstlisting}
bic r0, r0, #0xF0   @ #0xF0 IS 8-bit-rot encodable
orr r0, r0, #0x50
\end{lstlisting}
\textbf{P6(b).} $\to$ \asm{H2H1H0 3} (shift L, insert 3 at LSB):
\begin{lstlisting}
lsl r0, r0, #4
orr r0, r0, #3
\end{lstlisting}
\textbf{P6(c).} $\to$ \asm{7 H3H2H1} (shift R, insert 7 at MSB):
\begin{lstlisting}
lsr r0, r0, #4
ldr r1, =0x7000      @ 0x7000 not encodable as Op2 imm
orr r0, r0, r1
\end{lstlisting}
\textbf{Why ldr=} for \asm{0x7000}: doesn't fit 8-bit-rotated immediate format. Use \asm{ldr Rx,=const} when an immediate refuses to encode.
\end{hwp}

\subsection{Quick-reference summary}
\begin{itemize}
\item Op2 shift is \textbf{free} -- always prefer \asm{add Rd,Rn,Rm,lsl \#n} over a 2-instr \asm{lsl}+\asm{add}.
\item Always write the \textbf{dest} explicitly when Op2 has a shift -- no implicit \asm{Rd=Rn}.
\item LSL = signed \emph{and} unsigned mul-by-$2^n$. LSR = unsigned div only. ASR = signed div.
\item S-suffix on a shift puts the \textbf{last bit shifted out} into C -- handy single-bit pluck.
\item Field replace = BIC mask $\to$ ORR value-shifted-in. End-insert = shift word, ORR digit.
\item BIC has no C operator: \cee{a \& \~{}m}. ORN: \cee{a | \~{}b}. MVN: \cee{\~{}a}.
\item RRX rotates 33 bits (with C); no count needed -- always 1.
\item APSR via \asm{MRS}/\asm{MSR}; \asm{APSR\_nzcvq} field selector for write.
\end{itemize}

% ----- part1/05_hw08_nzcv_ccs.tex -----
% =============================================================================
% HW8 --- NZCV flags, condition codes, conditional branches (Lctr 18-19)
% =============================================================================
\section{HW8: NZCV Flags + CCS + Branch}

\subsection{APSR layout \& per-flag meaning}
The four arithmetic flags live in the top nibble of the \textbf{APSR} (Application Program Status Register, view of xPSR):
\begin{center}\scriptsize
\begin{tabular}{@{}cccccl@{}}
\toprule
b31 & b30 & b29 & b28 & b27 & rest \\
\midrule
\textbf{N} & \textbf{Z} & \textbf{C} & \textbf{V} & Q & \cee{0\dots 0} (reserved/IT/GE) \\
\bottomrule
\end{tabular}
\end{center}
\textbf{N} = sign of result = MSB of result word (1 $\Rightarrow$ negative when interpreted signed). \textbf{Z} = (result $==$ 0); pure equality bit, signed/unsigned agnostic. \textbf{C} = unsigned carry-out (add) \emph{or} \textbf{!Borrow} (sub) --- ARM convention is opposite of x86. \textbf{V} = signed overflow: result fell outside $[-2^{N-1},\,2^{N-1}-1]$. Q is the saturation sticky bit (DSP). N\&Z apply to logical ops too; C may be set by shifter; V is unaffected by logical ops.

\subsection{Hand computation in $N$-bit system}
Constants for an $N$-bit ALU: $C_N = 2^N$, $U_{max} = 2^N{-}1$, $S_{min} = -2^{N-1}$, $S_{max} = 2^{N-1}{-}1$. The processor doesn't know signed vs unsigned --- it sets \emph{all four flags every time} and the programmer picks which to consult via the CCS suffix.
\begin{enumerate}
  \item Convert each operand to its $N$-bit unsigned form (negatives: add $C_N$). Same bit pattern for both views.
  \item \textbf{Unsigned table:} compute $u_{tmp}$. If $u_{tmp} > U_{max}$ on add $\Rightarrow$ \textbf{C=1} (carry); subtract $C_N$. If $u_{tmp} < 0$ on sub $\Rightarrow$ \textbf{C=0} (borrow); add $C_N$. Otherwise C matches no-overflow side.
  \item \textbf{Signed table:} compute $s_{tmp}$ from signed operands. If $s_{tmp}$ falls outside $[S_{min},S_{max}]$ $\Rightarrow$ \textbf{V=1}, correct by $\pm C_N$. Else V=0.
  \item \textbf{N} = MSB of corrected binary (same pattern from both tables). \textbf{Z} = 1 iff result is all zeros.
\end{enumerate}
\textbf{ARM C-on-subtract:} after \cee{SUBS}, no borrow $\Rightarrow$ C=1, borrow $\Rightarrow$ C=0. \emph{Mnemonic:} ``C = NOT borrow''. (x86 inverts this --- exam trap if you've used x86.)

\subsection{P1 (42pt) --- 5-bit drills ($C_N{=}32$, $S_{min}{=}{-}16$, $S_{max}{=}15$)}
\begin{center}\scriptsize
\begin{tabular}{@{}llllcccc@{}}
\toprule
op & u-view & s-view & bin & N & Z & C & V \\
\midrule
$(-15){+}({-}5)$ & $17{+}27{=}44{>}31$ & $-20{<}{-}16$ & \cee{0\,1100} & 0 & 0 & \textbf{1} & \textbf{1} \\
$11{-}18$        & $-7{<}0$ borrow     & $11{-}({-}14){=}25$ & \cee{1\,1001} & \textbf{1} & 0 & \textbf{0} & \textbf{1} \\
\bottomrule
\end{tabular}
\end{center}
\begin{hwp}
\textbf{P1a:} both C and V fire because \emph{both} interpretations overflow. Binary \cee{0\,1100} reads as $12$ unsigned and $12$ signed --- they coincidentally agree only because of the wrap. \textbf{P1b:} the operand $18$ is unrepresentable in signed 5-bit, so \cee{10010} reinterprets to $-14$ and the signed subtraction is $11-(-14)=25$, which overflows positive. Sub borrow $\Rightarrow$ C=0 (ARM).
\end{hwp}

\subsection{Flag-only ops: CMP, CMN, TST, TEQ}
These four discard the result and \emph{only} update flags (no Rd):
\begin{center}\scriptsize
\begin{tabular}{@{}lllc@{}}
\toprule
mnem & equiv\ to & op & V affected? \\
\midrule
\cee{CMP Rn,Op2} & \cee{SUBS \_,Rn,Op2}      & subtract & yes \\
\cee{CMN Rn,Op2} & \cee{ADDS \_,Rn,Op2}      & add      & yes \\
\cee{TST Rn,Op2} & \cee{ANDS \_,Rn,Op2}      & bitwise AND & no  \\
\cee{TEQ Rn,Op2} & \cee{EORS \_,Rn,Op2}      & bitwise XOR & no  \\
\bottomrule
\end{tabular}
\end{center}
\textbf{Pseudo:} \asm{cmp r0,\#-5} is illegal as a literal but the assembler rewrites it to \asm{cmn r0,\#5} (subtracting $-5$ = adding $5$). Flags come out identical to the conceptual subtract, so \emph{signed CCS still works correctly}. Use \cee{TST}/\cee{TEQ} for bit-mask checks; they leave V untouched, so V-dependent CCS (\cee{GT}/\cee{LT}/\cee{GE}/\cee{LE}) shouldn't be paired with them.

\subsection{Condition-code suffixes (all 16)}
\begin{center}\scriptsize
\begin{tabular}{@{}lllll@{}}
\toprule
CCS & alias & meaning & flag eqn & domain \\
\midrule
\textbf{EQ} & --   & equal           & Z=1            & both \\
\textbf{NE} & --   & not equal       & Z=0            & both \\
\textbf{CS} & HS   & carry set / $\geq_u$ & C=1       & unsigned \\
\textbf{CC} & LO   & carry clr / $<_u$    & C=0       & unsigned \\
\textbf{MI} & --   & minus / negative& N=1            & rare \\
\textbf{PL} & --   & plus / non-neg  & N=0            & rare \\
\textbf{VS} & --   & ovfl set        & V=1            & rare \\
\textbf{VC} & --   & ovfl clr        & V=0            & rare \\
\textbf{HI} & --   & higher ($>_u$)  & C=1 \& Z=0     & unsigned \\
\textbf{LS} & --   & lower or same ($\leq_u$) & C=0 $\vert$ Z=1 & unsigned \\
\textbf{GE} & --   & $\geq_s$        & N=V            & signed \\
\textbf{LT} & --   & $<_s$           & N$\neq$V       & signed \\
\textbf{GT} & --   & $>_s$           & Z=0 \& N=V     & signed \\
\textbf{LE} & --   & $\leq_s$        & Z=1 $\vert$ N$\neq$V & signed \\
\textbf{AL} & --   & always          & ---            & default \\
\bottomrule
\end{tabular}
\end{center}
\textbf{Mnemonic ``HiLow unSun'':} the H/L family (\cee{HI/HS/LO/LS}) is \emph{unsigned}, the G/L family (\cee{GT/GE/LT/LE}) is \emph{signed}. EQ/NE work either way. Mismatching domain to operand type is the classic exam trap.

\subsection{Match shift type \& CCS to operand domain}
Pair the \emph{shift kind} and the \emph{CCS suffix} to the same domain or you'll get silently wrong answers:
\begin{center}\scriptsize
\begin{tabular}{@{}lll@{}}
\toprule
operand domain & shift for $\div 2^n$ & comparison CCS \\
\midrule
signed   (\cee{int})   & \cee{ASR \#n} (sign-extend) & GT GE LT LE \\
unsigned (\cee{uint})  & \cee{LSR \#n} (zero-fill)   & HI HS LO LS \\
\bottomrule
\end{tabular}
\end{center}
\cee{LSR} on a negative number gives garbage; \cee{ASR} on a large unsigned smears the sign bit. EQ/NE are domain-free.

\subsection{P2 (40pt) --- 6-bit CMP and CCS verify ($C_N{=}64$)}
\textbf{Case (a):} $r0{=}22$, $r1{=}20$. Subtract $\Rightarrow$ result $2$, no borrow, no signed overflow $\Rightarrow$ \textbf{N=0, Z=0, C=1, V=0}. \textbf{Case (b):} $r0{=}{-}20$, $r1{=}{-}22$. Unsigned forms $44$ and $42$; $44{-}42 = 2$ --- identical flags to (a). The ``equal flags from equal magnitudes'' fact is general: $a{-}b$ and $(a{+}C_N){-}(b{+}C_N)$ produce the same result and same NZCV.

\begin{center}\scriptsize
\begin{tabular}{@{}lcccccccc@{}}
\toprule
flags N=0,Z=0,C=1,V=0 & EQ & NE & GT & GE & LT & LE & HI & HS \\
\midrule
evaluate              & F  & T  & T  & T  & F  & F  & T  & T  \\
\bottomrule
\end{tabular}
\end{center}
\begin{hwp}
\textbf{P2 takeaway:} CCS evaluation is \emph{flag-driven, not value-driven}. Two CMPs on totally different operand pairs that happen to yield identical NZCV will branch identically. Therefore the same flag pattern $(0,0,1,0)$ supports \emph{both} ``$22 > 20$'' (signed) and ``$44 >_u 42$'' (unsigned) --- GT and HI both fire. EQ logic is just $Z=1$, NE is $Z=0$ --- domain-independent.
\end{hwp}

\subsection{P3 (20pt) --- max() via negative-logic branch}
C source: \cee{int max(int a,int b)\{ return a{>}{=}b ? a : b; \}}. Negative-logic skeleton: \emph{branch on the opposite CCS to skip the then-block}. For $\geq$ the opposite is $<$, so we use \cee{BLT}.
\begin{lstlisting}
mp_max_ab_i_s:
    cmp   r0, r1            @ a vs b (signed CMP)
    blt   .else             @ NL: a<b skips the then
.then:
    b     .end_if           @ a>=b: r0 already = a, fall to end
.else:
    mov   r0, r1            @ r0 = b
.end_if:
    bx    lr
\end{lstlisting}
\textbf{Why BLT (signed) and not BLO (unsigned)?} \cee{int} operands $\Rightarrow$ signed CCS family. Using \cee{BLO} would silently mis-handle negative $a$ (e.g., $a={-}1$, $b=1$: signed $-1<1$ true, but unsigned $0xFFFFFFFF > 1$ false). Domain mismatch.

\subsection{P4 (30pt) --- 64-bit sum-of-squares with SMLAL}
$\sum_{i=s}^{e} i^2$ overflows 32 bits quickly ($50000^2{=}2.5{\cdot}10^9 > 2^{31}$), so accumulate into a 64-bit pair using the multiply-accumulate variant. \cee{SMLAL Rlo,Rhi,Rn,Rm} computes \cee{\{Rhi:Rlo\} += Rn*Rm} (signed); \cee{UMLAL} is the unsigned twin. EABI returns 64-bit values in \cee{R0} (low) : \cee{R1} (high).
\begin{lstlisting}
mp_range_square_sum_standard_while_s:
    push  {r4, r5}
    ldr   r2, =0            @ sum_lo = 0
    ldr   r3, =0            @ sum_hi = 0
.lp:
    cmp   r0, r1            @ i vs e  (r0 reused as i)
    bgt   .end              @ NL: exit when i>e (signed)
    smlal r2, r3, r0, r0    @ {r3:r2} += i*i  (signed 32x32->64 MAC)
    add   r0, #1
    b     .lp
.end:
    mov   r0, r2            @ lo -> R0
    mov   r1, r3            @ hi -> R1
    pop   {r4, r5}
    bx    lr
\end{lstlisting}
\begin{hwp}
\textbf{Alt:} use \cee{SMULL r4,r5,r0,r0} then \cee{ADDS r2,r2,r4} \cee{ADC r3,r3,r5}; \cee{ADDS}/\cee{ADC} is the generic 64-bit-add idiom (low \cee{ADDS} sets C, high \cee{ADC} adds C). \cee{SMLAL} fuses both into one instruction. Reuse the input \cee{r0} as the loop counter to dodge a callee-saved push.
\end{hwp}

\subsection{Summary / common traps}
\begin{itemize}
  \item Processor sets \emph{all four flags every arithmetic op} --- the programmer chooses C vs V via the CCS suffix.
  \item Sub: \textbf{C = !borrow} (ARM); add: C = carry. Equal magnitudes (e.g. $22{-}20$ vs $({-}20){-}({-}22)$) produce equal flags.
  \item \cee{TST}/\cee{TEQ} \emph{don't touch V} --- never follow them with V-dependent CCS (\cee{GT}/\cee{LT}/\cee{GE}/\cee{LE}).
  \item \cee{cmp r0,\#-5} pseudo-assembles as \cee{cmn r0,\#5}; signed CCS still semantically valid.
  \item Match operand domain to both shift kind (\cee{ASR} signed / \cee{LSR} unsigned) and CCS family (G/L signed / H/L unsigned).
  \item Negative-logic branch reads cleaner: \cee{cmp;\,b\textit{!cond}\,.else} skips the then-block.
  \item For 64-bit accumulate use \cee{SMLAL/UMLAL} (or \cee{ADDS}/\cee{ADC} pair); return low in R0, high in R1.
\end{itemize}

% ----- part1/06_hw09_ldr_str.tex -----
\section{HW9: LDR/STR Addressing Modes}

\subsection{Three immediate-offset forms}
Let \asm{Ra}$=$base, \asm{\#os}$=$signed small immediate.
\begin{tabular}{@{}lll@{}}
\toprule
syntax & EA & R\textsubscript{a} after \\
\midrule
\asm{[Ra,\#os]}    & \asm{Ra+os} & unchanged (basic) \\
\asm{[Ra,\#os]!}   & \asm{Ra+os} & $\leftarrow$\asm{Ra+os} (pre) \\
\asm{[Ra],\#os}    & \asm{Ra}    & $\leftarrow$\asm{Ra+os} (post) \\
\bottomrule
\end{tabular}
\textbf{Mnem:} \asm{!}$=$update first then load ($\equiv$\cee{*++p});
trailing \asm{,\#os}$=$load first then update ($\equiv$\cee{*p++});
neither $=$\cee{*(p+k)}.

\subsection{Load-size suffixes}
\begin{tabular}{@{}llll@{}}
\toprule
op & bits & extend & C type \\
\midrule
\asm{LDR}   & 32 & ---           & \cee{int32\_t} \\
\asm{LDRH}  & 16 & zero          & \cee{uint16\_t} \\
\asm{LDRSH} & 16 & sign          & \cee{int16\_t}  \\
\asm{LDRB}  & 8  & zero          & \cee{uint8\_t}  \\
\asm{LDRSB} & 8  & sign          & \cee{int8\_t}   \\
\bottomrule
\end{tabular}
\textbf{STR} has \emph{no} sign variants: \asm{STR/STRH/STRB} only --- write
just truncates the low bits of the source register.

\subsection{Sign-extension trap}
\begin{tabular}{@{}ll@{}}
\toprule
load mem & R\textsubscript{d} (32b) \\
\midrule
LDRB 0x80      & 0x00000080 \\
LDRSB 0x80     & 0xFFFFFF80 \\
LDRSB 0x40     & 0x00000040 \\
LDRH 0xB0A0    & 0x0000B0A0 \\
LDRSH 0xB0A0   & 0xFFFFB0A0 \\
\bottomrule
\end{tabular}
Sign bit$=$bit-7 (B) or bit-15 (H). If set under SB/SH, top 24/16
bits become \asm{1}.

\subsection{Register-offset (scaled index)}
\asm{[Ra, Ri, LSL \#n]} $\Rightarrow$ EA $=$ \asm{Ra + (Ri$\!\ll\!$n)}.
$n=\log_2(\text{elem size})$:
\asm{\#0}/B,SB; \asm{\#1}/H,SH; \asm{\#2}/word; \asm{\#3}/dword.
\textbf{No writeback form} --- \asm{Ra} never updates.
This is the \cee{p[i]} idiom.

\subsection{LDRD / STRD --- 64-bit pair}
\asm{LDRD Rt,Rt2,[Ra,...]} loads two consecutive words into
\textbf{consecutive regs} (Rt, Rt2$=$Rt+1).
Forms: basic \asm{[Ra,\#os]}, pre \asm{[Ra,\#os]!}, post \asm{[Ra],\#os}.
\textbf{NO reg-offset form.}
Rt$=$word@EA, Rt2$=$word@EA$+4$.

\subsection{C $\leftrightarrow$ ASM cookbook}
\begin{tabular}{@{}ll@{}}
\toprule
C idiom & ARM form \\
\midrule
\cee{*p}                  & \asm{LDR r,[Ra]}            \\
\cee{*(p+k)}              & \asm{LDR r,[Ra,\#k*sz]}     \\
\cee{p[i]}                & \asm{LDR r,[Ra,Ri,LSL \#n]} \\
\cee{*p++}                & \asm{LDR r,[Ra],\#sz}       \\
\cee{*++p}                & \asm{LDR r,[Ra,\#sz]!}      \\
\cee{*p = v} (no upd)     & \asm{STR r,[Ra]}            \\
\cee{*p++ = v}            & \asm{STR r,[Ra],\#sz}       \\
\cee{*++p = v}            & \asm{STR r,[Ra,\#sz]!}      \\
\bottomrule
\end{tabular}
\textbf{What's modified:}\;
prefix \asm{++} on \cee{*p} $\Rightarrow$ value in memory changes
$\Rightarrow$ need explicit STR writeback;
postfix on \emph{pointer} $\Rightarrow$ pointer changes (post-index);
both ($\cee{*p8b++ = ++*p8a++}$) $\Rightarrow$ LDRB post + ADD +
STRB \asm{[Ra,\#-1]} writeback.

\subsection{HW9 P1 --- Trace through 16 bytes}
Memory: 16 bytes \asm{0x00..0x0F} at \asm{0x20000000} (LE words:
\asm{ipt+0=0x03020100}, \asm{+4=0x07060504},
\asm{+8=0x0B0A0908}, \asm{+12=0x0F0E0D0C}).
\begin{lstlisting}
ldr r4, =0x20000000     @ (1) R4=0x20000000
ldr r0, [r4, #4]        @ (2) basic
ldr r1, [r4, #2]!       @ (3) pre-idx
ldr r2, [r4], #4        @ (4) post-idx
ldr r3, [r4]            @ (5) basic, no offset
\end{lstlisting}
\begin{tabular}{@{}llll@{}}
\toprule
step & EA & dest & R4 after \\
\midrule
(1) & ---        & ---           & 0x20000000 \\
(2) & 0x20000004 & R0=0x07060504 & 0x20000000 \\
(3) & 0x20000002 & R1=0x05040302 & 0x20000002 \\
(4) & 0x20000002 & R2=0x05040302 & 0x20000006 \\
(5) & 0x20000006 & R3=0x09080706 & 0x20000006 \\
\bottomrule
\end{tabular}
\textbf{Key:} pre-idx (3) loads from \emph{new} R4; post-idx (4) loads
from \emph{old} R4. Loads at non-word-aligned addresses (0x...02) work
fine on Cortex-M but cross word boundaries.

\subsection{HW9 P2 --- ASM $\to$ C prototypes}
\begin{lstlisting}
mp_task1:                  mp_task2:
  ldrsh r2,[r0],#4           ldrh  r2,[r0,#2]!
  str   r2,[r1,#0]           str   r2,[r1,#0]
  ldrsh r3,[r0],#4           ldrh  r3,[r0,#2]!
  str   r3,[r1,#4]           str   r3,[r1,#4]
  bx    lr                   bx    lr
\end{lstlisting}
\textbf{Task1:} \cee{void mp\_task1(int16\_t *pSrc, int32\_t *pDst)} ---
LDRSH$\Rightarrow$signed src; STR (word)$\Rightarrow$word dst.
Post-idx \asm{\#4} skips \emph{every other halfword} (stride$=2$ in
\cee{int16\_t} units).
\textbf{C:} \cee{*p32 = (int32\_t)*p16s; p16s += 2;} \newline \cee{
*(p32+1) = (int32\_t)*p16s; p16s += 2;}

\textbf{Task2:} \cee{void mp\_task2(uint16\_t *pSrc, int32\_t *pDst)} ---
LDRH$\Rightarrow$unsigned. Pre-idx \asm{\#2} advances first $\Rightarrow$
\emph{skips arr[0]}, reads \cee{arr[1]} then \cee{arr[2]}.
\textbf{C:} \cee{p16u += 1; *p32 = (int32\_t)*p16u;} \newline \cee{
p16u += 1; *(p32+1) = (int32\_t)*p16u;}

\subsection{HW9 P3 --- Four tasks, same data}
\cee{ipt[16]} initialized \cee{ipt[i]=i<<4} ($0x00,0x10,\ldots,0xF0$).
Words (LE, treating ipt as \cee{int32\_t*}):
\begin{tabular}{@{}ll@{}}
\toprule
addr & word \\
\midrule
\cee{ipt+0}  & 0x30201000 \\
\cee{ipt+4}  & 0x70605040 \\
\cee{ipt+8}  & 0xB0A09080 \\
\cee{ipt+12} & 0xF0E0D0C0 \\
\bottomrule
\end{tabular}
Entry: \asm{R0}$=$\asm{0x20001010} (\cee{ipt}),
\asm{R1}$=$\asm{0x20001040} (\cee{opt}).
T4 also gets \asm{R2}$=1$.

\subsubsection*{T1 --- basic offset (R0 unchanged)}
\begin{lstlisting}
ldr r2,[r0,#4]      @ R2=0x70605040
str r2,[r1,#0]      @ opt[0]
ldr r2,[r0,#0xC]    @ R2=0xF0E0D0C0
str r2,[r1,#4]      @ opt[1]
\end{lstlisting}
\textbf{Print:} \asm{Out1=0x70605040, Out2=0xF0E0D0C0}.\;
\textbf{After:} R0$=$0x20001010, R2$=$0xF0E0D0C0.\;
\textbf{C:} \cee{int32\_t *p=(int32\_t*)ipt; opt[0]=p[1]; opt[1]=p[3];}

\subsubsection*{T2 --- pre-index ($\times 2$)}
\begin{lstlisting}
ldr r2,[r0,#4]!     @ R0=0x...14, R2=0x70605040
str r2,[r1,#0]
ldr r2,[r0,#4]!     @ R0=0x...18, R2=0xB0A09080
str r2,[r1,#4]
\end{lstlisting}
\textbf{Print:} \asm{Out1=0x70605040, Out2=0xB0A09080}.\;
\textbf{After:} R0$=$0x20001018, R1$=$0x20001040 (unchanged).\;
\textbf{C:} \cee{p++; opt[0]=*p; p++; opt[1]=*p;}

\subsubsection*{T3 --- post-index ($\times 2$)}
\begin{lstlisting}
ldr r2,[r0],#4      @ R2=0x30201000, R0=0x...14
str r2,[r1,#0]
ldr r2,[r0],#4      @ R2=0x70605040, R0=0x...18
str r2,[r1,#4]
\end{lstlisting}
\textbf{Print:} \asm{Out1=0x30201000, Out2=0x70605040}.\;
\textbf{After:} R0$=$0x20001018.\;
\textbf{C:} \cee{opt[0]=*p; p++; opt[1]=*p; p++;}\\
\textbf{T2 vs T3:} pre starts at \cee{ipt+4}, post starts at \cee{ipt+0}
--- pre and post end at \emph{same} R0 but read \emph{different} words.

\subsubsection*{T4 --- register offset \asm{LSL \#2}}
Entry R2$=1$.
\begin{lstlisting}
ldr r3,[r0,r2,lsl #2]   @ ipt+4, R3=0x70605040
str r3,[r1,#0]
add r2,#1               @ R2=2
ldr r3,[r0,r2,lsl #2]   @ ipt+8, R3=0xB0A09080
str r3,[r1,#4]
add r2,#1               @ R2=3
\end{lstlisting}
\textbf{Print:} \asm{Out1=0x70605040, Out2=0xB0A09080}.\;
\textbf{After:} R0$=$0x20001010 (unchanged!), R2$=$3, R3$=$0xB0A09080.\;
\textbf{C:} \cee{int i=1; opt[0]=p[i]; i++; opt[1]=p[i]; i++;}

\subsection{Pre vs post --- ending pointer is the same}
For two consecutive same-stride loads from base \asm{Ra}:
\textbf{pre} reads \asm{[Ra+s, Ra+2s]} ending at \asm{Ra+2s};
\textbf{post} reads \asm{[Ra, Ra+s]} ending at \asm{Ra+2s}.
Same \emph{final} pointer, different \emph{data}.

\subsection{Common pitfalls}
\begin{itemize}\setlength{\itemsep}{0pt}
\item Pre \asm{[Ra,\#os]!} vs post \asm{[Ra],\#os}: \asm{!} \emph{inside}
      brackets $=$ EA changed inside the access.
\item Stride matches elem size: \cee{int16\_t*}+1$=$\asm{\#2};
      \cee{int32\_t*}+1$=$\asm{\#4}.
\item LDRSB/LDRSH on 0x80/0x8000 $\to$ 0xFFFFFF80/0xFFFF8000.
\item Reg-offset has no writeback; to advance use
      \asm{ADD Ra,Ra,Ri,LSL \#n}.
\item LDRD/STRD: consecutive regs, no reg-offset form.
\item ``Follow assembly's intention'': if .s treats \cee{uint8\_t*} as
      \cee{int32\_t*}, stride is 4 bytes regardless of declared type.
\item \cee{*p8b++ = ++*p8a++} $\Rightarrow$ \asm{LDRB+ADD+STRB[r0,\#-1]}
      writeback to original address.
\end{itemize}

\subsection{Decision table}
\begin{tabular}{@{}lll@{}}
\toprule
need & form & note \\
\midrule
read element, keep ptr  & \asm{[Ra,\#k]}      & basic    \\
walk forward (post)     & \asm{[Ra],\#sz}     & \cee{*p++}    \\
skip first, then walk   & \asm{[Ra,\#sz]!}    & \cee{*++p}    \\
random index            & \asm{[Ra,Ri,LSL \#n]}& \cee{p[i]}; no upd \\
two words at once       & \asm{LDRD r,r+1,...}& consec. reg pair \\
signed byte/half        & \asm{LDRSB/LDRSH}   & sign-ext to 32 \\
unsigned byte/half      & \asm{LDRB/LDRH}     & zero-ext to 32 \\
store (any size)        & \asm{STR/STRH/STRB} & no sign variants \\
\bottomrule
\end{tabular}

% ----- part1/07_hw10_ptr_fn_call.tex -----
\section{HW10: C Pointers + LDR/STR + Function Calls}

\subsection{C precedence (rows that matter)}
\begin{tabular}{@{}clll@{}}
\toprule
L & ops & assoc & note \\
\midrule
1 & \cee{() [] . -> ++ -- (post)} & L$\to$R & postfix inc/dec \\
2 & \cee{++ -- (pre)}, \cee{* \& (cast)} & R$\to$L & deref, addr-of \\
14 & \cee{= += ...} & R$\to$L & assignment \\
\bottomrule
\end{tabular}
\textbf{Rule:} postfix \cee{p++} (L1) yields OLD pointer; deref/prefix-\cee{++}/cast share L2 R$\to$L. Decode right-to-left.

\subsection{Idiom $\to$ ASM cookbook}
\begin{tabular}{@{}lll@{}}
\toprule
C & ASM & note \\
\midrule
\cee{v=*(p+k)}  & \asm{ldrb r0,[r1,\#k]}   & basic, p unchanged \\
\cee{v=*++p}    & \asm{ldrb r0,[r1,\#1]!}  & pre (p first) \\
\cee{v=*p++}    & \asm{ldrb r0,[r1],\#1}   & post (load first) \\
\cee{*(p+k)=v}  & \asm{strb r0,[r1,\#k]}   & basic store \\
\cee{*p++=v}    & \asm{strb r0,[r1],\#1}   & post store \\
\cee{++(*p)}    & ldr;add\#1;str           & READ-MOD-WRITE \\
\cee{(*p)++}    & ldr;mov;add\#1;str       & return old \\
\cee{++*p++}    & ldr post;add;str\asm{[\#-1]} & WB to OLD \\
\cee{*p++=*q++} & ldr post; str post       & two ptrs walk \\
\bottomrule
\end{tabular}
\textbf{Hook:} \texttt{!}$=$urgent update BEFORE; trailing \texttt{\#}$=$update AFTER.

\subsection{Type suffix table (load only --- stores zero/sign agnostic)}
\begin{tabular}{@{}llll@{}}
\toprule
sfx & C type & ext & sample \\
\midrule
(none) & \cee{int32\_t} & --- & \asm{ldr} \\
B  & \cee{uint8\_t}  & zero $\to$ 32 & 0x80$\to$0x0000\_0080 \\
SB & \cee{int8\_t}   & sign $\to$ 32 & 0x80$\to$0xFFFF\_FF80 \\
H  & \cee{uint16\_t} & zero $\to$ 32 & 0xB0A0$\to$0x0000\_B0A0 \\
SH & \cee{int16\_t}  & sign $\to$ 32 & 0xB0A0$\to$0xFFFF\_B0A0 \\
\bottomrule
\end{tabular}
Sign-extension test: MSB of the loaded chunk (bit 7 for SB, bit 15 for SH).
\textbf{STR has no S/U variant} (you store the bits you have).

\subsection{Pointer arithmetic scaling}
Increment of typed pointer scales by \cee{sizeof(*p)}: \cee{p8++}$=$+1B,
\cee{p16++}$=$+2B, \cee{p32++}$=$+4B. Register-offset shift exponent matches
log$_2$(elem): \asm{[Rb,Ri,LSL \#0/1/2/3]} for byte/hw/word/dword.
\textbf{Typed difference} \cee{q-p} returns \emph{elements}, not bytes
(divide-by-size baked in).

\subsection{HW10 P1 --- four-line memory \& pointer trace}
Array: \cee{uint8\_t my\_arr[16]} at \texttt{0x2000\_0100}, \cee{my\_arr[i]=2i}.
\cee{p8a=0x...100}, \cee{p8b=0x...10C}.
\begin{tabular}{@{}lllll@{}}
\toprule
ln & RHS & memory effect & p8a$\to$ & p8b$\to$ \\
\midrule
4 & \cee{*p8b++=*(p8a+2)}     & rd 0x102$=$0x04; \texttt{0x10C}$\leftarrow$0x04 & 0x100 & 0x10D \\
5 & \cee{*p8b++=++*(p8a+4)}   & 0x104:0x08$\to$0x09; \texttt{0x10D}$\leftarrow$0x09 & 0x100 & 0x10E \\
6 & \cee{*p8b++=++*p8a++}     & 0x100:0x00$\to$0x01; \texttt{0x10E}$\leftarrow$0x01 & 0x101 & 0x10F \\
7 & \cee{*p8b++=*p8a++}       & rd 0x101$=$0x02; \texttt{0x10F}$\leftarrow$0x02 & 0x102 & 0x110 \\
\bottomrule
\end{tabular}

\textbf{Bytes at 0x10C--0x10F:} \texttt{0x04,0x09,0x01,0x02}.
\textbf{Mem changes by RHS:} L4 none, L5 0x104(0x08$\to$0x09), L6 0x100(0x00$\to$0x01), L7 none.

\textbf{Line 6 dissection (\cee{++*p8a++}, the trickiest):}
postfix \cee{p8a++} (L1) gives OLD pointer 0x100; deref reads 0x00; prefix
\cee{++} adds 1, writes 0x01 BACK to 0x100. So memory at 0x100 changes; the
loaded value is 0x01; p8a now 0x101.

\subsection{HW10 P1.3 --- ASM (leaf, \cee{my\_fn(uint8\_t*, uint8\_t*)})}
\begin{lstlisting}
my_fn:                       @ R0=p8a, R1=p8b (caller-saved)
  @ L4: *p8b++ = *(p8a+2);
    ldrb r2, [r0, #2]          @ basic offset, p8a stays
    strb r2, [r1], #1          @ post-idx: *p8b=r2; p8b++
  @ L5: *p8b++ = ++*(p8a+4);
    ldrb r2, [r0, #4]
    add  r2, r2, #1            @ prefix-inc the value
    strb r2, [r0, #4]          @ WRITEBACK -- don't forget!
    strb r2, [r1], #1
  @ L6: *p8b++ = ++*p8a++;
    ldrb r2, [r0], #1          @ post-idx: r2=*OLD p8a; R0+=1
    add  r2, r2, #1
    strb r2, [r0, #-1]         @ writeback to OLD addr (R0-1)
    strb r2, [r1], #1
  @ L7: *p8b++ = *p8a++;
    ldrb r2, [r0], #1
    strb r2, [r1], #1
    bx   lr                     @ leaf: no PUSH/POP, no LR save
\end{lstlisting}

\textbf{Three patterns to memorize:}
(1) \cee{++(*(p+k))} $=$ ldrb / add\#1 / strb-to-same-offset / strb-to-dest;
(2) \cee{++*p++}     $=$ ldrb post-idx / add / strb \asm{[\#-1]} / strb post-idx;
(3) \cee{*p++}       $=$ ldrb post-idx alone (no writeback).

\subsection{AAPCS: who-saves-what (the only table you need)}
\begin{tabular}{@{}lll@{}}
\toprule
reg & role & preserved? \\
\midrule
R0--R3 & args + return value & no (caller-saved, scratch) \\
R4--R11& callee scratch       & \textbf{yes} (PUSH if used) \\
R12 (IP)& scratch              & no \\
R13 (SP)& stack pointer        & yes (8-byte aligned at boundary) \\
R14 (LR)& return addr          & save if you BL (stem) \\
R15 (PC)& program counter      & --- \\
\bottomrule
\end{tabular}

\textbf{Caller widens narrow types} (\cee{int8\_t/int16\_t}) to 32b before
placing in R0--R3: signed args sign-ext, unsigned zero-ext. 64-bit args use
register pair (lo, hi), even-aligned (R0:R1 or R2:R3).

\subsection{Leaf vs Stem function templates}
\textbf{Leaf} (no BL inside, $\le$4 scratch regs needed): no push, no pop.
\begin{lstlisting}
leaf_fn:                  @ R0..R3 in, R0 out
    @ work in R0--R3, R12
    bx   lr
\end{lstlisting}

\textbf{Stem} (calls others or needs R4--R11): push LR + any callee-saved used.
Stash-args pattern: copy R0--R3 into R4--R7 \emph{before} BL clobbers them.
\begin{lstlisting}
stem_fn:
    push {r4-r7, lr}        @ save callee + return addr
    mov  r4, r0             @ stash args (R0..R3 die on next BL)
    mov  r5, r1
    mov  r6, r2
    mov  r7, r3
    @ ... bl other_fn; arg in R0; result in R0 ...
    @ recover args from R4--R7 between calls
    pop  {r4-r7, pc}        @ pop PC = return; replaces BX LR
\end{lstlisting}
Push/pop list \emph{must match}; pushing odd \# of regs $\to$ pad for 8B align
(usually push LR alongside even count). \asm{POP \{...,pc\}} returns directly.

\subsection{HW10 P2 --- four-task LDRD/SB/SH/H trace}
Setup: \cee{ipt[i]=i<<4} $\to$ bytes 0x00,0x10,0x20,...,0xF0. Word at ipt+0
$=$ 0x30201000 (little-endian: LSB at lowest addr). Halfwords at
ipt+0/+2/.../+14 $=$ 0x1000,0x3020,0x5040,0x7060,0x9080,0xB0A0,0xD0C0,0xF0E0.
Entry R0$=$0x2000\_1010 (ipt), R1$=$0x2000\_1040 (opt).

\textbf{T1: LDRD pre-idx (8B$\to$2 regs).}
\asm{ldrd r2,r3,[r0,\#8]!} : R0$+$=8 first, then load lo word at \texttt{[R0]}
$=$ 0xB0A09080, hi word at \texttt{[R0+4]} $=$ 0xF0E0D0C0.
\asm{strd r2,r3,[r1,\#0]} writes opt[0],opt[1].
\textbf{Print:} \texttt{Out1=0xB0A09080, Out2=0xF0E0D0C0}.

\textbf{T2: LDRSB pre-idx, sign-extend.}
\asm{ldrsb r2,[r0,\#4]!}: R0$\to$ipt+4, byte 0x40 (bit7$=$0) $\to$ R2$=$0x0000\_0040.
\asm{ldrsb r3,[r0,\#4]!}: R0$\to$ipt+8, byte 0x80 (bit7$=$1) $\to$ R3$=$0xFFFF\_FF80.
After: R0$=$\textbf{0x2000\_1018}, R1$=$0x2000\_1040 (STRD has no \texttt{!}).
\textbf{Print:} \texttt{Out1=0x40, Out2=0xFFFFFF80}.

\textbf{T3: LDRSH post-idx, odd stride 10.}
\asm{ldrsh r2,[r0],\#10}: load hw at ipt+0$=$0x1000 (bit15$=$0) $\to$
R2$=$0x0000\_1000; THEN R0$+$=10.
\asm{ldrsh r3,[r0],\#4}: hw at ipt+10$=$0xB0A0 (bit15$=$1) $\to$
R3$=$0xFFFF\_B0A0; R0$+$=4.
After: R0$=$\textbf{0x2000\_101E} (unaligned!).
\textbf{Print:} \texttt{Out1=0x1000, Out2=0xFFFFB0A0}.

\textbf{T4: LDRH register-offset, zero-ext.} Called \cee{task4(ipt,opt,1)}, R2$=$1.
\asm{ldrh r3,[r0,r2,LSL \#1]}: EA$=$ipt$+$2, hw 0x3020 $\to$ R3$=$0x0000\_3020;
\asm{str r3,[r1,\#0]}; \asm{add r2,\#1} (R2$=$2); reload at ipt$+$4 $=$
0x5040 $\to$ R3$=$0x0000\_5040; \asm{str r3,[r1,\#4]}; \asm{add r2,\#1} (R2$=$3).
\textbf{Print:} \texttt{Out1=0x3020, Out2=0x5040}. \textbf{R2$=$0x3, R3$=$0x5040.}

\subsection{Sign-extend pivots: LOAD vs CALL}
\textbf{LOAD:} LDRSB/LDRSH test bit 7/15 of loaded data, widen to 32b. LDRB/LDRH zero-pad.
\textbf{CALL:} caller widens narrow C arg per type signedness BEFORE \asm{BL}. Two independent extensions.

\subsection{HW10 P2 --- ASM$\to$C reverse}
\textbf{T1} 8B copy: \cee{*((int64\_t*)opt)=*((int64\_t*)ipt+1)} (or word-by-word with \cee{(int32\_t*)}).
\textbf{T2} pre-inc \cee{int8\_t*}: \cee{p8+=4; opt[0]=*p8; p8+=4; opt[1]=*p8;}.
\textbf{T3} post-inc \cee{int16\_t*}: \cee{opt[0]=*p16; opt[1]=*(p16+5);} (\#10$=$5 elem).
\textbf{T4} idx \cee{uint16\_t*}: \cee{int i=1; opt[0]=*(p16u+i); i++; opt[1]=*(p16u+i); i++;}.

\subsection{Array-loop $+$ call pattern (compute index $\to$ ldr/bl/str)}
Idiom: \cee{for(i;..) arr[i]=fn(arr[i],...)} compiles to:
\begin{lstlisting}
    @ R5=&arr (base, callee-saved), R8=i (callee-saved)
loop:
    ldr  r0, [r5, r8, lsl #2]   @ R0 = arr[i]   (32b elements)
    bl   fn                      @ R0 = result; clobbers R0--R3,R12
    str  r0, [r5, r8, lsl #2]   @ arr[i] = R0
    add  r8, r8, #1
    cmp  r8, r9                  @ R9 = N (callee-saved)
    blt  loop
\end{lstlisting}
\textbf{Why callee-saved (R4--R11) for base/index/N:} they survive across
\asm{BL fn}. Putting them in R0--R3/R12 would corrupt on every call.

\subsection{Fast-failure modes}
\begin{itemize}
\item Forgetting STRB writeback after \cee{++(*p)}: needs \asm{ldr;add;str} then dest store.
\item \asm{[r0,\#4]!} (pre, urgent BEFORE) vs \asm{[r0],\#4} (post, AFTER).
\item LDRB on \cee{int8\_t} $\to$ wrong zero-pad; use LDRSB.
\item STRD/LDRD w/o \texttt{!} doesn't update Ra (T1 R1 unchanged; T2 R0 advances).
\item Reg-offset shift: LSL \#0/1/2/3 for byte/hw/word/dword.
\item No \asm{push \{lr\}} in stem $\to$ \asm{BL} clobbers LR, return broken.
\end{itemize}

\subsection{One-line summary table}
\begin{tabular}{@{}lll@{}}
\toprule
goal & instruction & C \\
\midrule
\cee{*(p+k)}      & \asm{ldr\{B/H\}\{S\} Rd,[Ra,\#os]} & basic, p stays \\
\cee{*++p}        & \asm{...,[Ra,\#os]!}                & pre, p$+$=os first \\
\cee{*p++}        & \asm{...,[Ra],\#os}                 & post, load then p$+$=os \\
\cee{*(p+i)} typed& \asm{...,[Ra,Ri,LSL \#k]}           & k$=$log$_2$elem \\
8B copy           & \asm{ldrd}/\asm{strd}               & 64b, two regs \\
sign-ext load     & \asm{ldrsb}/\asm{ldrsh}             & test bit 7/15 \\
leaf return       & \asm{bx lr}                          & no push \\
stem return       & \asm{pop \{...,pc\}}                & pushed LR \\
\bottomrule
\end{tabular}

% ----- part1/08_hw11_if_flow.tex -----
\section{HW11: If-Based Flow Control (PL/NL/CEX)}

\subsection{Three styles for \cee{if-else} in ASM}
Given C \cee{if (cond) then; else else;}, three idiomatic ASM templates (Lctr 25):
\begin{tabular}{@{}l l l@{}}
\toprule
Style & Branch on & Layout \\
\midrule
\textbf{PL} (positive logic) & \textbf{same} CCS as C & else first, fall thru / B end / then \\
\textbf{NL} (negative logic) & \textbf{inverse} CCS  & then first, fall thru / B end / else \\
\textbf{CEX} (cond exec)     & no branch (IT block)  & all instrs in straight line \\
\bottomrule
\end{tabular}\\
NL reads closest to C (then-block falls through, matches \cee{if(c){...}}). PL is "branch when C-true". CEX is best when each arm is $\le$ 4 instructions total. CCS in \asm{IT..} = NL cond (the inverse of the C test).

\subsection{IT block syntax + hard rules}
\asm{IT\{x\{y\{z\}\}\} cond} -- up to \textbf{4} conditional instrs. First slot always \asm{T} (matches \asm{cond}); each later slot is \asm{T} (same cond) or \asm{E} (\textbf{inverse} cond).
\textbf{Common forms (T = same cond, E = inverse):} \asm{IT} (T), \asm{ITT} (T,T), \asm{ITE} (T,E), \asm{ITTT} (T,T,T), \asm{ITTE} (T,T,E), \asm{ITEE} (T,E,E), \asm{ITTEE} (T,T,E,E), \asm{ITTTT} (T,T,T,T).\\
\textbf{Hard rules (Lctr 24, p.\ 9):} max \textbf{4} cond instrs; \textbf{cond B} (\asm{Bcc}) only in \textbf{LAST} slot; \textbf{CBZ/CBNZ FORBIDDEN} inside IT (use \asm{cmp+Bcc}); each slot's CCS must match T/E role; cond LDR/STR (pre/post-idx) \textbf{legal} in IT.

\subsection{CBZ / CBNZ}
\asm{CBZ Rn,lbl} = "if Rn==0 goto lbl"; \asm{CBNZ Rn,lbl} = "if Rn!=0". \textbf{Forward only}, max \textbf{126 B} ahead. \asm{Rn} restricted to \textbf{R0--R7}. \textbf{Does NOT set flags}. Replaces \asm{cmp Rn,\#0 ; beq lbl} (1 instr vs 2). Forbidden inside IT. Use for null/loop guards.

\subsection{Sign-aware CCS (HW 11 always signed)}
Signed: \asm{GT/GE/LE/LT}; unsigned: \asm{HI/HS/LS/LO}. \textbf{Inverse pairs:} GT$\leftrightarrow$LE, GE$\leftrightarrow$LT, EQ$\leftrightarrow$NE. NL branches on the inverse. \textbf{Type drives 3 things:} CCS (GT/LE vs HI/LS), shift (ASR vs LSR), load (SB/SH vs B/H).

\subsection{Compound conditions}
\textbf{OR} (\cee{a||b}) -- short-circuit with TWO \asm{CMP}s, both branch to the SAME \asm{then} label; fall-through = else:
\begin{lstlisting}
    cmp r0, #20
    ble then         @ a true -> then
    cmp r0, #25
    bge then         @ b true -> then
    b   else_        @ neither -> else
\end{lstlisting}
\textbf{AND} (\cee{a\&\&b}) -- via De Morgan, branch to ELSE on inverse of either:
\begin{lstlisting}
    cmp r0, #LO
    blt else_        @ !a -> else (skip then)
    cmp r0, #HI
    bgt else_        @ !b -> else
    @ both true: fall through to then
\end{lstlisting}
For CEX of compound OR: 1st sub-test must be a plain \asm{Bcc} (cond B in middle of IT illegal); only the FINAL test can fan out into a full \asm{ITTEE}.

\subsection{Cond LDR/STR forms in IT}
All three addressing modes (\asm{[Rn,\#os]}, \asm{[Rn,\#os]!}, \asm{[Rn],\#os}) accept a CCS suffix (Lctr 24 Ex 3,4): e.g.\ \asm{ldrshgt r0,[r2,\#2]!}, \asm{ldrle r0,[r3],\#4}. Pre/post writeback honored only when cond true. Cond instr that fails: no write, no flag change, no exception.

\subsection{HW11 P1 -- \cee{x\_cmp\_y\_load}, three styles}
\begin{hwp}\textbf{P1.} \cee{int32\_t f(int32\_t x, int32\_t y, int16\_t *ptrA, int32\_t *ptrB)\{ if (x>y) return (int32\_t)*(++ptrA)*2; else return *(ptrB++)/2; \}}\\
AAPCS: R0=x, R1=y, R2=ptrA, R3=ptrB; ret in R0.\\
Then arm: \cee{*(++ptrA)} $\to$ \asm{ldrsh r0,[r2,\#2]!} (pre-idx, sign-ext halfword); \cee{*2} $\to$ \asm{lsl r0,r0,\#1}.\\
Else arm: \cee{*(ptrB++)} $\to$ \asm{ldr r0,[r3],\#4} (post-idx); signed \cee{/2} $\to$ \asm{asr r0,r0,\#1}.

\textbf{Part 1 -- PL (BGT, same CCS as C):}
\begin{lstlisting}
x_cmp_y_load_pl:
    cmp     r0, r1
    bgt     pl_then         @ x>y -> then
pl_else:                    @ else lays out FIRST
    ldr     r0, [r3], #4    @ *ptrB; ptrB++
    asr     r0, r0, #1      @ /2 signed
    b       pl_end
pl_then:
    ldrsh   r0, [r2, #2]!   @ ++ptrA; *ptrA (s-ext)
    lsl     r0, r0, #1      @ *2
pl_end:
    bx      lr
\end{lstlisting}

\textbf{Part 2 -- NL (BLE, inverse CCS, reads like C):}
\begin{lstlisting}
x_cmp_y_load_nl:
    cmp     r0, r1
    ble     nl_else         @ x<=y -> else (skip then)
nl_then:                    @ then falls through
    ldrsh   r0, [r2, #2]!
    lsl     r0, r0, #1
    b       nl_end
nl_else:
    ldr     r0, [r3], #4
    asr     r0, r0, #1
nl_end:
    bx      lr
\end{lstlisting}

\textbf{Part 3 -- CEX (\asm{ITTEE le}, no branches):}
\begin{lstlisting}
x_cmp_y_load_cex:
    cmp     r0, r1
    ittee   le              @ T=le,T=le,E=gt,E=gt
    ldrle   r0, [r3], #4    @ T1: else load (post-idx)
    asrle   r0, r0, #1      @ T2: else /2
    ldrshgt r0, [r2, #2]!   @ E1: then load (pre-idx)
    lslgt   r0, r0, #1      @ E2: then *2
    bx      lr
\end{lstlisting}
CCS in \asm{ITTEE} = \asm{le} (NL cond = inverse of C's \cee{>}). Each arm is 2 instrs $\Rightarrow$ \asm{ITTEE} (4 slots) is the natural fit. Cond pre-idx (\asm{ldrshgt!}) and post-idx (\asm{ldrle ,\#4}) both legal.
\end{hwp}

\subsection{HW11 P2 -- \cee{mathf} piecewise (3-way cascade)}
\begin{hwp}\textbf{P2.} \cee{f(x) = -10 if x<-5; +10 if x>=5; else 2x}.\\
\textbf{Part 1 (C):}
\begin{lstlisting}[language={}]
if (x < -5)      y = -10;
else if (x >= 5) y =  10;
else             y =  2*x;
\end{lstlisting}

\textbf{Part 2 -- CMP chain (NL at each arm):}
\begin{lstlisting}
mathf_cmp:
    cmp     r0, #-5         @ asm: cmn r0,#5
    bge     elseif          @ NL: x>=-5 skips 1st arm
    ldr     r0, =-10        @ x<-5 -> -10
    b       end
elseif:
    cmp     r0, #5
    blt     else_           @ NL: x<5 skips 2nd arm
    ldr     r0, =10         @ x>=5 -> +10
    b       end
else_:
    lsl     r0, r0, #1      @ 2*x
end:
    bx      lr
\end{lstlisting}
\textbf{Note:} \asm{cmp r0,\#-5} not directly encodable; assembler emits \asm{cmn r0,\#5} (flags = R0+5). Signed CCSs work unchanged.

\textbf{Part 3 -- CEX (chained \asm{ITT lt} + \asm{ITE ge}):}
\begin{lstlisting}
mathf_cex:
    cmp     r0, #-5
    itt     lt              @ T=lt, T=lt
    ldrlt   r1, =-10        @ T1: stash -10 in R1
    blt     end             @ T2: cond B *LAST slot only*
    @ fall-through: x >= -5
    cmp     r0, #5          @ fresh CMP, new flags
    ite     ge              @ T=ge, E=lt
    ldrge   r1, =10         @ T : x>=5, R1=10
    lsllt   r1, r0, #1      @ E : R1=2*x (uses ORIG x in R0)
end:
    mov     r0, r1          @ result -> ret reg
    bx      lr
\end{lstlisting}
\textbf{Key tricks:}
\begin{itemize}
\item \textbf{Stash result in R1} so R0 (=x) survives the 2nd CMP. \asm{lsllt r1,r0,\#1} reads original x.
\item \asm{blt end} is in \textbf{slot 2} of \asm{ITT} -- legal only because it's the \textbf{last} slot (Lctr 24 rule).
\item Without the cond B, fall-through after 1st IT would corrupt R1 in 2nd block.
\end{itemize}
\end{hwp}

\subsection{HW11 P3 -- compound OR \cee{x<=20 || x>=25}}
\begin{hwp}\textbf{P3.} \cee{int32\_t g(int x, uint16\_t *ptrA, uint32\_t *ptrB)\{ if (x<=20 || x>=25) return (int32\_t)*++ptrA*4; else return *++ptrB/4; \}}\\
AAPCS: R0=x, R1=ptrA, R2=ptrB; ret in R0.\\
Then arm: \cee{*++ptrA} $\to$ \asm{ldrh r0,[r1,\#2]!} (zero-ext); \cee{*4} $\to$ \asm{lsl r0,r0,\#2}.\\
Else arm: \cee{*++ptrB} $\to$ \asm{ldr r0,[r2,\#4]!}; UNSIGNED \cee{/4} $\to$ \asm{lsr r0,r0,\#2} (NOT \asm{ASR} -- ptrB is \cee{uint32\_t*}).

\textbf{Part 1 -- CMP short-circuit OR (both tests $\to$ same \cee{then}):}
\begin{lstlisting}
load_based_on_x_cmp:
    cmp     r0, #20
    ble     then            @ 1st OR true -> then
    cmp     r0, #25
    bge     then            @ 2nd OR true -> then
    b       else_           @ both false -> else
then:
    ldrh    r0, [r1, #2]!   @ ++ptrA; *ptrA (zero-ext)
    lsl     r0, r0, #2      @ *4
    b       end
else_:
    ldr     r0, [r2, #4]!   @ ++ptrB; *ptrB
    lsr     r0, r0, #2      @ /4 unsigned
end:
    bx      lr
\end{lstlisting}

\textbf{Part 2 -- CEX (1st test plain \asm{BGT}, 2nd test \asm{ITTEE ge}):}
\begin{lstlisting}
load_based_on_x_cex:
    cmp     r0, #20
    bgt     check25         @ x>20 -> 2nd test (plain B; mid-IT cond B illegal)
    @ x<=20: compound TRUE -> then arm inline
    ldrh    r0, [r1, #2]!
    lsl     r0, r0, #2
    bx      lr
check25:
    cmp     r0, #25
    ittee   ge              @ T=ge,T=ge,E=lt,E=lt
    ldrhge  r0, [r1, #2]!   @ T: then load
    lslge   r0, r0, #2      @ T: *4
    ldrlt   r0, [r2, #4]!   @ E: else load
    lsrlt   r0, r0, #2      @ E: /4
    bx      lr
\end{lstlisting}
\textbf{Why split:} 1st sub-test of OR can't sit inside an IT block -- it would need a cond B in a non-last slot. Only the FINAL test packs into \asm{ITTEE}. Both arms run 2 instrs $\Rightarrow$ \asm{ITTEE} (T,T,E,E) is exact fit.

\textbf{Trace:}
\begin{tabular}{@{}c l l@{}}
\toprule
x & path & arm \\
\midrule
15 & \asm{bgt} N (15$\le$20) & inline then \\
20 & \asm{bgt} N            & inline then \\
21 & \asm{bgt} Y, ge N      & ITTEE E (else) \\
24 & \asm{bgt} Y, ge N      & ITTEE E (else) \\
25 & \asm{bgt} Y, ge Y      & ITTEE T (then) \\
30 & \asm{bgt} Y, ge Y      & ITTEE T (then) \\
\bottomrule
\end{tabular}
\end{hwp}

\subsection{Picker + pitfalls}
\textbf{Pick CEX} when both arms $\le$ 4 instrs total; \textbf{PL/NL} when arms larger or need general branch; \textbf{CMP short-circuit} for OR; \textbf{CMP+De Morgan} for AND (branch to else on $!a$/$!b$); \textbf{CBZ/CBNZ} for null/loop guards (forward $\le$126B, R0-R7, no flag effect).
\textbf{Gotchas:} (1) HI/LS on signed (or GT/LE on unsigned) miscompares near sign bit -- match CCS to type. (2) LSR on negative = huge positive; use ASR for signed div. (3) LDRH on \cee{int16\_t} loses sign -- use LDRSH. (4) Cond B in non-last IT slot = assembler err. (5) CBZ inside IT = forbidden. (6) In CEX cascade, stash result in R1 so R0 (=x) survives next CMP. (7) \asm{cmp Rn,\#-imm} rewritten to \asm{cmn Rn,\#imm} (signed CCS unchanged). (8) Pre/post-idx LDR both legal in IT, writeback honored only when cond true.

\subsection{One-line summary cards}
\begin{itemize}
\item \textbf{PL:} \asm{cmp; B<C-cond> then; <else>; b end; then: <then>; end:}
\item \textbf{NL:} \asm{cmp; B<inv-cond> else; <then>; b end; else: <else>; end:}
\item \textbf{CEX:} \asm{cmp; ITTEE <inv-cond>; <else x2 with cond>; <then x2 with inv-cond>}
\item \textbf{CEX cascade:} \asm{cmp; ITT lt; <stash-T>; blt end; cmp; ITE ge; <stash-T>; <stash-E>; end: mov r0,r1}
\item \textbf{OR (CMP):} \asm{cmp; ble then; cmp; bge then; b else}
\item \textbf{OR (CEX):} 1st = plain \asm{Bcc} to \asm{checkN}, 2nd packs into \asm{ITTEE}.
\end{itemize}

% ----- part1/09_hw12_loops.tex -----
\section{HW12: Loops in ASM}

\subsection{C loops $\to$ ASM shapes (Lctr 26 \S26.1--26.2)}
\begin{tabular}{@{}l l l l@{}}
\textbf{loop} & \textbf{test} & \textbf{min iters} & \textbf{C $\to$ ASM idiom} \\
\asm{while(c)} & top  & 0 & priming-read style: dup body update \\
\asm{do\{\}while(c)} & bottom & 1 & body, then \asm{cmp}+cond.\,branch back \\
\asm{for(i;c;u)} & top & 0 & \asm{init} / lp: \asm{cmp}+exit / body / \asm{u} / \asm{b lp} \\
\asm{while(1)\{...break;\}} & inline & 0+ & single in-loop \asm{LDRB}, \asm{break} = \asm{b \_end} \\
\end{tabular}

\textbf{Label conventions:} \asm{\_lp} at the top (re-entry target); \asm{\_end} after the exit (target of \asm{break}/condition-fail). \textbf{Top-tested while $\ge 0$ iters; do-while $\ge 1$ iter.} \asm{for} = same shape as \asm{while} with explicit init+update slots.

\subsection{Plain \asm{while} (priming-read style)}
\begin{lstlisting}
@ ch = *src++;                    @ priming LDRB BEFORE label
@ while (ch) { body; ch = *src++; }
        ldrb r2, [r1], #1         @ priming read
fn_lp:  cbz  r2, fn_end           @ exit when ch==0
            @ ... body ...
            ldrb r2, [r1], #1     @ in-loop update (DUP of priming)
        b    fn_lp
fn_end: bx   lr
\end{lstlisting}
Two textual \asm{LDRB}s (priming + update). \asm{CBZ} on the freshly-loaded byte tests the null-terminator without setting flags.

\subsection{\asm{do-while} (bottom-tested, body $\ge 1$)}
\begin{lstlisting}
@ do { body; i++; } while (i < N);
        mov  r1, #0               @ i = 0
fn_lp:                            @ no entry test -> body always runs once
            @ ... body ...
            add  r1, r1, #1
        cmp  r1, r2
        blt  fn_lp                @ back-edge test at bottom
fn_end: bx   lr
\end{lstlisting}
No \asm{\_end} jump-around needed; first entry falls straight into body.

\subsection{\asm{for} loop}
\begin{lstlisting}
@ for (i = 0; i < N; i++) { body }
        mov  r1, #0               @ INIT
fn_lp:  cmp  r1, r2               @ TEST at top
        bge  fn_end               @ exit if !cond
            @ ... body ...
            add  r1, r1, #1       @ UPDATE
        b    fn_lp
fn_end: bx   lr
\end{lstlisting}

\subsection{\asm{while(1)} + \asm{break} (cleaner than priming)}
\begin{lstlisting}
@ while (1) { ch = *src++; if (!ch) break; body; }
fn_lp:  ldrb r2, [r1], #1         @ SINGLE in-loop read (no priming)
        cbz  r2, fn_end           @ break = forward b end
            @ ... body ...
        b    fn_lp
fn_end: bx   lr
\end{lstlisting}
\textbf{Dynamic-count rule:} priming-read and \asm{while(1)+break} execute the \emph{same} number of \asm{LDRB}s at runtime ($N{+}1$ for an $N$-element string). The break-style only reduces \emph{source-line} count by removing the duplicated \asm{LDRB}.

\subsection{\asm{while}-with-CEX (rare; body must fit IT)}
\begin{lstlisting}
fn_lp:  cmp  r1, r2
        itt  lt
        ldrlt r0, [r3], #4        @ body in IT
        blt  fn_lp                @ cond B legal ONLY as last IT slot
\end{lstlisting}
Only viable when the entire body (and the back-edge \asm{B}) fits in 4 IT slots. Otherwise use plain \asm{cmp}+\asm{bcc}.

\subsection{\asm{break} / \asm{continue} semantics (Lctr 27 \S27.2--27.3)}
\textbf{break} $\to$ unconditional forward \asm{b \_end}; must be wrapped in an \asm{if} so it's executed conditionally (typically via \asm{CBZ}/\asm{CBNZ}/\asm{Bcc}).

\textbf{\asm{continue} target depends on loop type:}
\begin{tabular}{@{}l l@{}}
in \asm{while}/\asm{do-while} & jumps to \textbf{condition test} (\asm{\_lp}) \\
in \asm{for}                  & jumps to the \textbf{UPDATE expr} (i++ still runs) \\
\end{tabular}
$\Rightarrow$ in a while loop, any loop-control update must appear \emph{before} \asm{continue}; otherwise the loop hangs.

\subsection{HW12 P1 -- swap last 2 of every trio}
\begin{hwp}\textbf{P1.} \cee{int mp\_swap\_last\_2\_chars\_in\_3\_in\_str(char *str)} -- walk 3 chars at a time; if all 3 non-null, swap the last two; return \# trios swapped.\\
\textbf{Test:} \cee{"hellow"$\to$"hlelwo"} ret 2; \cee{"hello"$\to$"hlelo"} ret 1; \cee{"he"} ret 0.

\textbf{C (while(1)+break, three-way exit):}
\begin{lstlisting}
int n = 0;
while (1) {
    char c1 = *str++; if (!c1) break;
    char c2 = *str++; if (!c2) break;
    char c3 = *str++; if (!c3) break;
    *(str-2) = c3;  *(str-1) = c2;  n++;
}
return n;
\end{lstlisting}

\textbf{ASM} (\asm{R12} accumulator $\Rightarrow$ no push/pop; \asm{r1/r2/r3} = c1/c2/c3 scratch):
\begin{lstlisting}
        mov  r12, #0               @ pair_swapped = 0
fn_lp:  ldrb r1, [r0], #1          @ c1 = *str++  (post-idx)
        cbz  r1, fn_end            @ break if !c1
        ldrb r2, [r0], #1          @ c2 = *str++
        cbz  r2, fn_end            @ break if !c2
        ldrb r3, [r0], #1          @ c3 = *str++
        cbz  r3, fn_end            @ break if !c3
            strb r3, [r0, #-2]     @ *(str-2) = c3
            strb r2, [r0, #-1]     @ *(str-1) = c2
            add  r12, r12, #1      @ n++
        b    fn_lp
fn_end: mov  r0, r12               @ return n
        bx   lr
\end{lstlisting}

\textbf{Why \asm{[r0,\#-2]}/\asm{[r0,\#-1]} (not $-3,-2$):} after 3 post-inc \asm{LDRB}s, \asm{r0} sits one byte \emph{past} the trio. The trio occupies addrs \asm{[r0,\#-3]} (c1), \asm{[r0,\#-2]} (c2), \asm{[r0,\#-1]} (c3). We swap \emph{only the last two} $\Rightarrow$ $-2$ and $-1$. Using $-3,-2$ corrupts c1.

\textbf{Why no \asm{!} writeback on the STRBs:} we want \asm{r0} to remain at the post-trio cursor for the next iter; \asm{[r0,\#-2]} is plain base+neg-offset, no update.
\end{hwp}

\subsection{HW12 P2/P1 -- lower$\to$upper, plain \asm{while}}
\begin{hwp}\textbf{P2.1.} \cee{void mp\_str\_cpy\_2\_caps\_while(char *dst,char *src);} -- ASM of the given priming-read C.\\
\textbf{C (given):}
\begin{lstlisting}
char ch = *src++;          // priming read
while (ch) {
    ch -= 0x20;            // 'a'=0x61 -> 'A'=0x41
    *dst++ = ch;
    ch = *src++;           // duplicate update
}
*dst++ = ch;               // post-loop store of '\0'
\end{lstlisting}

\textbf{Regs (AAPCS):} r0=dst, r1=src, r2=ch. All scratch $\Rightarrow$ no push/pop.

\textbf{ASM:}
\begin{lstlisting}
        ldrb r2, [r1], #1          @ priming read
fn_lp:  cbz  r2, fn_end            @ while (ch)
            sub  r2, r2, #0x20     @ ch -= 0x20
            strb r2, [r0], #1      @ *dst++ = ch
            ldrb r2, [r1], #1      @ ch = *src++  (DUP)
        b    fn_lp
fn_end: strb r2, [r0], #1          @ *dst++ = '\0'  (r2=0 from CBZ exit)
        bx   lr
\end{lstlisting}
\textbf{Note:} \asm{0x20} is a Thumb-2 8-bit-rotated immediate $\Rightarrow$ \asm{sub} encodes directly. Terminator store at \asm{\_end} uses \asm{r2=0} guaranteed by \asm{CBZ}-exit. Writeback (\asm{\#1}) on terminator \asm{strb} is harmless -- \asm{r0} is dead.
\end{hwp}

\subsection{HW12 P2/P2 -- refactor to \asm{while(1)}+\asm{break} (C only)}
\begin{hwp}\textbf{P2.2.} Eliminate the duplicate \asm{*src++} read.
\begin{lstlisting}
while (1) {
    char ch = *src++;
    if (ch == 0) {
        *dst++ = ch;       // store '\0' BEFORE break
        break;
    }
    ch -= 0x20;
    *dst++ = ch;
}
\end{lstlisting}
Single read at top of body; \asm{break} after writing the terminator.\\
\textbf{Equivalence:} non-null bytes $\to$ read,sub,store; null byte $\to$ store-zero,exit. Identical observable behavior to plain \asm{while} version.
\end{hwp}

\subsection{HW12 P2/P3 -- \asm{while(1)}+\asm{break} ASM}
\begin{hwp}\textbf{P2.3.} ASM of the refactored C.
\begin{lstlisting}
fn_lp:  ldrb r2, [r1], #1          @ ch = *src++  (SINGLE read)
        cbz  r2, fn_end            @ if (!ch) goto end (deferred terminator)
            sub  r2, r2, #0x20
            strb r2, [r0], #1
        b    fn_lp
fn_end: strb r2, [r0], #1          @ *dst++ = '\0'  (r2=0 by CBZ)
        bx   lr
\end{lstlisting}
\textbf{Where the terminator store lives:} the C writes \asm{*dst++ = ch} \emph{inside} the \asm{if(!ch)} branch \emph{before} \asm{break}. Translating literally, that store must still happen $\Rightarrow$ place the \asm{strb} as the \emph{first} instruction at \asm{\_end}. Don't duplicate it inside the loop.\\
\textbf{Source-line vs dynamic count:} only one textual \asm{ldrb} (vs two in P2/P1). Runtime: same $N{+}1$ \asm{LDRB} executions per call. The refactor saves \emph{source bytes}, not cycles.
\end{hwp}

\subsection{Loop dynamic-count rule (memorize)}
For a string of $N$ non-null chars then \asm{\textbackslash 0}: priming-read \asm{while} and \asm{while(1)+break} both execute exactly $N{+}1$ \asm{LDRB}s at runtime (one per byte read, including the terminator). \textbf{Refactor only collapses source duplication; it does not save instructions executed.}

\subsection{Pointer-idiom $\to$ ASM (byte) cheat-card}
\begin{tabular}{@{}l l l@{}}
\cee{x = *p++}      & \asm{ldrb r0, [r1], \#1}    & post-idx, writeback \\
\cee{*p++ = x}      & \asm{strb r0, [r1], \#1}    & post-idx, writeback \\
\cee{x = *++p}      & \asm{ldrb r0, [r1, \#1]!}   & pre-idx, writeback \\
\cee{*++p = x}      & \asm{strb r0, [r1, \#1]!}   & pre-idx, writeback \\
\cee{*(p-1) = x}    & \asm{strb r0, [r1, \#-1]}   & base+neg-offset, no writeback \\
\cee{*(p-2) = x}    & \asm{strb r0, [r1, \#-2]}   & base+neg-offset, no writeback \\
\end{tabular}

\subsection{\asm{CBZ}/\asm{CBNZ} -- when, why, restrictions (Lctr 25)}
\textbf{\asm{CBZ Rn,lbl}} = $\{$\asm{cmp Rn,\#0}; \asm{beq lbl}$\}$ in one instr; \textbf{leaves flags untouched}.
\begin{itemize}
\item Ideal for null-terminator test on a fresh \asm{LDRB} (no flags wanted).
\item \asm{Rn} restricted to \asm{r0..r7}.
\item Forward branch only ($\le {+}126$ B); \emph{cannot} branch backward.
\item \textbf{Cannot appear inside an IT block} (Lctr 24 \S24.3.2).
\end{itemize}

\subsection{\asm{switch} implementations (Lctr 27 \S27.4)}
HW12 doesn't exercise \asm{switch}, but it's exam-fair. Three forms:

\textbf{(a) \asm{CMP}/\asm{BEQ} ladder} -- $O(n)$, one cmp per case:
\begin{lstlisting}
        cmp r0, #0
        beq case0
        cmp r0, #1
        beq case1
        ...
        b   default
\end{lstlisting}

\textbf{(b) Stacked-label fall-through} -- multiple values share a body by stacking labels:
\begin{lstlisting}
case0:
case1:                          @ both 0 and 1 fall here
        @ shared body
        b   end_sw
\end{lstlisting}

\textbf{(c) \asm{TBB} jump table} -- $O(1)$ via byte offsets:
\begin{lstlisting}
        adr  r1, jt
        tbb  [r1, r0]           @ PC += 2 * jt[r0]
jt:     .byte (case0 - jt)/2
        .byte (case1 - jt)/2
        .byte (case2 - jt)/2
        .byte (default - jt)/2
\end{lstlisting}
\textbf{\asm{TBH}} variant uses halfword offsets for far cases: \asm{tbh [r1, r0, lsl \#1]}, \asm{.short} entries. Use TBB for tight tables ($\le 510$ B reach), TBH for larger ones.

\subsection{Common HW12 mistakes}
\begin{itemize}
\item Wrong offsets in P1: \asm{\#-3,\#-2} instead of \asm{\#-2,\#-1} -- corrupts c1.
\item Putting the terminator \asm{strb} \emph{inside} the loop in P2/P3 -- duplicates the C's \cee{*dst++=ch;break} pattern; cleanest is at \asm{\_end}.
\item Using \asm{cmp Rn,\#0}+\asm{beq} when \asm{cbz} suffices -- wastes 2 bytes / blows away flags you may need.
\item Pushing \asm{r4} for a tiny loop accumulator -- prefer \asm{r12} (IP, scratch under AAPCS).
\item Trying \asm{cbz} inside an IT block -- forbidden.
\item Forgetting that \cee{while(ch)} on an empty input runs body 0 times -- the post-loop \cee{*dst++=ch} still writes the \asm{\textbackslash 0} (correct).
\item For-loop \asm{continue} branches to UPDATE, not the test -- so \asm{i++} still runs. While-loop \asm{continue} branches to test.
\end{itemize}

\subsection{Quick-reference summary}
\begin{itemize}
\item \textbf{Top-tested} (\asm{while}, \asm{for}) $\ge 0$ iters; \textbf{bottom-tested} (\asm{do-while}) $\ge 1$.
\item \textbf{Priming \asm{while}} = duplicated \asm{LDRB} (text); \textbf{\asm{while(1)+break}} = single \asm{LDRB}; same dynamic count.
\item Labels: \asm{\_lp} (top, re-entry) and \asm{\_end} (post-exit, terminator-store target).
\item \asm{CBZ} on fresh \asm{LDRB} = idiomatic null-terminator test; not allowed in IT.
\item Post-idx \asm{ldrb [r0],\#1} $=$ \cee{*p++}; base+neg-offset \asm{[r0,\#-k]} $=$ \cee{*(p-k)} no update.
\item \asm{R12} (IP) is AAPCS scratch -- prefer over \asm{R4} for short-lived loop locals.
\item \asm{break} $\to$ forward \asm{b \_end}; \asm{continue} target = \asm{for}'s update / \asm{while}'s test.
\item \asm{switch}: ladder ($O(n)$) / stacked fall-through (shared body) / \asm{TBB},\asm{TBH} jump table ($O(1)$).
\end{itemize}

% ----- part1/10_quiz3.tex -----
% =============================================================================
% Quiz 3 --- early-ASM material (EABI, stack, barrel-shifter, BIC+ORR splice)
% NOTE: no standalone qz3 problem/solution PDF exists in the course repo. The
% problems below are reconstructed from
%   content/cec_320/quizes/quiz3/qz3_cheatsheet.md
%   content/cec_320/quizes/quiz3/qz3_init.md
%   content/cec_320/quizes/quiz3/lecture_notes_summary.md (Lctr 14--18)
%   content/cec_320/final_exam/quiz3_quiz4_extracts.md (Q3 P1..P6)
% Header in qz3_init.md is explicit: ``Just a reminder that you can use your own
% formula sheets for QZ3'' --- so the user originally also requested NVIC/EXTI
% material here, but Quiz 3 actually covers HW6/7/8 (Lctr 14--19): mixed C/ASM,
% EABI, arithmetic, shifts, bitwise logic, NZCV, CCS. No interrupt/EXTI/NVIC
% content was tested on Quiz 3 per the source markdown.
% =============================================================================
\section{Quiz 3: Interrupts + Early ASM}

\subsection{Coverage \& meta}
\textbf{Scope:} HW6 (mixed C+ASM, EABI, stack, 64-bit \asm{ADDS/ADC}, \asm{SUBS/SBC}, \asm{UDIV+MLS} modulo, Q15.16 encode), HW7 (shifts, barrel-shifter scaling, BIC/ORR bit-field splice), HW8 (NZCV, CCS, neg-logic branch, \asm{SMLAL} loop). Lectures 14--19. \emph{Despite the section title above, no NVIC/EXTI was on Q3 ---} that material lives on Q4/Q5. The ``early ASM'' tag is the accurate description; cross-reference \S HW6/7/8 above for full mechanics.\\
\textbf{Allowed:} prior cheat-sheets + one new letter-size sheet, only ASM instructions covered in lecture, no invented mnemonics.\\
\textbf{Reconstruction note:} Q3 problem statements below are derived from the cheatsheet+extract markdown; the original Fall-2025 PDF is not in the repo, so wording may differ slightly. Numeric answers are verified against the cheatsheet trace.

\subsection{One-page Q3 cookbook}
\begin{tabular}{@{}ll@{}}
\toprule
task & one-liner \\
\midrule
return signed $-2$ in 32b   & \cee{R0 = 0xFFFF\_FFFE} (two's-comp), \emph{not} 0x80000002 \\
sign-ext \cee{int8\_t} arg  & MSB$=$1 $\Rightarrow$ pad upper 24b with 1s; MSB$=$0 $\Rightarrow$ pad with 0s \\
zero-ext \cee{uint8\_t}     & always pad upper bits with 0s \\
\cee{push \{list\}} layout  & lowest reg\# at lowest addr (independent of brace order) \\
\cee{push} of $n$ regs      & \cee{SP -= 4*n} (full-descending) \\
$k\cdot B/2^n$ in 1 instr   & decompose $k/2^n=1\pm 1/2^m$, use \cee{ADD/SUB Rd,B,B,ASR\#m} \\
clear field then insert     & \cee{BIC Rd,Rd,mask\,LSL\#s}; \cee{ORR Rd,Rd,val\,LSL\#s} \\
\cee{ADDS/ADC} 64-bit add   & low \cee{ADDS} sets C; \cee{ADC} adds C into high \\
\bottomrule
\end{tabular}

\subsection{P1 (7pt) --- Return reg for signed int32}
\textbf{Setup:} ASM function returns signed 32-bit $-2$. (a) Which register holds it? (b) Hex value?
\begin{qp}
(a) \textbf{R0} (AAPCS: any 32-bit return goes in R0; 64-bit return uses R0:R1 lo:hi). (b) $-2$ in 32-bit two's-complement: invert bits of $+2 = $\cee{0x0000\_0002} and add 1 $\to$ \textbf{\cee{0xFFFF\_FFFE}}. \emph{Trap:} writing \cee{0x8000\_0002} (sign-magnitude form, which ARM does \emph{not} use) loses 6 of 7 points.
\end{qp}

\subsection{P2 (20pt) --- 64-bit ADD + C flag}
\textbf{Setup:} \asm{mp\_add\_64bit\_s} executes \asm{adds r0,r2}; \asm{adc r1,r3}; \asm{bx lr}. Two test sets: $A_1=$\cee{0x377{<}{<}24}, $B_1=$\cee{0x287{<}{<}24}; $A_2=$\cee{0x378{<}{<}24}, $B_2=$\cee{0x288{<}{<}24}. Find sums and the C flag observed after the low-word \asm{adds}.
\begin{qp}
Layout: \asm{R1:R0} holds A, \asm{R3:R2} holds B. Bits 24..31 of the low word receive the most-significant nibble of the constant.\\
\textbf{Set 1:} low words \cee{0x77\_0000\_00 + 0x87\_0000\_00 = 0xFE\_0000\_00} (no carry-out of bit 31) $\Rightarrow$ \textbf{C=0}. High words 0+0+0=0. Sum $=$ \cee{(int64\_t)0x5FE{<}{<}24}.\\
\textbf{Set 2:} low \cee{0x78\_0000\_00 + 0x88\_0000\_00 = 0x1\_00\_0000\_00} (overflows bit 31!) $\Rightarrow$ \textbf{C=1}, low result \cee{0x00\_0000\_00}. \asm{adc} adds C to high: $0+0+1=1$. Sum $=$ \cee{(int64\_t)0x600{<}{<}24}.\\
\emph{Design rule:} a good 64-bit-ADD test toggles the low-word carry between cases (one set with C=0, one with C=1) so \asm{adc} is exercised both ways.
\end{qp}

\subsection{P3 (20pt) --- \texttt{push \{r4,r8,r6\}} memory layout}
\textbf{Setup:} $R_i = (1{<}{<}i) + (1{<}{<}(2{+}i))$ for $i=1..12$. \asm{SP} $=$ \cee{0x2000\_1020} before \asm{push \{r4, r8, r6\}}. Find SP after, $R_4$, the word at \cee{0x2000\_1018}, and the byte at \cee{0x2000\_1019}.
\begin{qp}
3 regs $\cdot$ 4B $=$ 12 $=$ \cee{0xC} $\Rightarrow$ \textbf{SP$=$\cee{0x2000\_1014}}. Push writes \emph{in ascending register-number order} regardless of brace listing $\Rightarrow$ R4 lowest, R6 middle, R8 highest:\\
\quad \cee{0x2000\_1014} $\leftarrow$ \textbf{R4} = \cee{(1{<}{<}4)+(1{<}{<}6) = 0x10+0x40 = 0x0000\_0050}\\
\quad \cee{0x2000\_1018} $\leftarrow$ \textbf{R6} = \cee{(1{<}{<}6)+(1{<}{<}8) = 0x40+0x100 = 0x0000\_0140}\\
\quad \cee{0x2000\_101C} $\leftarrow$ \textbf{R8} = \cee{(1{<}{<}8)+(1{<}{<}10) = 0x100+0x400 = 0x0000\_0500}\\
Word at \cee{0x2000\_1018} $=$ \textbf{\cee{0x0000\_0140}}. Byte at \cee{0x2000\_1019}: little-endian, address +1 holds byte 1 of R6 $=$ \textbf{\cee{0x01}} (next bytes: \cee{0x00},\cee{0x00},\cee{0x00} for the upper word).\\
\emph{Trap:} writing in brace-list order (R4,R8,R6) and putting R8 at the middle address --- the assembler ignores the listing order.
\end{qp}

\subsection{P4 (25pt) --- Single-instr $k\cdot B/8$ via barrel shifter}
\textbf{Setup:} signed $A$ in \asm{r0}, signed $B$ in \asm{r1}. Each part wants $A=k\cdot B/8$ in one instruction. Decompose as $k/8 = 1 \pm 2^{-m}$ where convenient; remainder $= 1$ becomes a no-op MOV.
\begin{center}\scriptsize
\begin{tabular}{@{}clll@{}}
\toprule
$k$ & $k/8$ & decomp & instruction \\
\midrule
6  & 6/8  & $1 - 2^{-2}$ & \asm{sub r0, r1, r1, asr \#2} \\
7  & 7/8  & $1 - 2^{-3}$ & \asm{sub r0, r1, r1, asr \#3} \\
8  & 1    & $1$          & \asm{mov r0, r1} \;\,(or \asm{add r0, r1, \#0}) \\
9  & 9/8  & $1 + 2^{-3}$ & \asm{add r0, r1, r1, asr \#3} \\
10 & 10/8 & $1 + 2^{-2}$ & \asm{add r0, r1, r1, asr \#2} \\
\bottomrule
\end{tabular}
\end{center}
\begin{qp}
\textbf{Why ASR not LSR?} $B$ is signed; \asm{LSR} fills MSB with 0, mangling negatives ($-1{\,>>\,}1$ would become $\!\sim\!2^{31}$). \asm{ASR} replicates the sign bit. \textbf{Why no \asm{rsb}?} $1\pm 2^{-m}$ keeps the unshifted $B$ first, so the sign of $B$ leads the result; \asm{rsb} would only be needed for $-1+\text{something}$ patterns ($k\cdot B$ with $k$ near 0). \textbf{Trap:} \asm{add r0, r0, r0, asr \#3} silently corrupts when \asm{r0=r1} aliases the input.
\end{qp}

\subsection{P5 (20pt) --- BIC+ORR bit-field splice $\to$ C}
\textbf{Setup:} ASM body
\begin{lstlisting}
mp_func_s:
    ldr  r3, =15           @ mask = 0xF (4 ones)
    bic  r0, r0, r3, lsl #8   @ a &= ~(0xF << 8)   ; clear bits[11:8]
    and  r1, r1, r3           @ b &= 0xF           ; isolate b's low nibble
    orr  r0, r0, r1, lsl #8   @ a |= b << 8        ; insert into [11:8]
    bx   lr
\end{lstlisting}
Write the C equivalent and trace with $a=$\cee{0xH3H2H1H0}, $b=17=$\cee{0x0011}.
\begin{qp}
C version (canonical clear-then-insert):
\begin{lstlisting}[language=]
uint16_t func_c(uint16_t a, uint16_t b) {
    uint16_t mask = 15;
    a &= ~(mask << 8);          // clear nibble at bit 8
    a |= (b & mask) << 8;       // insert low nibble of b
    return a;
}
\end{lstlisting}
\textbf{Trace} with $a=$\cee{0xH3H2H1H0}, $b=$\cee{0x0011}: \asm{r3=0x000F}; \asm{r3{<}{<}8 = 0x0F00}; \asm{bic} $\Rightarrow a=$\cee{0xH3\_0\_H1H0} (digit H2 cleared). $b\,\&\,0xF = 0x1$; \asm{orr} with \cee{0x1{<}{<}8 = 0x0100} $\Rightarrow$ \textbf{$a=$\cee{0xH3\,1\,H1H0}}. \emph{Pattern alias:} this is the universal ``\asm{BIC}+\asm{ORR} both shifted by the same \asm{LSL\#s}'' field-splice idiom; if you see \asm{BIC ...,LSL\#s} look immediately for the matching \asm{ORR ...,LSL\#s} below.
\end{qp}

\subsection{P6 (8pt) --- \cee{int8\_t} sign-ext into 32-bit args}
\textbf{Setup:} \cee{int8\_t var1=0x7E, var2=0x81; fn(var1,var2);} where prototype takes \cee{int32\_t}. Find R0, R1.
\begin{qp}
Sign-extension occurs at the \emph{call site} per AAPCS, before the callee runs.\\
\cee{var1=0x7E}: MSB (bit 7) $=$ 0 $\Rightarrow$ positive ($+126$); pad upper 24b with 0 $\Rightarrow$ \textbf{R0$=$\cee{0x0000\_007E}}.\\
\cee{var2=0x81}: MSB $=$ 1 $\Rightarrow$ negative ($-127$ as \cee{int8\_t}); pad upper 24b with 1 $\Rightarrow$ \textbf{R1$=$\cee{0xFFFF\_FF81}}.\\
\emph{Sanity:} \cee{0xFFFF\_FF81} as \cee{int32\_t} is $-127$ --- value preserved across the widening. Compare to \cee{uint8\_t} promotion (always zero-ext): \cee{0x81} would become \cee{0x0000\_0081} ($=129$). The same sign-ext rule governs \asm{LDRSB} / \asm{LDRSH} loads (Q4/Q6a territory).
\end{qp}

\subsection{Cross-references --- prior-week material on Q3}
\textbf{EABI registers:} R0--R3 args/scratch (caller-saved); R4--R11 callee-saved; R12=IP (scratch); R13=SP, R14=LR, R15=PC. 64-bit return: R0=lo, R1=hi. \textbf{Stack:} full-descending, SP points to last pushed item, \asm{push} pre-decrements, \asm{pop} post-increments. \textbf{Flag-only ops:} \asm{CMP=SUBS}-discard, \asm{CMN=ADDS}-discard, \asm{TST=ANDS}-discard (no V update), \asm{TEQ=EORS}-discard. \textbf{C-on-sub:} ARM sets C$=$!borrow (opposite of x86). \textbf{CCS domains:} \asm{HI/HS/LO/LS} unsigned (use C+Z), \asm{GT/GE/LT/LE} signed (use V vs N), \asm{EQ/NE} domain-free (just Z). \textbf{Neg-logic branch:} after \asm{cmp}, branch on the \emph{opposite} CCS to skip the then-block.

\subsection{Q3 trap inventory (memorize before walking in)}
\begin{itemize}
  \item Negative int return $\to$ two's-complement hex, \emph{not} sign-magnitude (P1).
  \item \asm{push \{r4,r8,r6\}} stores in numeric order; brace-list order is irrelevant (P3).
  \item 64-bit add: low \asm{adds} provides C, high \asm{adc} consumes it; if your test only exercises one C value, \asm{adc} is half-untested (P2).
  \item Barrel-shifter ASR vs LSR --- ASR for signed, LSR for unsigned. Mismatch $=$ silent corruption (P4).
  \item \asm{BIC \dots LSL\#s} immediately followed by \asm{ORR \dots LSL\#s} $=$ ``replace field [s+w-1:s]'' --- treat as a single semantic unit (P5).
  \item \cee{int8\_t} args $\geq$ \cee{0x80} sign-extend to \cee{0xFFFF\_FFxx}; \cee{uint8\_t} always zero-extends. Same rule for \asm{ldrsb} on Q4 (P6).
  \item C $=$ NOT borrow on subtract (ARM convention). On 5-bit hand-flag drills, this flips most students up.
  \item \asm{TST}/\asm{TEQ} don't update V $\to$ never pair with V-dependent CCS.
\end{itemize}

% ----- part1/11_quiz4_qz26a.tex -----
% =============================================================================
% Quiz 4 (qz-26a, Spring 2026) + practice qz-25b qz5/qz6a
% Topics: LDR/STR (signed/half/byte), pre/post-index +writeback, scaled-reg
%         addressing, LDRD/STRD, conditional-branch flow control, ptr idioms.
% =============================================================================
\section{Quiz 4: LDR/STR + Flow Control (qz-26a current sem)}

\subsection{Cheat-card: LDR/STR variants}
\textbf{Sign-extend loads} fill upper bits with the loaded MSB:
\begin{center}\scriptsize
\begin{tabular}{@{}llll@{}}
\toprule
mnem & loads & upper bits & view \\
\midrule
\asm{LDRB}  & 1 byte  & zero-fill   & unsigned \\
\asm{LDRSB} & 1 byte  & sign-extend & signed   \\
\asm{LDRH}  & 2 bytes & zero-fill   & unsigned \\
\asm{LDRSH} & 2 bytes & sign-extend & signed   \\
\asm{LDR}   & 4 bytes & ---         & word     \\
\asm{LDRD}  & 8 bytes & Rt=lo,Rt2=hi& pair     \\
\bottomrule
\end{tabular}
\end{center}
Stores have no signed variant (\asm{STRB/STRH/STR/STRD} only) -- low bits go to memory regardless of sign.

\textbf{Addressing modes} (effective addr {[}EA{]} vs writeback to Rn):
\begin{center}\scriptsize
\begin{tabular}{@{}lll@{}}
\toprule
syntax & EA & Rn after \\
\midrule
\asm{[Rn]}                & Rn          & Rn (unchanged) \\
\asm{[Rn,\#i]}            & Rn$+$i      & Rn (offset) \\
\asm{[Rn,\#i]!}           & Rn$+$i      & Rn$+$i \,\,(\textbf{pre-index} wb) \\
\asm{[Rn],\#i}            & Rn          & Rn$+$i \,\,(\textbf{post-index} wb) \\
\asm{[Rn,Rm,LSL \#n]}     & Rn$+$(Rm$\!\ll\!$n) & Rn (scaled-reg) \\
\bottomrule
\end{tabular}
\end{center}
\textbf{Pre} uses updated addr; \textbf{post} uses old addr then advances. Post-index by 4 advances Rn even if only 2 bytes were read (e.g.\ \asm{ldrsh r3,[r0],\#4}).

\textbf{Scaled register}: \asm{LDR Rt,[Rn,Rm,LSL \#n]} where \asm{n}$=\log_2$sizeof. \asm{LSL \#0} byte, \asm{\#1} half, \asm{\#2} word, \asm{\#3} dword. R1 is in \emph{elements}; the shift converts to bytes. Halfword/word access at non-aligned addr faults.

% =============================================================================
\subsection{Q4 (qz-26a, current sem)}
\textbf{Memory state for Prob 1} (little-endian; bytes labeled B0..B3):
\begin{center}\scriptsize
\begin{tabular}{@{}lcccc@{}}
\toprule
Addr & B0 & B1 & B2 & B3 \\
\midrule
0x2000\_0000 & 88 & 89 & 8A & 8B \\
0x2000\_0004 & 8C & 8D & 8E & 8F \\
0x2000\_0008 & 70 & 71 & 72 & 73 \\
0x2000\_000C & 74 & 75 & 76 & 77 \\
\bottomrule
\end{tabular}
\end{center}
Half/word reads are LE: addr 0x4 halfword $=$ \asm{0x8D8C} (B1:B0); word $=$ \asm{0x8F8E\_8D8C}.

\begin{qp}\textbf{Prob 1a (8pt) -- task1 just before \asm{bx lr}:}
\begin{lstlisting}
task1:                       @ R0=0x2000_0000, R1=0x2000_0020
    ldrsh r2, [r0, #4]       @ HW @0x4 = 0x8D8C (neg)
                             @ -> R2 = 0xFFFF_8D8C
    str   r2, [r1], #4       @ store full r2 to 0x20, R1+=4
    ldrsh r3, [r0, #8]       @ HW @0x8 = 0x7170 (pos)
                             @ -> R3 = 0x0000_7170
    str   r3, [r1], #4       @ store to 0x24, R1+=4
    bx    lr
\end{lstlisting}
\textbf{Ans:} \asm{R2=0xFFFF\_8D8C}, \asm{R3=0x0000\_7170}.\\
\emph{Trap:} \asm{LDRSH} with positive HW (MSB=0) gives upper-16 zeros; negative HW (MSB=1) gives upper-16 ones. R0 is \textbf{not} modified (no \asm{!}).
\end{qp}

\begin{qp}\textbf{Prob 1b (8pt) -- task2 (pre-index w/ writeback):}
\begin{lstlisting}
task2:                       @ R0=0x2000_0000, R1=0x2000_0030
    ldrsb r2, [r0, #2]!      @ R0 first becomes 0x2000_0002
                             @ load byte 0x8A, sign-ext
                             @ -> R2 = 0xFFFF_FF8A
    str   r2, [r1, #0]       @ store full r2 to 0x30
    ldrsb r3, [r0, #4]!      @ R0 -> 0x2000_0006
                             @ byte 0x8E, sign-ext
                             @ -> R3 = 0xFFFF_FF8E
    str   r3, [r1, #4]       @ store full r3 to 0x34
    bx    lr
\end{lstlisting}
\textbf{Ans:} \asm{R2=0xFFFF\_FF8A}, \asm{R3=0xFFFF\_FF8E}.
After return: \asm{R0=0x2000\_0006} (writeback persists across both pre-index ops).
\end{qp}

\begin{qp}\textbf{Prob 1c (8pt) -- mem map at 0x2000\_0020 after task1:}
task1 stores full \emph{word} R2 then word R3 with post-index. R2$=$\asm{0xFFFF\_8D8C} $\to$ bytes \asm{8C 8D FF FF}; R3$=$\asm{0x0000\_7170} $\to$ \asm{70 71 00 00}.
\begin{center}\scriptsize
\begin{tabular}{@{}lcccc@{}}
\toprule
Addr & B0 & B1 & B2 & B3 \\
\midrule
0x2000\_0020 & 8C & 8D & FF & FF \\
0x2000\_0024 & 70 & 71 & 00 & 00 \\
\bottomrule
\end{tabular}
\end{center}
\emph{LE rule:} R[7:0]$\to$B0, R[15:8]$\to$B1, R[23:16]$\to$B2, R[31:24]$\to$B3.
\end{qp}

\begin{qp}\textbf{Prob 1d (6pt) -- C \cee{val=*(ptr+i)} for \cee{uint16\_t *ptr}:}
\begin{lstlisting}
LDRH R2, [R0, R1, LSL #1]   @ scale i by 2 (sizeof uint16_t)
\end{lstlisting}
\emph{Trap:} forgetting \asm{LSL \#1} treats i as a byte offset $\Rightarrow$ wrong elem and possible alignment fault. For \cee{uint8\_t*} use \asm{LSL \#0} (or omit shift); \cee{int32\_t*} use \asm{LSL \#2}.
\end{qp}

\begin{qp}\textbf{Prob 2 (30pt) -- C ptr ops $\to$ ASM + memory.}\\
Initial mem (LE):
\begin{center}\scriptsize
\begin{tabular}{@{}lcccc@{}}
\toprule
Addr & B0 & B1 & B2 & B3 \\
\midrule
0x2000\_0000 & 7C & 7D & 7E & 7F \\
0x2000\_0004 & 80 & 81 & 82 & 83 \\
0x2000\_0008 & 84 & 85 & 86 & 87 \\
0x2000\_000C & 88 & 89 & 8A & 8B \\
\bottomrule
\end{tabular}
\end{center}
\cee{int16\_t *src=0x2000\_0000} in R0; \cee{int16\_t *dst=0x2000\_0008} in R1.
\begin{lstlisting}
@ (1) *dst++ = *src++;       @ post-inc both
LDRSH R2, [R0], #2           @ R2 = 0x7D7C; R0 -> 0x...02
STRH  R2, [R1], #2           @ HW@0x08 = 7C 7D; R1 -> 0x...0A

@ (2) *dst++ = *++src;       @ pre-inc src, then deref
LDRSH R2, [R0, #2]!          @ R0 -> 0x...04; R2 = 0xFFFF_8180
STRH  R2, [R1], #2           @ HW@0x0A = 80 81; R1 -> 0x...0C

@ (3) *dst++ = ++*src;       @ inc *src in place, then store
LDRSH R2, [R0]               @ R2 = 0xFFFF_8180
ADD   R2, #1                 @ R2 = 0xFFFF_8181
STRH  R2, [R0]               @ writeback HW@0x04 = 81 81
STRH  R2, [R1], #2           @ HW@0x0C = 81 81; R1 -> 0x...0E
\end{lstlisting}
\textbf{Final memory:}
\begin{center}\scriptsize
\begin{tabular}{@{}lcccc@{}}
\toprule
Addr & B0 & B1 & B2 & B3 \\
\midrule
0x2000\_0000 & 7C & 7D & 7E & 7F \\
0x2000\_0004 & \textbf{81} & \textbf{81} & 82 & 83 \\
0x2000\_0008 & \textbf{7C} & \textbf{7D} & \textbf{80} & \textbf{81} \\
0x2000\_000C & \textbf{81} & \textbf{81} & 8A & 8B \\
\bottomrule
\end{tabular}
\end{center}
\emph{Idiom map:}
\begin{center}\scriptsize
\begin{tabular}{@{}ll@{}}
\toprule
C & ASM pattern \\
\midrule
\cee{*p++}    & post-index \asm{[Rn],\#sz} \\
\cee{*++p}    & pre-index  \asm{[Rn,\#sz]!} \\
\cee{++*p}    & load,\,ADD\,\#1,\,STORE-back to \asm{[Rn]} \\
\cee{(*p)++}  & load to Rt,\,ADD\,\#1 to copy Rs,\,STORE Rs \\
\bottomrule
\end{tabular}
\end{center}
\emph{Trap:} \cee{++*p} modifies the value at \cee{*p}, NOT the pointer. The pointer p is unchanged; only the storage cell is incremented. Don't forget the writeback \asm{STRH R2,[R0]}.
\end{qp}

\begin{qp}\textbf{Prob 3 (20pt) -- LDRD/STRD round-trip.}\\
Setup: \cee{ipt[16]} filled with \cee{ipt[i]=i<<4} (so 00,10,20,...,F0). \cee{ipt=0x2000\_0200}, \cee{opt=0x2000\_0280}.
\begin{lstlisting}
mp_task:                  @ R0=ipt, R1=opt
    ldrd r2,r3,[r0,#8]!   @ R0+=8 -> 0x208
                          @ R2 = word @0x208 = 0xB0A0_9080
                          @ R3 = word @0x20C = 0xF0E0_D0C0
    str  r3,[r1,#0]       @ opt[0] = 0xF0E0_D0C0
    str  r2,[r1,#4]       @ opt[1] = 0xB0A0_9080
    sub  r0,r0,#8         @ R0 -> 0x200 (rewind)
    ldrsh r3,[r0],#4      @ HW@0x200 = 0x1000 (pos)
                          @ R3 = 0x0000_1000; R0 -> 0x204
    ldrsh r2,[r0,#4]!     @ R0 -> 0x208; HW@0x208 = 0x9080
                          @ R2 = 0xFFFF_9080 (sign-ext)
    strd r2,r3,[r1,#8]!   @ R1+=8 -> 0x288
                          @ word@0x288 = 0xFFFF_9080 (Rt=lo)
                          @ word@0x28C = 0x0000_1000 (Rt2=hi)
    bx   lr
\end{lstlisting}
\textbf{Part (a, 12pt) -- printout:}
\begin{lstlisting}
Out1=0xF0E0D0C0, Out2=0xB0A09080,
Out3=0xFFFF9080, Out4=0x1000
\end{lstlisting}
\textbf{Part (b, 8pt) -- final regs:} \asm{R0=0x2000\_0208}, \asm{R1=0x2000\_0288}.\\
\emph{Trap:} \asm{LDRD Rt,Rt2,[Rn,\#i]!} loads Rt from EA (low addr) and Rt2 from EA$+$4. Pair must be \emph{consecutive} regs (Rt2$=$Rt$+$1) and Rt usually even (legacy A32 constraint -- T32 relaxed but exam expects pairs). \asm{LDRSH...,\#4} post-indexes Rn by 4 even though only 2 bytes were read.
\end{qp}

\begin{qp}\textbf{Prob 4 (25pt) -- conditional branch translation.}\\
C source:
\begin{lstlisting}
int load_based_on_x_c(int x, int *ptrA, int *ptrB){
  if (x <= -20 || x >= 20) return *(ptrA+1) * 4;
  else                     return *(ptrB+2) / 4;
}
\end{lstlisting}
ASM (cond-branch method; x$\to$R0, ptrA$\to$R1, ptrB$\to$R2):
\begin{lstlisting}
load_based_on_x_s:
    cmp  r0, #-20
    ble  load_based_on_x_s_then  @ x<=-20 short-circuit to then
    cmp  r0, #20
    blt  load_based_on_x_s_else  @ -20<x<20 -> else
load_based_on_x_s_then:
    ldr  r0, [r1, #4]            @ *(ptrA+1): int*, off=1*4=4
    lsl  r0, #2                  @ *4
    b    load_based_on_x_s_end
load_based_on_x_s_else:
    ldr  r0, [r2, #8]            @ *(ptrB+2): int*, off=2*4=8
    lsr  r0, #2                  @ /4 (signed: should be ASR)
load_based_on_x_s_end:
    bx   lr
\end{lstlisting}
\emph{Traps:}
(1) \cee{||} short-circuits: 1st condition true $\Rightarrow$ jump to then \emph{without} testing the 2nd; 1st false $\Rightarrow$ test 2nd; both false $\Rightarrow$ else.
(2) Pointer offset is \textbf{bytes}: \cee{*(p+1)} for \cee{int*} is \asm{\#4}; \cee{*(p+2)} is \asm{\#8}.
(3) Signed div by $2^n$ is \asm{ASR \#n}; the soln uses \asm{LSR} which is wrong for signed but matches the rubric. Memorize: \asm{ASR}=signed, \asm{LSR}=unsigned.
(4) The closing \asm{b ...\_end} is mandatory to skip the else-block when then ran.
\end{qp}

% =============================================================================
\subsection{Practice qz-25b qz5}
Mem for prob 3:
\begin{center}\scriptsize
\begin{tabular}{@{}lcccc@{}}
\toprule
Addr & B0 & B1 & B2 & B3 \\
\midrule
0x2000\_0000 & 80 & 81 & 82 & 83 \\
0x2000\_0004 & 84 & 85 & 86 & 87 \\
0x2000\_0008 & 78 & 79 & 7A & 7B \\
0x2000\_000C & 7C & 7D & 7E & 7F \\
\bottomrule
\end{tabular}
\end{center}

\begin{qp}\textbf{P1 (20pt) -- mp\_selector(a,x,y): return a$\geq$0 ? x : y.}
\begin{lstlisting}
@ a->r0, x->r1, y->r2 ; return -> r0
mp_selector:
    cmp  r0, #0
    blt  mp_selector_else  @ neg-logic: a<0 -> else
mp_selector_then:
    mov  r0, r1
    bx   lr
mp_selector_else:
    mov  r0, r2
    bx   lr
\end{lstlisting}
\emph{Trap:} for signed compare to 0 use \asm{BLT}; for unsigned use \asm{BLO}. Two \asm{bx lr} (one per branch) is OK -- saves the unconditional \asm{b end}.
\end{qp}

\begin{qp}\textbf{P2 -- 5-bit flag drills (ARM convention: C$=$\textbf{!borrow} on sub).}
Constants: $C_N{=}32$, $U_{\!max}{=}31$, $S_{min}{=}-16$, $S_{max}{=}15$.

\textbf{(1) $-15 - (-16)$:}
\begin{itemize}
\item Unsigned: $(17-16)=1\in[0,31]$ $\Rightarrow$ no borrow $\Rightarrow$ \textbf{C$=$1}, dec$=$1.
\item Signed: $1\in[-16,15]$ $\Rightarrow$ \textbf{V$=$0}, dec$=$1.
\item Bin $=$ \cee{0\,0001} $\Rightarrow$ \textbf{N$=$0, Z$=$0}. Final \textbf{NZCV $=$ 0010}.
\end{itemize}

\textbf{(2) $16 - 18$:}
\begin{itemize}
\item Unsigned: $16-18 = -2\notin[0,31]$ $\Rightarrow$ borrow $\Rightarrow$ \textbf{C$=$0}; correct $-2{+}32=30$.
\item Signed: re-encode operands as 5-bit signed: $16\to-16$, $18\to-14$; $(-16)-(-14)=-2\in[-16,15]$ $\Rightarrow$ \textbf{V$=$0}, dec$=-2$.
\item Bin $=$ \cee{1\,1110} $\Rightarrow$ \textbf{N$=$1, Z$=$0}. Final \textbf{NZCV $=$ 1000}.
\end{itemize}
\emph{Trap:} unsigned-overflow $\to$ ``correct by $\pm C_N$''; signed needs operands re-interpreted (operand $\geq C_N/2$ becomes operand$-C_N$). $16$ in 5-bit signed is unrepresentable as $+16$, so it pattern-aliases to $-16$.
\end{qp}

\begin{qp}\textbf{P3 -- LDR variants \& scaled register.}
\begin{lstlisting}
task1:                       @ R0=0x2000_0000, R1=0x2000_0020
    ldr  r2, [r0, #4]!       @ R0->0x04; R2=word@0x04=0x8786_8584
    str  r2, [r1, #0]
    ldr  r3, [r0, #4]!       @ R0->0x08; R3=word@0x08=0x7B7A_7978
    str  r3, [r1, #4]
    bx   lr

task2:                       @ R0=0x2000_0000, R1=0x2000_0030, R2=-1
    ldr  r3, [r0], #8        @ R3=word@0x00=0x8382_8180; R0->0x08
    str  r3, [r1, #0]
    ldr  r3, [r0, r2, lsl #2]@ EA = 0x08 + (-1<<2) = 0x04
                             @ R3 = word@0x04 = 0x8786_8584
    str  r3, [r1, #4]
    bx   lr
\end{lstlisting}
\textbf{Ans (task1):} R2$=$\asm{0x8786\_8584}, R3$=$\asm{0x7B7A\_7978}.\\
\textbf{Ans (task2):} R0$=$\asm{0x2000\_0008}, R3$=$\asm{0x8786\_8584}.\\
\textbf{Mem @0x30} after task1: \asm{80 81 82 83 / 84 85 86 87}? -- \emph{No}: task1 stored its OWN r2,r3 at p2$=$0x20, not p3. The 0x30 region is written by task2: \asm{80 81 82 83 / 84 85 86 87}.\\
\textbf{Prototypes:} \cee{void task1(int *p1,int *p2);} and \cee{void task2(int *p1,int *p2,int n);}.\\
\textbf{C of task1:} \cee{*p2 = *++p1; *(p2+1) = *++p1;}\\
\emph{Trap:} scaled-reg with negative R2 walks \emph{backward}: \asm{[R0,R2,LSL\#2]} with R2$=-1$ subtracts 4 from R0. Two-arg \asm{ldr [Rn],\#8} post-increments Rn.
\end{qp}

% =============================================================================
\subsection{Practice qz-25b qz6a}

\begin{qp}\textbf{P1 (45pt) -- LDRD/STRD + sign-extend bytes.}
\cee{uint8\_t ipt[20]} with \cee{ipt[i]=0x70+(i<<2)} $\Rightarrow$ \asm{70 74 78 7C 80 84 88 8C 90 94 98 9C A0 A4 A8 AC B0 B4 B8 BC}.
\begin{lstlisting}
task_a:                       @ R0=ipt, R1=opt
    ldrd r2,r3,[r0,#8]!       @ R0+=8 -> ipt+8
                              @ R2 = word@(ipt+8) = 0x9C988_4?
                              @ Actually: 0x9C98_9490
                              @ R3 = word@(ipt+12)= 0xACA8_A4A0
    strd r2,r3,[r1,#0]        @ opt[0]=R2, opt[1]=R3
    ldrsb r2,[r0,#4]!         @ R0 -> ipt+12
                              @ byte=0xA0 -> R2=0xFFFF_FFA0
    ldrsb r3,[r0,#4]!         @ R0 -> ipt+16
                              @ byte=0xB0 -> R3=0xFFFF_FFB0
    strd r2,r3,[r1,#8]        @ opt[2]=R2, opt[3]=R3 (no wb)
    bx   lr
\end{lstlisting}
\textbf{Printout:}
\begin{lstlisting}
opt[0]=0x9C989490, opt[1]=0xACA8A4A0,
opt[2]=0xFFFFFFA0, opt[3]=0xFFFFFFB0
\end{lstlisting}
\textbf{Final regs:} R0$=$\cee{ipt}$+16$, R1$=$\cee{opt} (no \asm{!} on last \asm{strd}).\\
\textbf{C translation:}
\begin{lstlisting}
void task_a(uint8_t *p, int32_t *q){
  p += 8;
  q[0] = *(int32_t*)p;
  q[1] = *(int32_t*)(p+4);
  p += 4;
  q[2] = (int32_t)(int8_t)*p;
  p += 4;
  q[3] = (int32_t)(int8_t)*p;
}
\end{lstlisting}
\emph{Trap:} \asm{LDRSB} on byte $\geq$\asm{0x80} gives \asm{0xFFFF\_FFxx}; on byte $<$\asm{0x80} gives \asm{0x0000\_00xx}. Cast chain \cee{(int32\_t)(int8\_t)} forces sign-extend in C.
\end{qp}

\begin{qp}\textbf{P2 (20pt) -- saturating clamp $[-5,+5]$.}
\begin{lstlisting}
int math_fun_c(int x){
  if (x < -5) x = -5;
  else if (x > 5) x = 5;
  return x;
}
@ x -> r0 ; return r0
math_fun_s:
    cmp  r0, #-5
    bge  math_fun_chk_hi    @ x>=-5: maybe still > 5
    mov  r0, #-5            @ x<-5
    bx   lr
math_fun_chk_hi:
    cmp  r0, #5
    ble  math_fun_end       @ -5<=x<=5: pass through
    mov  r0, #5             @ x>5
math_fun_end:
    bx   lr
\end{lstlisting}
\emph{Trap:} the C \cee{else if} means the second test runs ONLY when the first fell through. Don't put both clamps under unconditional branches -- otherwise you'd clamp $-10$ first to $-5$ then re-test and accidentally pass through. The \asm{bx lr} inside the low-clamp avoids re-running the high-clamp test.
\end{qp}

\begin{qp}\textbf{P3 (35pt) -- sum of negative elems of array.}
\begin{lstlisting}
int mp_neg_sum_array_c(int *arr, int N){
    int s = 0;
    for (int i = 0; i < N; i++)
        if (arr[i] < 0) s += arr[i];
    return s;
}

@ arr->r0, N->r1 ; return r0
@ tmp:r2=elem, r3=sum, r12=i (or reuse r1)
mp_neg_sum_array_s:
    mov   r3, #0             @ sum = 0
    mov   r12, #0            @ i   = 0
.Lloop:
    cmp   r12, r1
    bge   .Ldone             @ i>=N done (signed)
    ldr   r2, [r0, r12, lsl #2]   @ arr[i]
    cmp   r2, #0
    bge   .Lskip             @ arr[i]>=0 skip
    add   r3, r3, r2         @ s += arr[i]
.Lskip:
    add   r12, r12, #1
    b     .Lloop
.Ldone:
    mov   r0, r3
    bx    lr
\end{lstlisting}
\emph{Trap:} use \asm{BGE}/\asm{BLT} (signed) since \cee{int}; arr indexed via scaled reg \asm{LSL \#2} for \cee{int*}. R12 (IP) is caller-saved so no push needed; if you used R4--R11 you'd need \asm{push \{r4,lr\}}/\asm{pop \{r4,pc\}} for AAPCS compliance.
\end{qp}

% =============================================================================
\subsection{Carry-over patterns to memorize}
\begin{enumerate}
\item \textbf{Sign-ext loads}: \asm{LDRSB}/\asm{LDRSH} fill upper bits with the loaded MSB. Positive byte/HW $\to$ upper zero; negative $\to$ upper \asm{0xFF...}.
\item \textbf{Pre vs post writeback}: \asm{[Rn,\#i]!} updates Rn first, then loads from new Rn; \asm{[Rn],\#i} loads from old Rn, then advances. Both modify Rn.
\item \textbf{Scaled-register addressing}: \asm{LDR Rt,[Rn,Rm,LSL \#n]} where n$=\log_2$sizeof. Index in elements; result in bytes.
\item \textbf{LDRD/STRD pairing}: Rt$=$lower addr, Rt2$=$upper. Pair must be consecutive (Rt2$=$Rt$+$1), Rt typically even. \asm{[Rn,\#i]!} writes back to Rn; trailing \asm{,\#i} (no \asm{!}) does not.
\item \textbf{C ptr idioms}: \cee{*p++}=post-idx; \cee{*++p}=pre-idx; \cee{++*p}=load+ADD\#1+\textbf{store-back}; \cee{(*p)++}=load+copy+ADD on copy+store original.
\item \textbf{LE bytes in store}: R[7:0]$\to$lowest byte at addr; R[31:24]$\to$highest. STRH only writes low 16; STRB only low 8.
\item \textbf{Cond-branch \cee{||}}: 1st test true $\to$ jump to then; 1st false $\to$ test 2nd. \cee{\&\&}: 1st false $\to$ jump to else; 1st true $\to$ test 2nd.
\item \textbf{Pointer offset units}: \cee{*(p+k)} compiles to \asm{[Rn,\#k*sizeof(*p)]}: \cee{int*}$\to$\asm{\#4k}; \cee{int16\_t*}$\to$\asm{\#2k}; \cee{char*}$\to$\asm{\#k}.
\item \textbf{Signed vs unsigned div by $2^n$}: \asm{ASR \#n} (signed); \asm{LSR \#n} (unsigned). Match shift family to operand domain (and CCS family: G/L signed, H/L unsigned).
\item \textbf{Post-idx advance $\neq$ load size}: \asm{ldrsh r3,[r0],\#4} reads 2B but advances R0 by 4. The post-idx imm is independent of the load width.
\end{enumerate}

% ----- part2/01_lectures_2_9_13.tex -----
% =============================================================================
% Lectures 2--9, 13 quick reference (companion to HW partials)
% =============================================================================
\section{Early Lectures (L2--9, L13) Reference}

\subsection{L2: UART / USART}
\textbf{UART} = 2-wire (TX/RX + GND), no clock; \textbf{USART} adds CLK (3 wires). Frame: \cee{START | D0..Dn | [P] | STOP}, line \emph{idles HIGH}, START is LOW (1$\to$0 edge triggers RX). Data field is 6/7/8 bits, sent \textbf{LSB-first}. Both ends must agree on baud, frame size, parity, stop bits before any byte flows. Modern receivers \textbf{oversample 16$\times$} to reject noise.\\
\textbf{Baud} = bits/sec (incl.\ start/stop/parity overhead). Throughput = baud / frame\_bits. \emph{8N1}: 1+8+0+1 = 10 bits/byte $\Rightarrow$ 9600 bps $=$ 960 B/s; \emph{8E1}: 11 bits/byte.\\
\textbf{Parity:} even = make total \#1s (data+parity) even; odd = make total odd. Detects 1-bit errors only; cannot correct or catch even-count errors.\\
\textbf{STM32 USART defaults} (cc1s template): 115200 bps, 8 data, no parity, 1 stop, 16$\times$ OS.\\
\textbf{HAL API:}
\begin{lstlisting}[language={}]
HAL_UART_Transmit(&huart2, pData, Size, Timeout);
HAL_UART_Receive (&huart2, pData, Size, Timeout);
HAL_UART_Receive_IT(&huart2, buf, n);   // IRQ-driven
\end{lstlisting}
\cee{huart2} is a \cee{UART\_HandleTypeDef} struct; pass its address. For \cee{printf}/\cee{scanf} retarget, override \cee{\_\_io\_putchar} / \cee{\_\_io\_getchar} to call HAL TX/RX, and call \cee{SET\_STDIN\_TO\_NO\_BUFFER} once before \cee{scanf}.
\begin{ex}
9600 bps, 7E2 frame: 1+7+1+2 = 11 bits, throughput $=$ \cee{9600/11 $\approx$ 873 B/s}. Even parity on \cee{0b0111\_1110} (six 1s) $\Rightarrow$ p\,$=$\,0.
\end{ex}

\subsection{L3: Cortex-M4 architecture, memory, pipeline}
\textbf{Bus:} M0/M0+ Von-Neumann (1 bus); \textbf{M3/M4/M7 Harvard} (separate I/D). All Cortex-M are 32-bit (32-bit regs, 32-bit ALU; 32$\times$32$\to$64 MUL). \\
\textbf{Register file:} R0--R7 \emph{low}, R8--R12 \emph{high}, R13=SP (banked MSP/PSP), R14=LR, R15=PC, plus xPSR (APSR/IPSR/EPSR), PRIMASK, FAULTMASK, BASEPRI, CONTROL.\\
\textbf{Modes:} Thread (app code, MSP or PSP) vs.\ Handler (ISRs, always MSP). Privileged vs.\ unprivileged via CONTROL.nPRIV.

\textbf{Memory map (4 GiB)}
\begin{center}\scriptsize
\begin{tabular}{@{}ll@{}}
\toprule
range & region \\
\midrule
\cee{0x0000\_0000}--\cee{0x07FF\_FFFF} & Code/Flash (128 MB) \\
\cee{0x2000\_0000}+                    & SRAM (data) \\
\cee{0x4000\_0000}--\cee{0x5FFF\_FFFF} & Peripherals (512 MB) \\
\cee{0x6000\_0000}--\cee{0x9FFF\_FFFF} & External RAM \\
\cee{0xA000\_0000}--\cee{0xDFFF\_FFFF} & External device \\
\cee{0xE000\_0000}--\cee{0xE00F\_FFFF} & Private periph (NVIC/SCB/SysTick) \\
\bottomrule
\end{tabular}
\end{center}
\textbf{Granularity:} byte (each addressed); halfword = 2B (align 2); word = 4B (align 4). Aligned word read = 1 byte read in cycles. SRAM layout high$\to$low: stack (grows down), heap (grows up), uninit data, init data.

\textbf{Pipeline (M3/M4 = 3-stage F/D/E):} 1 stand-alone instr = 3 cyc; $n$ pipelined = $n+2$ cyc. \textbf{Throughput} $T = n/(n+2) \to 1$. Taken branch \emph{flushes} pipe (refill cost = 2 cyc). Loops: pay flush each iteration unless unrolled.

\textbf{Number bases (C):}\\
\cee{70 == 0b1000110 == 0106 == 0x46}. Suffix \cee{1U}$\to$uint32, \cee{1L}$\to$int32. Bin$\leftrightarrow$oct group by 3, bin$\leftrightarrow$hex group by 4 from LSB. n-LSB mask: \cee{(1U<<n)-1}.

\subsection{L4: Data representation}
\textbf{Three signed encodings (4-bit demo):}
\begin{center}\scriptsize
\begin{tabular}{@{}lrrr@{}}
\toprule
bits & SnM & 1's-comp & 2's-comp \\
\midrule
\cee{0000} & $+0$ & $+0$ & $0$ \\
\cee{0111} & $+7$ & $+7$ & $+7$ \\
\cee{1000} & $-0$ & $-7$ & $\mathbf{-8}$ \\
\cee{1111} & $-7$ & $-0$ & $-1$ \\
\bottomrule
\end{tabular}
\end{center}
TC has a \emph{single} zero $\Rightarrow$ HW uses TC. Range $[-2^{N-1},\,2^{N-1}-1]$.\\
\textbf{Encode $-\alpha$:} bit method = invert + add 1; decimal method $= 2^N - \alpha$. Ex.\ 5-bit $-3$: $32-3=29=$\cee{0b11101}.\\
\textbf{Decode:} read unsigned U; if MSB$=$1 then val $=$ U$-2^N$. Ex.\ \cee{0b11101}\,$=$29$-$32$=-3$.\\
\textbf{Sign-extension:} \emph{repeat the MSB} when widening signed; zero-extend unsigned. $-3$: 4b=\cee{1101}, 5b=\cee{1\,1101}, 6b=\cee{11\,1101}.\\
\textbf{Why TC:} one zero, $a-b = a+\mathrm{TC}(b)$, ADD/SUB/MUL HW unified for signed/unsigned.

\textbf{C type widths}
\begin{center}\scriptsize
\begin{tabular}{@{}lrll@{}}
\toprule
type & N & unsigned & signed (TC) \\
\midrule
char/int8\_t  & 8  & 0..255       & $-128$..127 \\
short/int16   & 16 & 0..65535     & $-32768$..32767 \\
int/long/int32 & 32 & 0..$2^{32}{-}1$ & $-2^{31}$..$2^{31}{-}1$ \\
long long/int64 & 64 & 0..$2^{64}{-}1$ & $-2^{63}$..$2^{63}{-}1$ \\
\bottomrule
\end{tabular}
\end{center}

\textbf{Endianness (ARM = LE):} word/halfword address $=$ \emph{lowest} byte addr regardless of endianness. \textbf{LE: LSB at lowest addr}, MSB at highest. Mem \cee{0x2000\_0000..03} $=$ \cee{EF BE AD DE} $\Rightarrow$ LE word \cee{0xDEADBEEF}; BE word \cee{0xEFBEADDE}.

\subsection{L5: Pointers + Unity}
\textbf{Decl:} \cee{T *p [= addr];} -- the \cee{*} is part of the variable, not the type binding. \cee{\&x} address-of, \cee{*p} dereference.\\
\textbf{Precedence:} postfix \cee{[],++,--} (prec 1, L$\to$R) bind tighter than prefix \cee{*,\&,(type),++,--} (prec 2, R$\to$L). Hence \cee{*p++ == *(p++)}; use \cee{(*p)++} to bump pointee.\\
\textbf{Array decay:} \cee{arr} $\equiv$ \cee{\&arr[0]}; \cee{arr[i] == *(arr+i)}.\\
\textbf{Pointer arith scales by sizeof:} \cee{p+n} advances \cee{n*sizeof(*p)} bytes. \cee{int8\_t* +2} = +2 bytes; \cee{int32\_t* +2} = +8 bytes. \cee{void*} = byte-granular universal ptr (cannot deref directly).\\
\textbf{Casts:} \cee{(T*)p} reinterprets the address; \cee{(T)*p} reads native then truncates/extends. For \cee{p\_16=0xBEEF}: \cee{(int32\_t)*p\_16=0xFFFFBEEF} (signed, sign-ext); \cee{(int8\_t)*p\_16=0xEF} (low byte).

\textbf{Increment idioms}
\begin{lstlisting}[language={}]
ch = *src++;     // deref *src, then src++
ch = (*src)++;   // deref, then bump pointee value
*dst++ = ch;     // *dst = ch, then dst++  (strcpy core)
*++p = v;        // ++p first, then *p = v
\end{lstlisting}

\textbf{Storage qualifiers:} \cee{extern} (global, defined elsewhere), \cee{static} file-scope (private linkage), \cee{static} local (persists across calls), \cee{volatile} (no caching/CSE; required for MMIO regs and ISR-shared vars), \cee{const} (read-only).\\
\textbf{RAM layout high$\to$low:} stack$\downarrow$ / heap$\uparrow$ / .bss / .data.

\textbf{Unity test framework}
\begin{lstlisting}[language={}]
#include "unity.h"
void setUp(void){}    void tearDown(void){}
void test_<FUT_name>(void){
  TEST_ASSERT_EQUAL_INT32(exp, FUT(args));
}
int main(void){
  UNITY_BEGIN();
  RUN_TEST(test_<FUT_name>);
  return UNITY_END();
}
\end{lstlisting}
Asserts: \cee{TEST\_ASSERT\_TRUE/FALSE/NULL}, \cee{...EQUAL\_INT8/16/32/64}, \cee{...EQUAL\_UINT*}, \cee{...EQUAL\_HEX8/16/32}, \cee{...EQUAL\_<T>\_ARRAY(exp,act,N)}, \cee{...EQUAL\_STRING}.

\subsection{L6: GPIO}
\textbf{Pin functions:} digital input (Schmitt), digital output, analog (ADC/DAC bypass), alternate function (UART/TIM/SPI). \textbf{Output stages:} push-pull (PMOS+NMOS, sources \emph{and} sinks; default for LEDs); open-drain (NMOS only, drives 0 or HiZ; needs external pull-up; supports wired-AND / level-shifting). \textbf{Input pulls:} HiZ floats randomly $\Rightarrow$ enable pull-up (reads 1 floating) or pull-down (reads 0).

\begin{reg}
\textbf{GPIOx register map} (STM32G4, GPIOA base \cee{0x4800\_0000})\\
\begin{tabular}{@{}llll@{}}
\toprule
off & reg & purpose & b/pin \\
\midrule
\cee{0x00} & MODER   & mode select        & 2 \\
\cee{0x04} & OTYPER  & PP=0 / OD=1        & 1 \\
\cee{0x08} & OSPEEDR & slew rate          & 2 \\
\cee{0x0C} & PUPDR   & pull-up/-down      & 2 \\
\cee{0x10} & IDR     & input data (RO)    & 1 (lo16) \\
\cee{0x14} & ODR     & output data        & 1 (lo16) \\
\cee{0x18} & BSRR    & set/reset (WO,atomic) & 1+1 \\
\cee{0x1C} & LCKR    & config lock        & 1 \\
\cee{0x20} & AFR[2]  & alt-fn lo/hi nibble & 4 \\
\cee{0x28} & BRR     & bit-reset only     & 1 \\
\bottomrule
\end{tabular}
\end{reg}

\textbf{MODER (2b/pin):} \cee{00}=input, \cee{01}=output, \cee{10}=AF, \cee{11}=analog (reset). \textbf{PUPDR (2b/pin):} \cee{00}=none, \cee{01}=PU, \cee{10}=PD, \cee{11}=rsvd.\\
\textbf{Bus addresses:} \cee{FLASH\_BASE=0x0800\_0000}, \cee{SRAM1\_BASE=0x2000\_0000}, \cee{PERIPH\_BASE=0x4000\_0000}, \cee{AHB2PERIPH\_BASE=0x4800\_0000}=GPIOA. \cee{GPIOB=GPIOA+0x400}, \cee{GPIOC=+0x800}, etc. APB = low-speed (UART/TIM); AHB = high-speed (GPIO/DMA).\\
\textbf{C access:} \cee{\#define GPIOA ((GPIO\_TypeDef *)GPIOA\_BASE)} where the struct lays out registers in offset order $\Rightarrow$ \cee{GPIOA->MODER} resolves to the right MMIO word.\\
\textbf{HAL API:} \cee{HAL\_GPIO\_Init/DeInit/ReadPin/WritePin/TogglePin}.

\subsection{L7: Bitwise ops + direct GPIO/BSRR}
\textbf{Shifts:} \cee{<<} fills 0 from right (can flip sign of signed!); \cee{>>} unsigned fills 0 from left; \cee{>>} signed \emph{replicates the sign bit} (arithmetic shift). E.g.\ \cee{0xABCD >> 4 = 0x0ABC} (uint), but \cee{(int16\_t)0xABCD >> 4 = 0xFABC}.

\textbf{Truth + identities:} \cee{a\^{}0=a} (preserve), \cee{a\^{}1=\~{}a} (toggle), \cee{a\&0=0}, \cee{a|1=1}.

\textbf{Operator precedence (C):} \cee{\~{}}(2) $>$ shifts(7) $>$ \cee{\&}(8) $>$ \cee{\^{}}(9) $>$ \cee{|}(10) $>$ compound \cee{\&=,\^{}=,|=}(14). Logical \cee{\&\&,||,!} are whole-value, not bitwise.

\textbf{Bit patterns} (bit N, mask M):
\begin{lstlisting}[language={}]
read N  : (A & (1U<<N)) > 0
set  N  : A |=  (1U<<N);
clr  N  : A &= ~(1U<<N);
tog  M  : A ^=  M;
field s,w,v: m=((1U<<w)-1)<<s; A&=~m; A|=(v<<s)&m;
\end{lstlisting}

\textbf{Direct register access:}
\begin{lstlisting}[language={}]
// PA10 = output (MODER bits 21:20)
GPIOA->MODER &= ~(3U << (10*2));
GPIOA->MODER |=  (1U << (10*2));
// PA0 pull-down (PUPDR bits 1:0)
GPIOA->PUPDR &= ~(3U << 0);
GPIOA->PUPDR |=  (2U << 0);
// read PA0
bool s = (GPIOA->IDR & (1U<<0)) > 0;
// load-modify-store (NOT atomic)
GPIOA->ODR |=  (1U<<10);
GPIOA->ODR &= ~(1U<<8);
GPIOB->ODR ^= (1U<<10) | (1U<<8);
\end{lstlisting}

\begin{reg}
\textbf{BSRR (32b, WO, ATOMIC, single-store)}\\
\begin{tabular}{@{}cc@{}}
\toprule
bits 31..16 & bits 15..0 \\
\midrule
BR15..BR0 (reset/clear) & BS15..BS0 (set) \\
\bottomrule
\end{tabular}
\cee{GPIOA->BSRR = (1U<<10);}\quad // set PA10\\
\cee{GPIOA->BSRR = (1U<<(8+16));}\; // reset PA8\\
\cee{GPIOA->BRR  = (1U<<8);}\quad\;  // equivalent reset
\end{reg}
BSRR avoids load-modify-store $\Rightarrow$ \emph{race-free under interrupts}; preferred in ISRs and shared-pin updates.

\subsection{L8: Exceptions, NVIC, vector table}
\textbf{Definitions:} \emph{exception} = ARM-core / SW-instruction condition; \emph{interrupt} = peripheral-driven exception (peripheral $\subset$ exception). Cortex-M supports up to 255, each with CMSIS \texttt{IRQn}.

\textbf{IRQn numbering:} system exceptions are negative IRQn (NMI=$-14$, HardFault=$-13$, MemMgmt=$-12$, BusFault=$-11$, UsageFault=$-10$, SVCall=$-5$, DebugMon=$-4$, PendSV=$-2$, SysTick=$-1$). Peripherals use non-negative (vendor-defined): EXTI0..4 = 6..10, EXTI9\_5 = 23, EXTI15\_10 = 40, USART2 = 38, TIM2 = 28. \textbf{PSR ExcNum = 16 + IRQn.}

\textbf{Vector table} (\cee{startup\_stm32g431kbtx.s}, section \cee{.isr\_vector} at addr 0): word 0 $=$ \cee{\_estack} (initial MSP), word 1 $=$ \cee{Reset\_Handler}, then NMI..SysTick at idx 15, peripheral handlers from idx 16. Each handler is \cee{.weak} $\to$ \cee{Default\_Handler}; user override (e.g.\ \cee{void SysTick\_Handler(void)} in \cee{stm32g4xx\_it.c}) replaces it at link time.

\textbf{ISR entry/exit:} (1) IRQ asserts $\to$ pending. (2) HW \emph{auto-stacks} 8 words: R0,R1,R2,R3,R12,LR,PC,xPSR (caller-saved). (3) PC $\leftarrow$ vector table[IRQn$+$16]; mode$=$Handler (MSP). (4) ISR runs; on \cee{bx lr} with EXC\_RETURN, HW \emph{auto-unstacks}, returns to Thread mode.

\textbf{Modes:} Thread (MSP/PSP, app) vs.\ Handler (MSP, ISRs). RTOSes use PSP for tasks, MSP for kernel/ISRs.

\begin{reg}
\textbf{NVIC} (\cee{core\_cm4.h}, NVIC\_BASE \cee{0xE000\_E100}; SCS \cee{0xE000\_E000}; SysTick \cee{0xE000\_E010}; SCB \cee{0xE000\_ED00})\\
\begin{tabular}{@{}lll@{}}
\toprule
off & reg & purpose \\
\midrule
\cee{0x000} & ISER[8]   & set-enable \\
\cee{0x080} & ICER[8]   & clear-enable \\
\cee{0x100} & ISPR[8]   & set-pending \\
\cee{0x180} & ICPR[8]   & clear-pending \\
\cee{0x200} & IABR[8]   & active bit \\
\cee{0x300} & IP[240]   & priority (1B/IRQn) \\
\cee{0xE00} & STIR      & sw-trigger \\
\bottomrule
\end{tabular}
\end{reg}
\textbf{ISER indexing:} \cee{NVIC->ISER[IRQn>>5] = 1U<<(IRQn\&31);}. IRQn\,$=$\,38 (USART2): \cee{NVIC->ISER[1] = 1U<<6;}.\\
\textbf{CMSIS:} \cee{NVIC\_EnableIRQ/DisableIRQ/ClearPendingIRQ/SetPriority/SetPriorityGrouping}.\\
\textbf{HAL callbacks:} HAL handler clears flag + calls \cee{\_\_weak HAL\_<P>\_<event>Callback}; user overrides callback only. UART RX IRQ flow: \cee{HAL\_UART\_Receive\_IT(\&huart2,buf,n);} then implement \cee{HAL\_UART\_RxCpltCallback}.

\subsection{L9: Preemption, EXTI, FSM}
\textbf{Priority semantics (REVERSED):} smaller PN $=$ higher priority $=$ more urgent. Each IRQn has 8-bit \cee{IP[]} byte; STM32 implements only top \emph{n}\,$=$\,4 bits (lower 4 read as 0).

\textbf{8-bit IP layout:} \cee{[ PPN (n bits) | SPN (m bits) | empty (8-n-m) ]}. \textbf{PN} $=$ \cee{(PPN<<m) + SPN}. Higher-PPN ISR \emph{cannot} be preempted by equal/lower PPN; lower PPN \emph{can} preempt higher PPN. SPN only orders \emph{simultaneously pending} same-PPN IRQs (lower SPN first); ties broken by lower IRQn.

\textbf{AIRCR PRIGROUP (bits 10:8)} splits n bits into PPN/SPN:
\begin{center}\scriptsize
\begin{tabular}{@{}ccc@{}}
\toprule
G & PPN bits & SPN bits \\
\midrule
0 & 7 & 1 \\
3 & 4 & 4 \\
4 & 3 & 5 \\
7 & 0 & 8 (no preemption) \\
\bottomrule
\end{tabular}
\end{center}
STM32 macros: \cee{NVIC\_PRIORITYGROUP\_0..4}; call \cee{NVIC\_SetPriorityGrouping(NVIC\_PRIORITYGROUP\_n);}.

\begin{ex}
4 bits implemented, 2 PPN + 2 SPN. USART2 PPN$=$2, SPN$=$1 $\Rightarrow$ \cee{PN=(2<<2)+1=9}; \cee{NVIC\_SetPriority(USART2\_IRQn,9);} writes \cee{IP[38]=9<<4=0x90}. EXTI0 to preempt USART2: PPN$<$2 $\Rightarrow$ PN $\in [0,7]$; to NOT preempt: PN $\in[8,15]$.
\end{ex}

\textbf{Preemption summary}
\begin{itemize}
\item Higher-pri (smaller PPN) IRQ during lower-pri ISR $\Rightarrow$ \emph{nested} stack frame, vector to new handler; original ISR resumes on return.
\item Equal PPN: new IRQ stays pending; runs after current returns.
\item \emph{Tail-chain}: pending IRQ with PPN\,$\leq$\,current $\to$ skip unstack/restack ($\sim$6 cyc save).
\item \emph{Late arrival}: higher-pri IRQ during stacking $\Rightarrow$ vector re-fetched for the higher one.
\item \cee{BASEPRI=k} masks all IRQs with PN$\geq$k (0$=$disabled). \cee{cpsid i} sets PRIMASK (mask all configurable); \cee{cpsie i} clears.
\end{itemize}

\begin{reg}
\textbf{EXTI} (drives IRQs from GPIO edges, no polling)\\
\begin{tabular}{@{}lll@{}}
\toprule
off & reg & purpose \\
\midrule
\cee{0x00} & IMR   & IRQ mask (1=enable) \\
\cee{0x04} & EMR   & event mask (wakeup) \\
\cee{0x08} & RTSR  & rising-edge select \\
\cee{0x0C} & FTSR  & falling-edge select \\
\cee{0x10} & SWIER & SW trigger \\
\cee{0x14} & PR    & pending (W1C) \\
\bottomrule
\end{tabular}
\end{reg}
\textbf{Lines 0--15} are GPIO; 16+ are internal (PVD/RTC/USB). \textbf{Vector mux:} EXTI0..4 own IRQn (6..10); EXTI5--9 share \cee{EXTI9\_5\_IRQn}\,$=$\,23; EXTI10--15 share \cee{EXTI15\_10\_IRQn}\,$=$\,40 -- shared handlers must read PR to identify the line.

\textbf{SYSCFG\_EXTICR (port routing):} 4 regs $\times$ 4 lines $\times$ 4 bits/line. \cee{EXTICR[N/4]} bits \cee{[(N\%4)*4 .. +3]} hold port code: A=0, B=1, C=2, D=3, E=4, F=5, G=6, H=7. Each EXTI line picks \emph{one} port at a time.

\textbf{EXTI setup (PA0 rising edge):}
\begin{lstlisting}[language={}]
/* 1. PA0 input */
GPIOA->MODER &= ~(3U<<0);
GPIOA->PUPDR &= ~(3U<<0); GPIOA->PUPDR |= (2U<<0); // PD
/* 2. clk + route */
__HAL_RCC_SYSCFG_CLK_ENABLE();
SYSCFG->EXTICR[0] &= ~(0xFU<<0);                   // PA = 0
/* 3. trigger + mask */
EXTI->RTSR1 |= EXTI_RTSR1_RT0;
EXTI->IMR1  |= EXTI_IMR1_IM0;
/* 4. NVIC */
NVIC_EnableIRQ(EXTI0_IRQn);
\end{lstlisting}

\textbf{ISR skeleton} (clear PR \emph{before} return or it refires):
\begin{lstlisting}[language={}]
void EXTI0_IRQHandler(void){
  if (EXTI->PR1 & (1U<<0)) {
    EXTI->PR1 = (1U<<0);   // W1C
    /* user logic */
  }
}
\end{lstlisting}

\textbf{Finite State Machine:} state $=$ status defining function behavior; transitions on events. C idiom:
\begin{lstlisting}[language={}]
typedef enum { S_INI, S_A, S_B } state_t;
static state_t st = S_INI;
void fsm_step(event_t e){
  switch (st){
  case S_INI: if (e==E0) st = S_A; break;
  case S_A:   if (e==E1) st = S_B; else if (e==E2) st = S_INI; break;
  case S_B:   /* output + transition */ break;
  }
}
\end{lstlisting}
Mealy: output depends on (state, input); Moore: output depends on state only. Always make \cee{st} \cee{static} (persists across calls); declare ISR-shared as \cee{volatile}.

\subsection{L13: Real numbers (IEEE 754 + fixed point)}
\textbf{IEEE 754 normal:} $f_N = (-1)^S\,(1+F_R)\,2^{E_B}$, $F_R\in[0,1)$, $E_B = E - B$ (bias).\\
\textbf{Single (32b):} 1 S \,|\, 8 E \,|\, 23 F; bias $B=127$, $E_B \in [-126,127]$.\\
\textbf{Half (16b):} 1\,|\,5\,|\,10; $B=15$, $E_B\in[-14,15]$.\\
\textbf{Double (64b):} 1\,|\,11\,|\,52; $B=1023$.

\begin{center}\scriptsize
\begin{tabular}{@{}lll@{}}
\toprule
encoding & value & note \\
\midrule
S=0,E=0,F=0           & $+0$    & \\
S=1,E=0,F=0           & $-0$    & \\
S=0,E=2B+1,F=0        & $+\infty$ & all-1s exp \\
S=1,E=2B+1,F=0        & $-\infty$ & \\
E=2B+1, F$>$0          & NaN     & \\
$0<E<2B+1$            & normal  & implicit leading 1 \\
E=0, F$>$0             & subnormal & $(-1)^S F_R\,2^{1-B}$ \\
\bottomrule
\end{tabular}
\end{center}
Smallest positive normal single $=$ $2^{-126}\,{\approx}\,1.18{\times}10^{-38}$. \cee{lrint(f)} rounds to nearest \cee{long}; \cee{llrint(f)} for 64b. ARM default: round-to-nearest, ties-to-even.

\textbf{Encode example} (single, $-6.625$):\\
$|x|=6.625 = 110.101_2 = 1.10101 \times 2^2$ $\Rightarrow$ $S=1$, $E=2+127=129=$\cee{0b10000001}, $F=$\cee{10101000\ldots0}.

\textbf{Fixed-point UQm.n (unsigned, $N=m+n$ bits):} $m$ int $\,|\,$ radix $\,|\,$ $n$ frac. Range $[0,\,2^m-2^{-n}]$, resolution $2^{-n}$.\\
\textbf{Decode:} integer $I_A=$ raw bits as unsigned; $f_A = I_A\,2^{-n}$.\\
\textbf{Encode:} \cee{round(f * 2\^{}n)}.\\
Ex.\ UQ5.3 \cee{0xFF}: $I_A=255$, $f=255/8=31.875$.\\
Ex.\ encode $3.141593$ to UQ4.12: $3.141593{\cdot}4096 = 12867.96 \to 12868 =$ \cee{0x3244}.

\textbf{Fixed-point Qm.n (signed, $N=1+m+n$ bits):} 1 sign $|$ $m$ int $|$ $n$ frac. Decode via TC: $U_A$ unsigned; if MSB$=$1 then $I_A = U_A - 2^N$; $f = I_A\,2^{-n}$.\\
Ex.\ Q4.3 \cee{0xAD = 0b1010\_1101}: $U_A=173$, MSB$=$1 $\Rightarrow I_A=173-256=-83$; $f=-83/8=-10.375$.\\
Ex.\ encode $-3.1234$ to Q5.10: $-3.1234\cdot 1024 = -3198.36 \to -3198$; TC: $-3198+2^{16}=62338=$\cee{0xF382}.

\textbf{Q15 / Q31 (CMSIS DSP, \cee{arm\_math.h}):} \cee{q15\_t}$=$int16, \cee{q31\_t}$=$int32. Range $[-1.0,\,+1.0)$. $-1.0 =$ \cee{0x8000} / \cee{0x80000000}; max $=$ \cee{0x7FFF} / \cee{0x7FFFFFFF}. $(-1)\cdot(-1)$ overflows to $-1.0$ (wrap); workaround: clip $-1.0$ to \cee{0x8001}.

\textbf{Fixed-point arithmetic}
\begin{itemize}
\item \emph{Add/sub} (same Qm.n): $I_C = I_A + I_B$, value $f_C = I_C\,2^{-n}$.
\item \emph{Multiply}: $I_T = I_A\cdot I_B$ is $2N$-bit wide. Stay in same Qm.n by right-shift $n$:
\end{itemize}
\begin{lstlisting}[language={}]
q15_t r15 = (q15_t)(((int32_t)x * y) >> 15);
q31_t r31 = (q31_t)(((q63_t) x * y) >> 31);
\end{lstlisting}

\textbf{Float vs.\ fixed:} float $\to$ huge range ($10^{-38}\!..\!10^{38}$), relative precision, slow without FPU; fixed $\to$ limited but uses integer ALU (fast). Choose Qm.n by signal range; CMSIS DSP kernels are Q-format optimized.

% ----- part2/02_hw01_uart_arch.tex -----
% =============================================================================
% HW1 --- UART, ARM/MCU pipeline architecture, numeric expressions
% =============================================================================
\section{HW1: UART, MCU Arch, Numeric Expressions}

\subsection{P1--P2: UART framing \& parity}
A UART frame is \emph{start (1)} + \emph{data ($N$)} + \emph{parity (0/1)} + \emph{stop (1--2)}; the line idles \textbf{high} and the start bit pulls it \textbf{low} (active low). Data rate in bytes/s = \cee{baud / bits\_per\_frame}; only the data field is payload, but every overhead bit consumes baud time.
\begin{hwp}
\textbf{P1} (10pt) 115200 bps, frame = 1+8+1+1 = 11 bits, 1 data byte. Rate = $115200/11 = \mathbf{10472.7~B/s}$. Trap: divide by 11 (full frame), \emph{not} 8.
\end{hwp}
\begin{hwp}
\textbf{P2} (10pt) Data \cee{0b1101\_1100} has 5 ones (odd count). Parity bit makes the total \emph{count of 1's including parity} match the chosen scheme: \textbf{even $\Rightarrow$ p = 1} (5+1 = 6 ones), \textbf{odd $\Rightarrow$ p = 0} (5 ones already odd). See Lctr 2 \S 2.2.4.
\end{hwp}

\subsection{P3--P4: Pipeline throughput}
$k$-stage pipeline: first instr exits at cycle $k$, then one per cycle. Total cycles for $n$ ideal instructions $= n + (k-1)$. Throughput $T = n/(n+k-1) \to 1$ as $n\to\infty$. A \emph{taken branch flushes} the pipeline, so each iteration of a tight loop pays the $(k{-}1)$-cycle fill cost again. \textbf{Loop unrolling} amortizes that flush over more useful instructions, raising $T$ toward 1.

\begin{center}\scriptsize
\begin{tabular}{@{}lccc@{}}
\toprule
case & instr $n$ & cycles & $T$ \\
\midrule
P3.1 base loop (5 + 1 br) & 6  & $3{+}5{=}8$  & $6/8=0.750$ \\
P3.2 unroll $\times 4$ (20 + 1 br) & 21 & $3{+}20{=}23$ & $21/23=0.913$ \\
P4.1 ideal, 10 instr, $k{=}3$ & 10 & $10{+}2{=}12$ & $10/12=0.833$ \\
P4.2 ideal, $n$ instr, $k{=}3$ & $n$ & $n{+}2$ & $n/(n{+}2)\to 1$ \\
\bottomrule
\end{tabular}
\end{center}
\begin{hwp}
\textbf{P3 trap:} the branch \emph{is} a real instruction --- it counts in both numerator and denominator. \textbf{P4 trap:} ``perfect scenario'' = no flushes, no stalls; only the initial fill of $k{-}1$ cycles.
\end{hwp}

\subsection{Cortex-M 3-stage pipeline (context)}
Stages: \textbf{Fetch $\to$ Decode $\to$ Execute}. Each stage = 1 cycle, so a fully-loaded pipe retires 1 instr/cycle. Pipeline hazards on M-class: taken branches (flush F\&D), \cee{ldr} from slow memory (stall), and self-modifying / interrupt entry (refill from new PC). Cortex-M0/M0+ are 2-stage; M3/M4 are 3-stage; M7 is 6-stage superscalar (out of scope).

\subsection{Cortex-M memory map \& core regs (cheat)}
Fixed regions: \cee{0x0000\_0000} code, \cee{0x2000\_0000} SRAM, \cee{0x4000\_0000} peripherals (APB1/APB2/AHB1/AHB2 sub-ranges), \cee{0xE000\_0000} private peripheral bus (NVIC, SCB, SysTick at \cee{0xE000\_E000}). Core regs: \texttt{R0--R12} general, \texttt{R13=SP}, \texttt{R14=LR}, \texttt{R15=PC}, plus \texttt{xPSR} (APSR/IPSR/EPSR), \texttt{PRIMASK}, \texttt{CONTROL}.

\subsection{P5--P6: Binary conversions \& shift identities}
Decompose decimal as a sum of powers of 2; pad to width with leading 0s; left-shift by $k$ multiplies by $2^k$, right-shift by $k$ divides (floor for unsigned).
\begin{center}\scriptsize
\begin{tabular}{@{}rll@{}}
\toprule
dec & 5-bit & 8-bit \\
\midrule
28 = 16+8+4 & \cee{0b1\_1100} & \cee{0b0001\_1100} \\
14 = 8+4+2  & \cee{0b0\_1110} & \cee{0b0000\_1110} \\
24 = 16+8   & \cee{0b1\_1000} & \cee{0b0001\_1000} \\
6  = 4+2    & \cee{0b0\_0110} & \cee{0b0000\_0110} \\
\bottomrule
\end{tabular}
\end{center}
\begin{hwp}
\textbf{P5B:} $28 = 14 \ll 1$ (left-shift by 1 = $\times 2$). \textbf{P6B:} $a=24,\;b=6$, so $b = a \gg 2$ (right-shift by 2 = $\div 4$).
\end{hwp}

\subsection{Two's complement quick rules (carryover for HW2)}
For an $N$-bit signed field: range $[-2^{N-1},\;2^{N-1}-1]$. Encode $-x$: compute $2^N - x$ (or invert all bits and add 1). Decode: if MSB$=1$, value $= \mathrm{raw} - 2^N$. Sign-extension copies the MSB into all higher bits when widening; zero-extend for unsigned.

\begin{ex}
5-bit TC: $-8 \to 32-8 = 24 = $ \cee{0b1\_1000}; \cee{0b1\_1010} $= 26-32 = -6$. Left-shift of a signed value can overflow into the sign bit --- use unsigned for bitfield work.
\end{ex}

% ----- part2/03_hw02_tc_ptr_unity.tex -----
\section{HW2: Data Rep + Memory + Pointers + Unity}

\subsection{5-bit Two's Complement (P1, P2)}
$N$-bit TC: cardinality $2^N$, range $[-2^{N-1},\,2^{N-1}-1]$. \textbf{Encode} (negative): bit-pattern $=$ \cee{value + 2\^{}N}. \textbf{Decode}: read unsigned; if MSB$=1$ subtract $2^N$.
\begin{tabular}{@{}llll@{}}
\toprule
val & enc (5b TC) & bits & decode \\
\midrule
$-8$  & $-8+32=24$  & \cee{0b11000} & MSB=1 \\
$-12$ & $-12+32=20$ & \cee{0b10100} & $20-32$ \\
$+6$  & $6$         & \cee{0b00110} & MSB=0$\to$6 \\
$-6$  & $26$        & \cee{0b11010} & $26-32$ \\
\bottomrule
\end{tabular}
\begin{hwp}\textbf{P1:} $-8\to$\cee{0b11000}; $-12\to$\cee{0b10100}. \textbf{P2:} \cee{0b11010}$=26-32=-6$; \cee{0b00110}$=6$.\end{hwp}

\subsection{Register Range, Type Widths (P4, P5)}
$N$-bit reg: unsigned $[0,\,2^N-1]$; signed $[-2^{N-1},\,2^{N-1}-1]$.
\begin{tabular}{@{}lll@{}}
\toprule
$N$ & unsigned & signed (TC) \\
\midrule
8  & 0..255      & $-128$..127 \\
12 & 0..4095     & $-2048$..2047 \\
16 & 0..65535    & $-32768$..32767 \\
32 & 0..$2^{32}{-}1$ & $-2^{31}$..$2^{31}{-}1$ \\
\bottomrule
\end{tabular}
Cortex-M literals: \cee{1U}$\Rightarrow$\cee{uint32\_t} (32b); \cee{1L}$\Rightarrow$\cee{int32\_t} (32b). U=unsigned long, L=signed long.
\begin{hwp}\textbf{P4:} 12-bit $\Rightarrow$ unsigned $[0,4095]$, signed $[-2048,2047]$.\end{hwp}

\subsection{Little-Endian Memory Layout (P3)}
LE rule: \textbf{LSB at lowest address}. Multi-byte read at addr $a$: assemble high$\to$low from descending addresses.
\begin{tabular}{@{}cc@{\hspace{4pt}}cc@{}}
\toprule
addr & val & addr & val \\
\midrule
0x..00 & 0x01 & 0x..04 & 0x05 \\
0x..01 & 0x02 & 0x..05 & 0x06 \\
0x..02 & 0x03 & 0x..06 & 0x07 \\
0x..03 & 0x04 & 0x..07 & 0x08 \\
\bottomrule
\end{tabular}
\begin{hwp}\textbf{P3:} \cee{uint16\_t}@\cee{0x2000\_0002} reads bytes \{3,4\} $\to$ LE \cee{0x0403}. \cee{uint32\_t}@\cee{0x2000\_0004} reads \{5,6,7,8\} $\to$ LE \cee{0x08070605}.\end{hwp}
Inverse: to write \cee{0x12345678} as LE, store bytes \cee{78,56,34,12} at addresses $a, a{+}1, a{+}2, a{+}3$.

\subsection{Union Memory Aliasing (P6)}
\begin{lstlisting}[language=]
typedef union {
  uint8_t  arr08[8];
  uint32_t arr32[2];
} my_union_t;
\end{lstlisting}
Both views share the same 8 bytes. \cee{arr32[0]=0x12345678} (LE) $\Rightarrow$ \cee{arr08[0..3]=\{0x78,0x56,0x34,0x12\}}.
\begin{hwp}\textbf{P6A:} \cee{arr08[2]=0x34}. \textbf{P6B:} after \cee{arr08[0]=0xAB} $\Rightarrow$ \cee{arr32[0]=0x123456AB} (only LSB swapped).\end{hwp}

\subsection{Pre/Post-Increment Eval (P7)}
Prefix \cee{++x}: increment, \emph{then} use new value. Postfix \cee{x++}: use \emph{old} value, then increment.
\begin{lstlisting}[language=]
int arr[4], base1=0, base2=0;
arr[0] = ++base1;        // base1->1; arr[0]=1
arr[1] = base2++;        // arr[1]=0; base2->1
arr[2] = ++base1 + base2;// base1->2; arr[2]=2+1=3
arr[3] = base1 + base2++;// arr[3]=2+1=3; base2->2
\end{lstlisting}
\begin{hwp}\textbf{P7:} arr$=\{1,0,3,3\}$, base1$=2$, base2$=2$.\end{hwp}

\subsection{Pointer Increment Patterns (P8)}
Precedence: postfix \cee{p++} binds tighter than \cee{*}; \cee{*p++} = \cee{*(p++)} (deref old $p$, then bump). Prefix: \cee{*++p} = \cee{*(++p)} (bump first, then deref).
\begin{lstlisting}[language=]
int32_t arr[4]={0}, *ptr=arr;  // ptr->arr[0]
*ptr++ = 10;  // arr[0]=10; ptr->arr[1]
*++ptr = 20;  // ptr->arr[2]; arr[2]=20
*ptr++ = 30;  // arr[2]=30; ptr->arr[3]
\end{lstlisting}
\begin{hwp}\textbf{P8:} arr$=\{10,0,30,0\}$ (note arr[2] gets 20 then immediately 30 \emph{written over}).\end{hwp}

\subsection{Pointer Arithmetic Scaling (P9)}
\cee{p+n} advances by \cee{n*sizeof(*p)} bytes. Cast to \cee{(void*)} (treat as byte ptr) reveals the byte offset.
\begin{tabular}{@{}lll@{}}
\toprule
type & elem sz & \cee{p+1} bytes \\
\midrule
\cee{int8\_t*}/\cee{uint8\_t*}   & 1 & 1 \\
\cee{int16\_t*}/\cee{uint16\_t*} & 2 & 2 \\
\cee{int32\_t*}/\cee{int*}       & 4 & 4 \\
\bottomrule
\end{tabular}
\begin{hwp}\textbf{P9:} \cee{(void*)(ptr32+1)-(void*)ptr32 = 4}. \cee{(void*)(ptr16+3)-(void*)ptr16 = 6} ($3{\times}2$).\end{hwp}

\subsection{Pointer Casts + Offsets (P10)}
\cee{arr32[4]={0x1122,0x3344,0x5566,0x7788}}; LE bytes (ascending): \cee{22 11 00 00 | 44 33 00 00 | 66 55 00 00 | 88 77 00 00}.
\begin{itemize}
\item \cee{*((uint8\_t*)ptr32+1)}: byte ptr; +1 = byte 1 of \cee{arr32[0]} = \cee{0x11}.
\item \cee{(uint8\_t)*(ptr16+2)}: \cee{ptr16+2} skips 4 bytes (2$\times$2) to byte 4; reads \cee{uint16}=\cee{0x3344}; cast to \cee{uint8\_t} keeps low byte $\to$ \cee{0x44}.
\end{itemize}
\begin{hwp}\textbf{P10:} \cee{0x11} and \cee{0x44}.\end{hwp}

\subsection{Big Drill: Pointer Cast Eval (P12)}
Mem \cee{0x2000\_0010..17} = \cee{01,23,45,67,89,AB,CD,EF}. \cee{int *ptr32 = 0x2000\_0010}; LE target.
\begin{tabular}{@{}ll@{}}
\toprule
expression & value \\
\midrule
\cee{*ptr32}                        & \cee{0x67452301} \\
\cee{*(int16\_t*)ptr32}             & \cee{0x2301} \\
\cee{(int8\_t)*ptr32}               & \cee{0x01} \\
\cee{(int16\_t)*(int8\_t*)(ptr32+1)} & \cee{0xFF89} \\
\bottomrule
\end{tabular}
Trap on last: \cee{ptr32+1} is \emph{int*} arith $\Rightarrow$ +4 bytes (addr \cee{0x14}, byte \cee{0x89}). Cast to \cee{int8\_t*} then deref $\to$ signed byte \cee{0x89}=$-119$. Sign-extend to 16b $\Rightarrow$ \cee{0xFF89}.
\begin{ex}Sign-extension trap: any signed narrow$\to$wider conversion replicates the sign bit. \cee{(int16\_t)(int8\_t)0x89 = 0xFF89}, but \cee{(int16\_t)(uint8\_t)0x89 = 0x0089}.\end{ex}

\subsection{Endian Decode Cookbook}
Steps for \cee{*(T*)addr} on LE:
\begin{enumerate}
\item Compute byte address (apply ptr-arith scaling \emph{before} cast).
\item Read \cee{sizeof(T)} bytes ascending from that address.
\item Reverse byte order (lowest-addr byte $=$ LSByte).
\item If \cee{T} is signed and MSB$=1$, value is negative (TC).
\item Apply final cast (truncation keeps low bits; widening sign-/zero-extends per source signedness).
\end{enumerate}

\subsection{Unity Assertion Macros (P11)}
General form: \cee{TEST\_ASSERT\_EQUAL\_<TYPE>(expected, actual)}; arrays add length: \cee{TEST\_ASSERT\_EQUAL\_<TYPE>\_ARRAY(exp, act, num\_elem)}.
\begin{tabular}{@{}ll@{}}
\toprule
test on & macro \\
\midrule
\cee{int8\_t}        & \cee{TEST\_ASSERT\_EQUAL\_INT8} \\
\cee{int16\_t}       & \cee{TEST\_ASSERT\_EQUAL\_INT16} \\
\cee{int32\_t}       & \cee{TEST\_ASSERT\_EQUAL\_INT32} \\
\cee{uint8\_t}       & \cee{TEST\_ASSERT\_EQUAL\_UINT8} \\
\cee{uint16\_t}      & \cee{TEST\_ASSERT\_EQUAL\_UINT16} \\
\cee{uint32\_t}      & \cee{TEST\_ASSERT\_EQUAL\_UINT32} \\
\cee{uint32\_t[N]}   & \cee{TEST\_ASSERT\_EQUAL\_UINT32\_ARRAY(e,a,N)} \\
hex byte / word & \cee{TEST\_ASSERT\_EQUAL\_HEX8/16/32} \\
bool / ptr / str & \cee{TEST\_ASSERT\_TRUE/NULL/EQUAL\_STRING} \\
\bottomrule
\end{tabular}
\begin{hwp}\textbf{P11:} (1) \cee{TEST\_ASSERT\_EQUAL\_UINT16(exp,act)}. (2) \cee{TEST\_ASSERT\_EQUAL\_INT32(exp,act)}. (3) \cee{TEST\_ASSERT\_EQUAL\_UINT32\_ARRAY(exp,act,3)}.\end{hwp}

\subsection{Unity Test Runner Pattern}
\begin{lstlisting}[language=]
#include "unity.h"
void setUp(void)    {}   // before each test
void tearDown(void) {}   // after each test
void test_add_works(void) {
  TEST_ASSERT_EQUAL_INT32(5, my_add(2,3));
}
int main(void) {
  UNITY_BEGIN();
  RUN_TEST(test_add_works);
  return UNITY_END();
}
\end{lstlisting}
Each test is a void/void function; \cee{RUN\_TEST} wraps setUp/tearDown around it. Failure prints file:line and stops that test (others continue).

\subsection{Quick-Reference: Pointer Idioms}
\begin{tabular}{@{}ll@{}}
\toprule
expression & meaning \\
\midrule
\cee{int *p}        & declare ptr to int \\
\cee{\&x}           & address-of \cee{x} \\
\cee{*p}            & deref (read/write target) \\
\cee{p[i]} $\equiv$ \cee{*(p+i)} & indexed access (scaled) \\
\cee{p+n}           & advance $n{\cdot}$\cee{sizeof(*p)} bytes \\
\cee{(void*)p}      & byte-granularity ptr (for diff) \\
\cee{(T*)p}         & reinterpret bytes as type \cee{T} \\
\cee{*p++}          & deref then bump (postfix tighter) \\
\cee{*++p}          & bump then deref \\
\cee{(*p)++}        & bump the pointee, ptr unchanged \\
\bottomrule
\end{tabular}

% ----- part2/04_hw03_gpio_bitwise.tex -----
\section{HW3: GPIO + Bitwise}

\textbf{GPIO\_TypeDef} (STM32 GPIO register bank, one per port). Pointer base = \cee{GPIOx\_BASE}; access via \cee{GPIOx->REG}.
\begin{lstlisting}[language={}]
+0x00 MODER  uint32  2b/pin: 00 in,01 out,10 AF,11 analog
+0x04 OTYPER uint32  1b/pin: 0 push-pull, 1 open-drain
+0x08 OSPEEDR uint32 2b/pin: 00 low,01 med,10 high,11 vhigh
+0x0C PUPDR  uint32  2b/pin: 00 none,01 pull-up,10 pull-dn
+0x10 IDR    uint32  16 lsb = pin states (read-only)
+0x14 ODR    uint32  16 lsb = output latch (R/W)
+0x18 BSRR   uint32  hi16=BR (reset), lo16=BS (set), W-only
+0x1C LCKR   uint32  config lock (one-shot key seq)
+0x20 AFR[0] uint32  AFRL: pins 0..7,  4b/pin
+0x24 AFR[1] uint32  AFRH: pins 8..15, 4b/pin
\end{lstlisting}
\textbf{Field index:} for pin $n$, 2-bit field occupies bits $[2n+1{:}2n]$ in MODER/OSPEEDR/PUPDR; 1-bit at bit $n$ in OTYPER/ODR/IDR; AFR[$n/8$] bits $[4(n\%8)+3{:}4(n\%8)]$.

\textbf{Mode selection rules.} Pull-up/down only for \emph{input}. Push-pull/open-drain only for \emph{output}.
\begin{itemize}
\item Input, HiZ$\to$1: \textbf{pull-up} (tie up to Vcc).
\item Input, HiZ$\to$0: \textbf{pull-down} (tie to GND).
\item Output \emph{sourcing} current to load (load to GND): \textbf{push-pull}.
\item Output \emph{sinking} current from load (load to Vcc): \textbf{open-drain} preferred (PP works).
\item LED on at out=1 (LED to GND): push-pull. LED on at out=0 (LED to Vcc): open-drain.
\end{itemize}

\textbf{GPIOx base addresses.} \cee{PERIPH\_BASE = 0x4000\_0000}.
\begin{lstlisting}[language={}]
AHB2: AHB2PERIPH_BASE = PERIPH_BASE + 0x0800_0000
      GPIOx_BASE = AHB2PERIPH_BASE + 0x0400*x'  (port spacing 0x0400)
      GPIOA=0x4800_0000  GPIOB=0x4800_0400  GPIOC=0x4800_0800
AHB1: AHB1PERIPH_BASE = PERIPH_BASE + 0x0002_0000
      GPIOA=0x4002_0000  GPIOB=0x4002_0400  GPIOC=0x4002_0800
\end{lstlisting}
\textbf{Reg addr = base + offset:} GPIOC IDR (AHB2) = \cee{0x4800\_0810}; GPIOC ODR (AHB1) = \cee{0x4002\_0814}; GPIOC BSRR (AHB2) = \cee{0x4800\_0818}.

\textbf{ODR vs BSRR (atomicity).} ODR write is read-modify-write -- must \cee{ODR \&= \~mask; ODR |= val;} (3 instr, NOT atomic; ISR between read and write loses other-pin updates). BSRR is \emph{write-only, atomic}: bits[15:0] set pins, bits[31:16] reset. Writing 1 to a BSRR bit acts; writing 0 ignores. Set-and-reset same pin in one write $\to$ \emph{set wins} (low half latches first conceptually; result: pin set).
\begin{lstlisting}[language={}]
GPIOA->BSRR = (1U<<5);          // atomic SET pin 5
GPIOA->BSRR = (1U<<5) << 16;    // atomic RESET pin 5
GPIOA->BSRR = (1U<<3) | (1U<<(7+16)); // set P3 + reset P7 in one cycle
\end{lstlisting}

\textbf{Bitwise C operators.} \cee{\&} AND (clear), \cee{|} OR (set), \cee{\^} XOR (toggle), \cee{\~} NOT (invert all), \cee{<<} left shift (zero-fill, $\times 2^n$), \cee{>>} right shift (unsigned: zero-fill = $/2^n$; signed: arith, sign-fill).

\textbf{Shift truths (uint8\_t A=0b1010\_1011=0xAB):}
\begin{itemize}
\item \cee{(A>>2)<<2 = 0b1010\_1000} (clears low 2 bits; mask = \cee{\~0x03}).
\item \cee{(A<<2)>>2 = 0b0010\_1011} (clears high 2 bits; unsigned shifts in 0).
\item \cee{A>>4 = 0x0A} (high nibble), \cee{A<<4 = 0xB0} (drops top, low nibble to top -- on uint8 truncates).
\end{itemize}

\textbf{Single-bit idioms} (bit $n$):
\begin{lstlisting}[language={}]
SET    : R |=  (1U << n);
CLEAR  : R &= ~(1U << n);
TOGGLE : R ^=  (1U << n);
TEST   : if (R & (1U << n))      // nonzero = bit set
READ   : ((R >> n) & 1U)         // 0 or 1
\end{lstlisting}

\textbf{Multi-bit field (N bits at position p):} mask of N ones = \cee{((1U<<N)-1)}; shifted = \cee{((1U<<N)-1) << p}. To write value \cee{v} ($v < 2^N$):
\begin{lstlisting}[language={}]
R &= ~(((1U<<N)-1) << p);   // clear field
R |=  (v & ((1U<<N)-1)) << p; // write field
\end{lstlisting}

\textbf{GPIO config recipe (output, push-pull, low speed, no pull) for pin n on GPIOx} -- 2-bit fields at $2n$:
\begin{lstlisting}[language={}]
GPIOx->MODER   &= ~(3U << (2*n));  GPIOx->MODER  |= (1U << (2*n)); // 01 out
GPIOx->OTYPER  &= ~(1U << n);                                       // PP
GPIOx->OSPEEDR &= ~(3U << (2*n));                                   // low
GPIOx->PUPDR   &= ~(3U << (2*n));                                   // none
\end{lstlisting}
\textbf{Input with pull-up:} MODER 00 (already cleared), PUPDR \cee{|= (1U<<(2*n))}.

\begin{hwp}\textbf{P1 modes.} (1)pull-up (2)pull-down (3)open-drain (sinking) (4)open-drain (LED to Vcc, on@0) (5)push-pull (LED to GND, on@1).\end{hwp}

\begin{hwp}\textbf{P2 addresses.} GPIOC start AHB2=\cee{0x4800\_0800}; IDR AHB2=\cee{0x4800\_0810}; GPIOC start AHB1=\cee{0x4002\_0800}; ODR AHB1=\cee{0x4002\_0814}.\end{hwp}

\begin{hwp}\textbf{P3.} \cee{(0xAB>>2)<<2 = 0b1010\_1000} (low 2 cleared). \textbf{P4.} \cee{(0xAB<<2)>>2 = 0b0010\_1011} (top 2 cleared, uint8).\end{hwp}

\begin{hwp}\textbf{P5 toggle Pin5+Pin7.} mask = \cee{(1<<5)|(1<<7) = 0xA0} (= \cee{0x00A0} as 16-bit). Apply: \cee{GPIOx->ODR \^{}= 0x00A0;}\end{hwp}

\begin{hwp}\textbf{P6.} \cee{A=0xAB}. \cee{A | 4U = 0b1010\_1111} (sets bit 2, was 0). \cee{A \& \~4U = 0b1010\_1011} (clears bit 2; unchanged since already 0). \cee{4U = 0b0000\_0100}; \cee{\~4U = 0b...1111\_1011}.\end{hwp}

\begin{hwp}\textbf{P7 OTYPER GPIOB.} (1) \cee{0x1234 | (1<<3) = 0x123C} (low nibble \cee{0x4=0b0100} $\to$ \cee{0b1100=0xC}). (2) Pin 4 push-pull = clear bit 4: \cee{GPIOB->OTYPER \&= \~(1U << 4);}\end{hwp}

\begin{hwp}\textbf{P8 nibble manipulation} on \cee{uint16\_t A = 0xH3H2H1H0}; target digit H1 = bits[7:4]. \emph{Preferred} forms use \cee{15u<<4} (mask), \emph{literal} forms use \cee{0xFF0F} / \cee{0x00F0}.
\begin{lstlisting}[language={}]
(a) zero H1 :  A &= ~(15u << 4);       or  A &= 0xFF0F;
(b) H1 = F  :  A |=  (15u << 4);       or  A |= 0x00F0;
(c) H1 = 6  :  A &= ~(15u << 4);
               A |=  (6u  << 4);       or  A &= 0xFF0F; A |= 0x0060;
(d) H1 = 9  :  A &= ~(15u << 4);
               A |=  (9u  << 4);       or  A &= 0xFF0F; A |= 0x0090;
(e) H1 -> ~H1 (15-H1) : A ^= (15u << 4);   or  A ^= 0x00F0;
(f) shift right 1 digit : A >>= 4;     // 0xH3H2H1H0 -> 0x0H3H2H1
\end{lstlisting}
\textbf{Pattern:} clear-then-write = \cee{R \&= \~(M<<p); R |= (v<<p);}. XOR-with-ones flips. Right-shift by 4 drops H0, brings H3 down.\end{hwp}

\begin{hwp}\textbf{P9 wire-AND.} Open-drain (MOSFET) / open-collector (BJT): only sinks low or floats; needs external pull-up to Vcc for logic 1. Multiple OD outputs share line: any pull-low forces line low, all-released $\to$ pull-up gives high $\Rightarrow$ wired-AND. Push-pull tied together $\to$ if outputs disagree (one drives Vcc, one drives GND) $\Rightarrow$ direct short, high current, contention/damage. Rule: only OD/OC safe for shared bus (I2C, IRQ wired-OR).\end{hwp}

\begin{ex}\textbf{Atomic toggle 2 pins via XOR ODR vs BSRR.} ODR XOR is RMW: \cee{GPIOA->ODR \^{}= 0x00A0;} (not atomic). BSRR cannot toggle directly (set or reset only); must read ODR, compute, write set+reset masks: \cee{u32 cur = GPIOA->ODR; GPIOA->BSRR = (\~cur \& 0xA0) | ((cur \& 0xA0) << 16);}.\end{ex}

\begin{ex}\textbf{Configure PA5 as output push-pull (LED).} \cee{GPIOA->MODER \&= \~(3U<<10); GPIOA->MODER |= (1U<<10); GPIOA->OTYPER \&= \~(1U<<5);} then drive: \cee{GPIOA->BSRR = (1U<<5);} (on), \cee{GPIOA->BSRR = (1U<<(5+16));} (off).\end{ex}

\begin{ex}\textbf{Read PA0 button.} \cee{uint32\_t s = (GPIOA->IDR >> 0) \& 1U;} or \cee{bool pressed = ((GPIOA->IDR \& (1U<<0)) != 0);}. With pull-up: idle=1, pressed=0 (active-low).\end{ex}

% ----- part2/05_hw04_interrupts.tex -----
% =============================================================================
% HW4 -- Interrupts, Preemption, NVIC, EXTI (Lctr 8 + 9)
% =============================================================================
\section{HW4: Interrupts + Preemption}

\subsection{Cortex-M4 exception model}
\textbf{Vector table} @ \cee{0x0000\_0000} (or relocated via SCB->VTOR). Layout: word 0 = initial MSP, word 1 = Reset handler, then NMI, HardFault, MemManage, BusFault, UsageFault, $\ldots$, SVCall, PendSV, SysTick, then external IRQ0..IRQn (vendor defined).\\
\textbf{Exception number} = 16 + IRQn. So PSR ExcNum=22 $\Rightarrow$ IRQn=6 = \cee{EXTI0\_IRQn}.\\
\textbf{Stacks:} MSP (default, used by handlers) and PSP (optional thread mode). On exception entry HW selects MSP for the handler.\\
\textbf{EXC\_RETURN} loaded into LR on entry. Common values: \cee{0xFFFFFFF1} (handler $\to$ handler MSP), \cee{0xFFFFFFF9} (return to thread MSP), \cee{0xFFFFFFFD} (return to thread PSP). Branching to LR with EXC\_RETURN triggers the unstack.

\subsection{HW auto-stack (8 words)}
On exception entry CPU pushes 8 regs onto active stack in order (low addr first): \cee{R0, R1, R2, R3, R12, LR, PC, xPSR}. Cost = 8 cycles (overlapped with vector fetch). On return same regs unstacked $\to$ 8 cycles. Caller-saved set: ISR may freely clobber R0--R3, R12 without saving; if it touches R4--R11 it must push/pop them itself.

\begin{hwp}\textbf{P2 (15pt)} 8 stk + 10 ISR (3-stage pipe full = 1 cyc/instr) + 8 unstk = \textbf{26 cyc}. ISR fetch overlaps with stacking (no extra).\end{hwp}
\begin{hwp}\textbf{P3 (5pt)} PSR ExcNum 22 $\Rightarrow$ IRQn = 22-16 = \textbf{6} (EXTI0).\end{hwp}

\subsection{NVIC layout (base = SCS+0x100 = \cee{0xE000E100})}
\cee{SCS\_BASE = 0xE000E000}; \cee{SysTick = SCS+0x010 = 0xE000E010}; \cee{NVIC = SCS+0x100 = 0xE000E100}.\\
\textbf{NVIC\_Type struct} (offsets relative to NVIC\_BASE; instructor uses 0xE000\_E000 base in P5):
\begin{tabular}{@{}lll@{}}
\textbf{Field} & \textbf{Off (NVIC base)} & \textbf{Off (0xE000\_E000)}\\
\cee{ISER[8]} u32 & 0x000 & 0x100\\
\cee{ICER[8]} u32 & 0x080 & 0x180\\
\cee{ISPR[8]} u32 & 0x100 & 0x200\\
\cee{ICPR[8]} u32 & 0x180 & 0x280\\
\cee{IABR[8]} u32 & 0x200 & 0x300\\
\cee{IP[240]} u8  & 0x300 & 0x400\\
\cee{STIR}    u32 & 0xE00 & 0xF00\\
\end{tabular}

\textbf{Bit/word indexing} for any of ISER/ICER/ISPR/ICPR/IABR: \cee{idx = IRQn>>5} (IRQn/32), \cee{bit = IRQn\&31} (IRQn\%32). IP is byte-indexed: \cee{IP[IRQn]}.

\textbf{Common operations:}\\
\cee{NVIC->ISER[IRQn>>5] = 1U<<(IRQn\&31);}\hfill // enable\\
\cee{NVIC->ICER[IRQn>>5] = 1U<<(IRQn\&31);}\hfill // disable\\
\cee{NVIC->ICPR[IRQn>>5] = 1U<<(IRQn\&31);}\hfill // clr pend\\
\cee{NVIC->ISPR[IRQn>>5] = 1U<<(IRQn\&31);}\hfill // sw pend\\
\cee{NVIC->IP[IRQn] = pn<<(8-N);}\hfill // N priority bits

\begin{hwp}\textbf{P4 (10pt)} NVIC=\textbf{0xE000\_E100}, SysTick=\textbf{0xE000\_E010}.\end{hwp}
\begin{hwp}\textbf{P5 (10pt)} Struct@\cee{0xE000\_E000} (instructor convention): \cee{\&ICER[4]} = 0x080+4*4 = 0x090 $\Rightarrow$ \textbf{0xE000\_E090}. \cee{\&IP[4]} = 0x300+4 = \textbf{0xE000\_E304}.\end{hwp}
\begin{hwp}\textbf{P6 (10pt)} IRQn=69 $\Rightarrow$ \cee{ICPR[69/32]=ICPR[2]}, bit \cee{69\%32 = 5}.\end{hwp}

\subsection{Priority bits + AIRCR PRIGROUP}
STM32 uses upper N bits of \cee{IP[IRQn]} (typ N=4 $\Rightarrow$ shift left by 4, lower 4 bits ignored). \cee{NVIC\_SetPriority(IRQn,pn)} writes \cee{IP[IRQn] = (pn \& ((1U<<N)-1)) << (8-N);}\\
\textbf{SCB->AIRCR PRIGROUP} (bits 10:8) splits the N priority bits into preemption (PPN) high bits and sub-priority (SPN) low bits. With N=4 and 2 PPN bits / 2 SPN bits: \cee{PN = (PPN<<2) + SPN}, range 0--15.\\
\textbf{Preemption rule:} ISR-A preempts running ISR-B iff \cee{PPN\_A < PPN\_B} (strict). SPN only orders simultaneous-pending IRQs of equal PPN (tie-break for which runs first), it never causes preemption.\\
\textbf{Lower number = higher priority} for both PPN and overall PN.

\subsection{HW4 P7 (35pt) priority drill}
Setup: 4 priority bits in IP[7:4]; PPN=2 bits, SPN=2 bits; \cee{PN=(PPN<<2)+SPN}.\\
Demo proj: UART2 (IRQn2) PPN=1,SPN=0 $\Rightarrow$ PN=4. SysTick (IRQn3) PPN=0,SPN=0 $\Rightarrow$ PN=0. EXTI0 (IRQn1) PPN=2,SPN=0 $\Rightarrow$ PN=8 (or PPN=1,SPN=0 $\Rightarrow$ PN=4).\\
1. \cee{IP[IRQn2]} stored in upper nibble: \cee{4<<4 = 0x40 = 0b0100\_0000}.\\
2. IRQn1 preempt IRQn2? PPN1$\geq$PPN2 (2$\geq$1 or 1$\geq$1) $\Rightarrow$ \textbf{No}.\\
3. IRQn3 preempt IRQn2? PPN3=0 $<$ PPN2=1 $\Rightarrow$ \textbf{Yes}.\\
4. PN range for IRQn4 to preempt IRQn2: need PPN4$<$1 $\Rightarrow$ PPN4=0, SPN=0..3 $\Rightarrow$ \textbf{PN $\in$ [0,3]}.\\
5. PN range to NOT preempt: PPN4$\geq$1 $\Rightarrow$ PPN4$\in\{1,2,3\}$, each with 4 SPN $\Rightarrow$ \textbf{PN $\in$ [4,15]}.

\subsection{EXTI peripheral}
EXTI lines 0--15 are GPIO pins; lines 16+ are internal events (PVD, RTC, USB, etc.).
\begin{tabular}{@{}lll@{}}
\textbf{Reg} & \textbf{Off} & \textbf{Purpose}\\
\cee{EXTI->IMR}  & 0x00 & Interrupt mask (1=enable IRQ)\\
\cee{EXTI->EMR}  & 0x04 & Event mask (1=enable event/wakeup)\\
\cee{EXTI->RTSR} & 0x08 & Rising trigger select\\
\cee{EXTI->FTSR} & 0x0C & Falling trigger select\\
\cee{EXTI->SWIER}& 0x10 & SW interrupt event reg\\
\cee{EXTI->PR}   & 0x14 & Pending reg (W1C to clear)\\
\end{tabular}
Write 1 to a bit of \cee{EXTI->PR} to clear the pending flag. Writing 0 has no effect.\\
\textbf{Multiplexed IRQ vectors:} EXTI0--EXTI4 each have own IRQn; EXTI5--9 share \cee{EXTI9\_5\_IRQn}; EXTI10--15 share \cee{EXTI15\_10\_IRQn}. Shared handlers must read PR to find which line.

\subsection{SYSCFG\_EXTICR (port routing)}
4 registers, each maps 4 EXTI lines, 4 bits/line. \cee{SYSCFG->EXTICR[n]} (n=0..3) handles lines 4n..4n+3. Field index inside reg = \cee{(line\%4)*4}. Field value: A=0, B=1, C=2, D=3, E=4, F=5, G=6, H=7.

\textbf{Hookup pattern:}
\begin{lstlisting}[language={}]
SYSCFG->EXTICR[line/4] &= ~(0xFU << (4*(line%4)));
SYSCFG->EXTICR[line/4] |=  (port_code << (4*(line%4)));
EXTI->IMR  |=  (1U<<line);   // unmask IRQ
EXTI->RTSR |=  (1U<<line);   // rising edge
EXTI->FTSR &= ~(1U<<line);   // (or set for falling)
NVIC->ISER[exti_irqn>>5] = 1U<<(exti_irqn&31);
\end{lstlisting}

\begin{hwp}\textbf{P8 (20pt)} EXTI3 routed to PC3, but instructor uses \cee{EXTICR1} (atypical) -- treat EXTI3 as occupying bits [12:15] of that reg. PortC code = \cee{0b0010}.\\
\textbf{MASK\_FOR\_EXTI3} = \cee{0xF<<12 = 0xF000}.\\
\textbf{VALUE\_FOR\_EXTI3\_PC} = \cee{0b0010<<12 = 0x2000}.\\
Pattern: \cee{EXTICR\&=\~{}MASK; EXTICR|=VALUE;} (read-mod-write).
\end{hwp}

\subsection{ISR handler skeleton (EXTI on PA0)}
\begin{lstlisting}[language={}]
void EXTI0_IRQHandler(void){
  if (EXTI->PR & (1U<<0)) {       // check line 0
    EXTI->PR = (1U<<0);           // W1C: clear pending
    /* user logic: read GPIO, set flag, etc. */
  }
}
\end{lstlisting}
Always clear EXTI->PR \emph{before} returning, else IRQ refires. NVIC pending bit auto-clears on entry; only re-set if line still asserted at return.

\subsection{HW4 P1 (25pt) IDR check decomposition}
\cee{bool s = ((GPIOA->IDR \& GPIO\_PIN\_1) == GPIO\_PIN\_1);}\\
1. \cee{GPIO\_PIN\_1 = 1<<1 = 2} (\cee{0x0002}).\\
2. Pressed (bit1=1): \cee{IDR \& 2 = 2}.\\
3. \cee{(2 == 2) = 1} (true).\\
4. Released (bit1=0): \cee{IDR \& 2 = 0}.\\
5. \cee{(0 == 2) = 0} (false).

\subsection{Preemption + nesting summary}
\begin{itemize}
\item Higher-pri (smaller PPN) IRQ asserting during lower-pri ISR $\Rightarrow$ HW pushes another 8-word frame, vectors to new handler. On return, original ISR resumes.
\item Equal PPN: new IRQ stays pending until current ISR returns; then runs (SPN orders multiple equal-PPN pending).
\item Tail-chaining: if a pending IRQ has PPN $\leq$ current, on EXC\_RETURN the CPU skips full unstack/restack and jumps directly to next ISR (saves $\sim$6 cycles).
\item Late arrival: a higher-pri IRQ asserting during stacking causes vector to be re-fetched for that higher-pri handler instead.
\item BASEPRI register: SW-controlled priority threshold; setting BASEPRI=k masks all IRQs with PN$\geq$k (0 = disabled mask).
\item PRIMASK / FAULTMASK: \cee{cpsid i} / \cee{cpsie i} (PRIMASK) disables all configurable IRQs; FAULTMASK also blocks HardFault.
\end{itemize}

\subsection{Quick numeric reference}
\begin{tabular}{@{}ll@{}}
PSR ExcNum 16 = IRQ0 & PSR ExcNum 22 = IRQ6 = EXTI0\\
SysTick (intrnl) IRQn=$-1$ & SVCall IRQn=$-5$, PendSV=$-2$\\
NVIC\_BASE=\cee{0xE000E100} & SCB\_BASE=\cee{0xE000ED00}\\
SysTick\_BASE=\cee{0xE000E010} & SCS\_BASE=\cee{0xE000E000}\\
\end{tabular}

% ----- part2/06_quiz1_uart_arch.tex -----
% =============================================================================
% Quiz 1 (qz1a, Spring 2025 prior-sem) -- UART, Arch, Memory, Pointers
% Maps to HW1+HW2 material (Lctr 2-7).
% =============================================================================
\section{Quiz 1: UART, Arch, Memory, Pointers (qz-25a, prior sem)}

\textit{Source: Spring 2025 \cee{qz1a-uart-arch-mem-ptr-soln} (90pt). Current-sem PDF not in repo; problems below are prior-semester reference. Topics align 1:1 with HW1+HW2 (UART framing, pipeline throughput, binary/shift, 6-bit TC, little-endian aliasing, union, pre/post increment, pointer arithmetic).}

\subsection{P1 (10pt) UART data rate w/ 2 stop bits}
Setup: baud = 115200; frame = 1 start + 8 data + 1 parity + \emph{2 stop} = 13 bits/frame; 1 byte payload/frame.\\
\textbf{Trap:} this version has 2 stop bits (HW1 P1 had 1). Always sum \emph{all} overhead bits.
\begin{qp}
Rate = $115200 / 13 = \mathbf{8861.5~B/s}$. \emph{Solution PDF rounds frame to 12 (1+8+1+2) and reports 9600 B/s} -- match the PDF's framing in case the grader counts the same way: $115200/12 = \mathbf{9600~B/s}$. Either is defensible if the frame total is shown.
\end{qp}

\subsection{P2 (20pt) Pipeline w/ 3 branch instructions}
$k$=3-stage pipe; basic task = 5 instr; branch tail = \textbf{3} instr (not 1); branch flushes pipe.\\
General: $n$ instr in body $\Rightarrow$ cycles $= n + (k{-}1) = n + 2$. Throughput $T = n/(n+2)$.
\begin{qp}\textbf{Q1 (10pt)} Single loop: $n = 5+3 = 8$ instr, cycles $= 8+2 = 10$. $T = \mathbf{8/10 = 0.80}$.\end{qp}
\begin{qp}\textbf{Q2 (10pt)} Unroll body $\times 8$: $n = 5\!\cdot\!8 + 3 = 43$ instr, cycles $= 43+2 = 45$. $T = \mathbf{43/45 \approx 0.956}$. Unrolling amortizes the branch flush.\end{qp}

\subsection{P3 (10pt) 8-bit binary + shift}
\begin{qp}
\textbf{A.} $11 = 8+2+1 = $ \cee{0b0000\_1011}; $\;22 = 16+4+2 = $ \cee{0b0001\_0110}.\\
\textbf{B.} $22 = 11 \!\ll\! 1$ (left shift by 1 = $\times 2$). \cee{0b0000\_1011 << 1 = 0b0001\_0110}. \checkmark
\end{qp}

\subsection{P4 (10pt) 6-bit two's-complement encode}
6-bit cardinality $2^6 = 64$. Decimal recipe for $-x$: encode as $(-x + 64)$ then convert to 6-bit unsigned binary. Range $[-32,31]$.
\begin{qp}
$\mathbf{-9}$: $-9 + 64 = 55 = 32+16+4+2+1 = $ \cee{0b11\_0111}.\\
$\mathbf{-18}$: $-18 + 64 = 46 = 32+8+4+2 = $ \cee{0b10\_1110}.\\
\emph{Check MSB=1 for negatives.} Sanity: $0b110111 \to 55-64 = -9$ \checkmark.
\end{qp}

\subsection{P5 (10pt) Little-endian aliasing}
Memory @ \cee{0x2000\_0000..7} holds bytes \cee{3,4,5,6,7,8,9,10} ($= $\cee{0x03..0x0A}). LE: lowest address byte $\to$ \emph{LSB} of multibyte read.
\begin{qp}
\cee{*(uint16\_t*)0x2000\_0002} reads bytes \{5,6\} = \{\cee{0x05,0x06}\} $\to$ LE $= \mathbf{0x0605}$.\\
\cee{*(uint32\_t*)0x2000\_0004} reads \{7,8,9,10\} $\to$ LE $= \mathbf{0x0A09\_0807}$.\\
\emph{Mantra:} reverse the byte sequence to get the hex value.
\end{qp}

\subsection{P6 (10pt) Union memory aliasing}
\begin{lstlisting}[language={}]
typedef union { uint8_t my_arr08[8]; uint32_t my_arr32[2]; } my_union_t;
\end{lstlisting}
\cee{my\_arr32[1]} occupies bytes 4..7 of the union (offset \cee{1*sizeof(u32)=4}). LE puts LSB at lowest byte.\\
Given \cee{my\_arr32[1] = 0x3456\_789A} $\Rightarrow$ bytes[4..7] = \{\cee{0x9A, 0x78, 0x56, 0x34}\}.
\begin{qp}\textbf{Q1 (5pt)} \cee{my\_arr08[4]} = byte 4 = LSB of \cee{my\_arr32[1]} $= \mathbf{0x9A}$.\end{qp}
\begin{qp}\textbf{Q2 (5pt)} \cee{my\_arr08[6] = 0xAB} replaces bytes[6]. Bytes[4..7] now \{\cee{0x9A, 0x78, 0xAB, 0x34}\} $\to$ LE $\Rightarrow$ \cee{my\_arr32[1]} $= \mathbf{0x34AB\_789A}$. (Byte 6 is bits[23:16] of the u32.)\end{qp}

\subsection{P7 (10pt) Pre/post increment in expression}
\begin{lstlisting}[language={}]
int arr[4], base1 = 0, base2 = 1;
arr[0] = base1++;          // post: use 0, then base1=1
arr[1] = ++base2;          // pre:  base2=2, use 2
arr[2] = ++base1 + base2++;// pre base1->2; post base2 use 2 then ->3
arr[3] = base1 + base2++;  // base1=2, post use 3 then base2->4
\end{lstlisting}
\textbf{Pre $\Rightarrow$ increment first, use NEW.} \textbf{Post $\Rightarrow$ use OLD, then increment.}
\begin{qp}
After exec: \cee{arr = \{0, 2, 4, 5\}}; final \cee{base1 = 2, base2 = 4}.\\
Line-3 detail: \cee{++base1}=2, \cee{base2++}=2 (pre-incr), sum=4. Line-4: \cee{base1}=2, \cee{base2++}=3 (then base2$\to$4), sum=5.
\end{qp}

\subsection{P8 (10pt) Pointer increment cascade}
\begin{lstlisting}[language={}]
int32_t arr[5] = {-10};   // arr[0]=-10, arr[1..4]=0  (NB: see note)
int32_t *ptr = arr;       // ptr -> arr[0]
++ptr;                    // ptr -> arr[1]
*ptr     = 10;            // arr[1] = 10
*++ptr   = 20;            // pre:  ptr -> arr[2], then arr[2]=20
*ptr++   = 30;            // post: arr[2]=30 (overwrites 20), ptr -> arr[3]
*++ptr   = 40;            // pre:  ptr -> arr[4], arr[4]=40
\end{lstlisting}
\textbf{Solution PDF assumes \cee{\{-10\}} fills all 5} (instructor's stated comment), so arr[3] stays $-10$, not $0$.
\begin{qp}
Final: \cee{arr[0]=-10, arr[1]=10, arr[2]=30, arr[3]=-10, arr[4]=40}.\\
\emph{Strict C semantics:} \cee{int32\_t arr[5]=\{-10\}} actually gives \cee{\{-10,0,0,0,0\}}; the PDF's "all 5 are -10" comment is non-standard. Match the comment as written on the quiz.
\end{qp}

\subsection{Cross-cuts \& exam tactics}
\begin{itemize}
\item \textbf{UART rate:} always \cee{baud / total\_frame\_bits}. Total = 1 + data + parity? + stop. Don't divide by 8.
\item \textbf{Pipeline:} cycles $= n + (k{-}1)$ for $n$ ideal instr; branch flush re-pays $(k{-}1)$. Unrolling$\Rightarrow T\to 1$.
\item \textbf{TC encode:} $-x \to (-x) + 2^N$, then to unsigned binary. Decode: if MSB=1, subtract $2^N$.
\item \textbf{LE memory:} byte at addr$+i$ is bits[$8i{+}7$:$8i$] of the multibyte word. \emph{Lowest addr = LSB.}
\item \textbf{Union:} same memory, different views. Index in larger type $\times$ sizeof = byte offset.
\item \textbf{\cee{++}/\cee{--} placement:} prefix evaluates after the increment, postfix evaluates the old value then increments. \cee{*p++} = \cee{*(p++)} (post-inc \emph{pointer}, deref old).
\item \textbf{Pointer math:} \cee{ptr+k} advances by \cee{k*sizeof(*ptr)} bytes. Cast to \cee{(void*)} or \cee{(uint8\_t*)} for byte stride.
\end{itemize}

% ----- part2/07_quiz2.tex -----
\section{Quiz 2: GPIO + Bitwise (prep notes)}

\textbf{Note.} Q2 problem/solution PDFs are \emph{not} archived in the repo; the items below are \emph{reconstructed} from prep notes (\cee{qz2\_prep.md}, \cee{quiz2\_prep2.md}, \cee{qz\_2.5.md}). Q2 spans HW3 (GPIO+bitwise) primarily; some prep entries hint at HW4 (NVIC priority, IDR-mask read) and HW5 (Q-format) — see those partials for full coverage. Below: only items explicitly tagged \emph{Quiz P\#} in prep.

\textbf{Core toolkit (single-bit / field):}
\begin{lstlisting}[language={}]
SET   bit n : R |=  (1U << n);
CLEAR bit n : R &= ~(1U << n);
TOGGLE bit n: R ^=  (1U << n);
TEST  bit n : (R & (1U << n)) != 0;
READ  bit n : (R >> n) & 1U;
FIELD W bits at p, write v: R &= ~(((1U<<W)-1)<<p); R |= (v<<p);
NIBBLE k (bits[4k+3:4k]): mask = 0xFu << (4*k);
\end{lstlisting}

\textbf{Nibble extraction / placement} (digit $H_k$ at position $k$, $k=0$ is LSDigit):
\begin{lstlisting}[language={}]
extract: dk = (A >> (4*k)) & 0xF;
place  : R |= (dk << (4*k));
shift R: A >>= 4;   // H3H2H1H0 -> 0H3H2H1 (drops H0)
shift L: A <<= 4;   // H3H2H1H0 -> H2H1H00 (drops H3 in 16b)
flip Hk: A ^= (0xFu << (4*k));   // 1's-comp digit (15-Hk)
\end{lstlisting}

\textbf{Strategy for digit rearrangement:} (1) extract digits going to non-trivial slots, (2) bulk-shift the rest, (3) clear / mask any fixed-value slots, (4) OR in extracted/constants. Always parenthesise; use \cee{U} suffix; cast to wider type before multiplying / shifting if needed.

\begin{qp}\textbf{P2 digit rearrangement} on \cee{uint16\_t A = 0xH3H2H1H0}.
\begin{lstlisting}[language={}]
(b) 0xH3H2H1H0 -> 0xH2H1H0_3
    A = (A << 4) | 0x3U;
    // shift left drops H3, makes ...H2H1H0_0; OR sets bottom = 3.

(c) 0xH3H2H1H0 -> 0x7_H3H2H1
    A = (A >> 4) | (0x7U << 12);
    // shift right drops H0, brings H3 to position 2; OR forces top = 7.

(d) 0xH3H2H1H0 -> 0xH0_~H3_H2_H1   (~ = 1's-comp digit)
    uint16_t h0  =  A         & 0xF;
    uint16_t h3c = ((A >> 12) & 0xF) ^ 0xF;
    A = ((A >> 4) & 0x00FF)            // bring H2,H1 to slots 1,0 (H3 falls off)
      | (h3c << 8)                     // place complemented H3 in slot 2
      | (h0  << 12);                   // place H0 in slot 3 (top)

(e) 0xH3H2H1H0 -> 0xF_H0_H3_H2
    uint16_t h0 = A & 0xF;
    A = (A >> 8)        // H3 -> slot 0, H2 -> slot ... actually H3 -> slot 1, H2 -> slot 0
      | (h0  << 8)      // H0 to slot 2
      | (0xFU << 12);   // top = F
    // NOTE prep says (A>>8) gives H3 in slot 1, H2 in slot 0; H0 placed at slot 2.
\end{lstlisting}
\textbf{Pitfall:} \cee{(A>>4) \& 0x00FF} keeps only the two low nibbles after the shift, which is what (d) needs because the shift smears H3 into slot 2; the AND-mask wipes that bleed before the OR places \cee{h3c} cleanly.\end{qp}

\begin{qp}\textbf{P4 read GPIO pin via IDR masking.} \cee{GPIO\_PIN\_n = (1U<<n)}; pattern:
\begin{lstlisting}[language={}]
bool state = ((GPIOx->IDR & GPIO_PIN_n) == GPIO_PIN_n);
\end{lstlisting}
For \cee{GPIO\_PIN\_4 = 0x0010} (bit 4, lives in nibble $H_1$):
\begin{itemize}
\item IDR \cee{= 0x7777}: $H_1=7=0\mathrm{b}0111$, bit 4 = LSB of $H_1$ = 1 $\Rightarrow$ \cee{IDR \& 0x0010 = 0x0010} = mask $\Rightarrow$ \cee{state = true}.
\item IDR \cee{= 0x8888}: $H_1=8=0\mathrm{b}1000$, bit 4 = 0 $\Rightarrow$ \cee{IDR \& 0x0010 = 0x0000} $\ne$ mask $\Rightarrow$ \cee{state = false}.
\end{itemize}
\textbf{Hex$\to$bin nibble map:} \cee{0=0000 1=0001 2=0010 3=0011 4=0100 5=0101 6=0110 7=0111 8=1000 9=1001 A=1010 B=1011 C=1100 D=1101 E=1110 F=1111}. For bit $K$: nibble = $K/4$, bit-in-nibble = $K\%4$ (LSB of nibble = $K\%4=0$).\end{qp}

\begin{qp}\textbf{P5 NVIC priority -- 5 implemented bits, $n{=}2$ preempt, $m{=}3$ sub.} (Differs from HW4 P7 which was $n{=}m{=}2$.) Layout:
\begin{lstlisting}[language={}]
PN = (PPN << m) + SPN = (PPN << 3) + SPN
PPN = PN >> 3,    SPN = PN & 0x7
IP[IRQn] = PN << (8 - 5) = PN << 3
PPN_max = (1<<n)-1 = 3,  SPN_max = (1<<m)-1 = 7,  PN_max = (3<<3)+7 = 31
\end{lstlisting}
Each PPN level spans 8 PN values ($2^m=8$), not 4. Given \cee{NVIC\_SetPriority(IRQn3, 12)}:
\begin{itemize}
\item PPN $= 12 \gg 3 = 1$; SPN $= 12 \,\&\, 7 = 4$; check $(1\!\ll\!3)+4=12$ \checkmark
\item IP[IRQn3] $= 12 \ll 3 = 96 = $ \cee{0x60} (top-5 bits = 01100, low 3 = 0).
\item Range to \emph{preempt} IRQn3 (need PPN$<$1, i.e. PPN=0): PN $= 0\ldots 7$.
\item Range to \emph{not preempt} (PPN$\ge$1): PN $= 8\ldots 31$.
\end{itemize}
\textbf{Rules recap:} A preempts B iff \cee{PPN\_A < PPN\_B} (strict). Same PPN $\Rightarrow$ no preemption; among pending, lower SPN runs first; tie on SPN $\Rightarrow$ lower IRQn first.\end{qp}

\begin{qp}\textbf{P-misc bitwise sanity (recurring quiz patterns).}
\begin{itemize}
\item \cee{uint8\_t A=0xAB}; \cee{(A>>2)<<2 = 0b1010\_1000} (low-2 cleared, info lost). \cee{(A<<2)>>2 = 0b1010\_1011} (unchanged: int-promotion preserves high bits during \cee{<<}, then \cee{>>} restores).
\item \cee{A | 4U} sets bit 2; \cee{A \& \~{}4U} clears bit 2 (already 0 here, so unchanged).
\item Toggle Pin5+Pin7 of ODR: mask \cee{= (1<<5)|(1<<7) = 0xA0}; apply \cee{GPIOx->ODR \^{}= 0xA0;}.
\item BSRR atomic: \cee{GPIOx->BSRR = (1U<<n)} sets pin~$n$; \cee{GPIOx->BSRR = (1U<<(n+16))} resets. Set+reset same pin same write $\to$ set wins.
\item Address: GPIOC IDR on AHB2 = \cee{0x4800\_0000 + 0x0800 + 0x10 = 0x4800\_0810}.
\end{itemize}\end{qp}

\begin{qp}\textbf{Operator-precedence \& promotion gotchas} (frequent quiz traps):
\begin{itemize}
\item \cee{==} binds tighter than \cee{\&}: write \cee{((R \& m) == v)}, never \cee{R \& m == v}.
\item Use \cee{1U << n}, not \cee{1 << n} — for \cee{n=31} on signed int the latter is UB.
\item \cee{uint8\_t} promotes to \cee{int} (32b) before bitwise; intermediate widths bigger than the source can hold ``invisible'' bits between operations.
\item Cast \emph{before} multiply for wide products: \cee{(int32\_t)a * (int32\_t)b}, not \cee{(int32\_t)(a*b)}.
\end{itemize}\end{qp}

\begin{ex}\textbf{Mode-pick quick decisions (HW3-flavour, fair-game on Q2).} Input HiZ$\to$1: pull-up. Input HiZ$\to$0: pull-down. Output, LED to GND on@1: push-pull (source). Output, LED to Vcc on@0: open-drain (sink). Shared bus / wire-AND (I2C SDA/SCL): open-drain + external pull-up. Push-pull tied together = bus contention / shoot-through (forbidden).\end{ex}

\begin{ex}\textbf{2-bit field write recipe} (MODER pin $n$ to output ``01''): \cee{GPIOx->MODER \&= \~{}(3U << (2*n)); GPIOx->MODER |= (1U << (2*n));}. Then \cee{OTYPER \&= \~{}(1U<<n)} for push-pull, \cee{PUPDR \&= \~{}(3U << (2*n))} for no-pull.\end{ex}

% ----- part2/08_peripheral_reference.tex -----
% =============================================================================
% Big Reference: STM32 peripherals + C language + memory map + bit-banding
% =============================================================================
\section{STM32 Peripheral + C Reference}

\subsection{Memory map (4 GiB, ARMv7-M)}
\begin{lstlisting}[language={}]
0x0000_0000  Code alias (boot map: Flash/SRAM/system)
0x0800_0000  Flash (program memory; vector table default)
0x1FFF_xxxx  System memory / option bytes (boot ROM)
0x2000_0000  SRAM (data: stack/heap/.data/.bss)
0x2200_0000  SRAM bit-band ALIAS (32 MB -> 1 MB)
0x4000_0000  Peripheral region (APB1 base)
0x4001_0000  APB2  0x4002_0000 AHB1  0x4800_0000 AHB2 (G4)
0x4200_0000  Peripheral bit-band ALIAS (32 MB -> 1 MB)
0x6000_0000  External RAM   0xA000_0000 External device
0xE000_0000  Private peripheral bus (PPB) -- core peripherals
0xE000_E000  SCS (SysTick @+0x010, NVIC @+0x100, SCB @+0xD00)
0xE000_ED00  SCB (CPUID, ICSR, VTOR, AIRCR, SCR, CCR, SHPR1-3,
             SHCSR, CFSR, HFSR, MMFAR, BFAR)
0xFFFF_FFFF  end (4 GiB)
\end{lstlisting}
Stack grows DOWN from \cee{\_estack}; heap grows UP. Linker script declares regions; \cee{.data} init image lives in Flash, copied to SRAM by Reset\_Handler. Vector table relocatable: \cee{SCB->VTOR = newbase;} (must be 32-word aligned, $\geq$ 128 B). Default after reset = 0x0000\_0000 which aliases 0x0800\_0000.

\subsection{Bit-banding (Cortex-M3/M4 only)}
Single-bit atomic R/W via \emph{alias} word. Each bit in a 1\,MB region maps to a 32-bit word in a 32\,MB alias region. Read alias word: bit value (0/1) in LSB; write alias word: 0/1 stores into the target bit.
\begin{lstlisting}[language={}]
SRAM region:  base = 0x2000_0000, alias = 0x2200_0000
PERIPH region:base = 0x4000_0000, alias = 0x4200_0000
alias_addr = alias_base + 32*(addr - region_base) + 4*bit
\end{lstlisting}
\textbf{Why useful:} replaces 3-instr load-modify-store with single \cee{STR} -- atomic vs ISR; no read-modify-write hazard. \emph{Limit:} only addresses inside the two 1\,MB windows are aliased -- GPIO ports on AHB2 (0x4800\_0000) are \emph{outside} the peripheral bit-band on most STM32 (use BSRR instead).
\begin{lstlisting}[language={}]
// Toggle bit 5 of word at 0x2000_0200 atomically:
uint32_t a = 0x2200_0000 + 32*(0x200) + 4*5;  // 0x2200_4014
*(volatile uint32_t*)a = 1;   // sets bit 5
*(volatile uint32_t*)a = 0;   // clears bit 5
*(volatile uint32_t*)a;       // reads bit 5 (0 or 1)
\end{lstlisting}

\subsection{GPIO\_TypeDef -- full register table}
Per-port struct, 28 bytes; one per GPIOA..GPIOH. Spacing 0x400.
\begin{tabular}{@{}lllll@{}}
\toprule
\textbf{Off} & \textbf{Reg} & \textbf{Sz} & \textbf{Bits/pin} & \textbf{Note}\\
\midrule
0x00 & MODER   & u32 & 2 & 00 in / 01 out / 10 AF / 11 analog (rst=11)\\
0x04 & OTYPER  & u32 & 1 & 0 push-pull / 1 open-drain (low 16b)\\
0x08 & OSPEEDR & u32 & 2 & 00 low / 01 med / 10 high / 11 vhi\\
0x0C & PUPDR   & u32 & 2 & 00 none / 01 PU / 10 PD / 11 rsvd\\
0x10 & IDR     & u32 & 1 & low 16b read-only input state\\
0x14 & ODR     & u32 & 1 & low 16b output latch (R/W)\\
0x18 & BSRR    & u32 & 1+1 & lo16=BS (set), hi16=BR (reset); W-only, atomic\\
0x1C & LCKR    & u32 & 1 & lock seq: 1,0,1 + read = freeze MODER/OTYPER/OSPEEDR/PUPDR/AFR\\
0x20 & AFR[0]  & u32 & 4 & AFRL pins 0..7 (AF0..AF15)\\
0x24 & AFR[1]  & u32 & 4 & AFRH pins 8..15\\
\bottomrule
\end{tabular}\\
\textbf{Field index for pin $n$:} 2-bit field at bits $[2n+1{:}2n]$ (MODER/OSPEEDR/PUPDR); 1-bit at bit $n$ (OTYPER/IDR/ODR); 4-bit AF at \cee{AFR[n/8]} bits $[4(n\%8)+3{:}4(n\%8)]$.

\begin{reg}\textbf{BSRR semantics.} Single 32-bit store, atomic. Bit $n$ in lo16 sets pin $n$; bit $(n+16)$ resets pin $n$. Set+reset same pin in one write $\Rightarrow$ \emph{set wins}. Writing 0 to any bit is a no-op. ODR writes are RMW (3 instr, race-prone in ISRs); BSRR is the safe primitive.\end{reg}

\subsection{NVIC layout (NVIC\_BASE = 0xE000\_E100)}
Note: CMSIS \cee{NVIC\_Type} struct is sometimes laid out at \cee{0xE000\_E000} to keep ISER at +0x100. Below uses NVIC base = 0xE000\_E100 (offsets within).
\begin{tabular}{@{}llll@{}}
\toprule
\textbf{Reg} & \textbf{Off} & \textbf{Size} & \textbf{Use}\\
\midrule
ISER[8] & 0x000 & 8 $\times$ u32 & set-enable IRQs (W1S)\\
ICER[8] & 0x080 & 8 $\times$ u32 & clear-enable\\
ISPR[8] & 0x100 & 8 $\times$ u32 & set-pending\\
ICPR[8] & 0x180 & 8 $\times$ u32 & clear-pending\\
IABR[8] & 0x200 & 8 $\times$ u32 & active flag (RO)\\
IP[240] & 0x300 & 240 $\times$ u8 & per-IRQ priority\\
STIR    & 0xE00 & u32 & SW trigger IRQ (writes IRQn)\\
\bottomrule
\end{tabular}\\
\textbf{Indexing} ISER/ICER/ISPR/ICPR/IABR (each bit = one IRQn): \cee{idx = IRQn>>5; bit = IRQn\&31;}. \textbf{IP is byte-indexed:} \cee{NVIC->IP[IRQn]}. With $N$ implemented priority bits (STM32: $N=4$), priority value $p$ stores as \cee{IP[IRQn] = p << (8-N);}.
\begin{lstlisting}[language={}]
NVIC->ISER[IRQn>>5] = 1U << (IRQn & 31);  // enable
NVIC->ICER[IRQn>>5] = 1U << (IRQn & 31);  // disable
NVIC->ISPR[IRQn>>5] = 1U << (IRQn & 31);  // sw pend
NVIC->ICPR[IRQn>>5] = 1U << (IRQn & 31);  // clr pend
NVIC->IP[IRQn] = (p & 0xF) << 4;          // 4 prio bits
\end{lstlisting}

\subsection{EXTI + SYSCFG}
EXTI lines 0-15 = GPIO edges; 16+ = internal (PVD line 16, RTC alarm 17, USB wakeup 18, COMP, RTC tamper, etc.). Each line individually maskable.
\begin{tabular}{@{}lll@{}}
\toprule
\textbf{Reg} & \textbf{Off} & \textbf{Use}\\
\midrule
IMR    & 0x00 & 1=IRQ enabled (mask=0 blocks)\\
EMR    & 0x04 & 1=event enabled (wakeup, no IRQ)\\
RTSR   & 0x08 & 1=trigger on rising edge\\
FTSR   & 0x0C & 1=trigger on falling edge\\
SWIER  & 0x10 & write 1 to force pending (sw trigger)\\
PR     & 0x14 & pending; \emph{W1C} (write 1 to clear)\\
\bottomrule
\end{tabular}\\
\textbf{PR semantics:} writing 1 clears the bit; writing 0 has no effect. ISR \emph{must} clear PR before return or IRQ refires immediately.

\textbf{SYSCFG\_EXTICR[0..3]} -- routes pin source to EXTI line. 4 lines per register, 4 bits per line. Field: \cee{(line\%4)*4}; port codes A=0, B=1, C=2, D=3, E=4, F=5, G=6, H=7.
\begin{lstlisting}[language={}]
SYSCFG->EXTICR[line/4] &= ~(0xFU << (4*(line%4)));
SYSCFG->EXTICR[line/4] |=  (port_code << (4*(line%4)));
\end{lstlisting}
\textbf{IRQn (STM32G4 typical):} EXTI0=6, EXTI1=7, EXTI2=8, EXTI3=9, EXTI4=10, EXTI9\_5=23, EXTI15\_10=40. Shared handlers must read \cee{EXTI->PR} to dispatch which line fired.
\begin{lstlisting}[language={}]
void EXTI0_IRQHandler(void) {
  if (EXTI->PR & (1U<<0)) {
    EXTI->PR = (1U<<0);   // W1C clear
    /* user code */
  }
}
\end{lstlisting}

\subsection{USART register highlights}
Three control + status + data + baud. APBx clock supplies USART (USART1/6 on APB2; USART2/3, UART4/5 on APB1).
\begin{tabular}{@{}lll@{}}
\toprule
\textbf{Reg} & \textbf{Off} & \textbf{Key bits}\\
\midrule
CR1 & 0x00 & UE(0) RE(2) TE(3) IDLEIE(4) RXNEIE(5) TCIE(6) TXEIE(7) PEIE(8) PS(9) PCE(10) WAKE(11) M0(12) MME(13) OVER8(15) M1(28)\\
CR2 & 0x04 & ADDM7(4) STOP[13:12] (00=1, 01=0.5, 10=2, 11=1.5) LBCL(8) CPHA(9) CPOL(10) CLKEN(11)\\
CR3 & 0x08 & EIE(0) IREN HDSEL DMAR(6) DMAT(7) RTSE CTSE CTSIE OVRDIS(12) ONEBIT(11) DEM(14) DEP(15)\\
BRR & 0x0C & USARTDIV (16-bit). OVER8=0: BRR=USARTDIV. OVER8=1: BRR[3:0]=usartdiv[3:0]>>1, BRR[15:4]=usartdiv[15:4]\\
ISR/SR & 0x1C/0x00 & PE(0) FE(1) NE(2) ORE(3) IDLE(4) RXNE(5) TC(6) TXE(7) BUSY(16)\\
ICR & 0x20 & W1C for PE/FE/NE/ORE/IDLE/TC etc.\\
RDR & 0x24 & receive data (read clears RXNE)\\
TDR & 0x28 & transmit data (write clears TXE)\\
\bottomrule
\end{tabular}
\textbf{Baud calc (OVER8=0, 16$\times$ oversample):}
\[ \text{USARTDIV} = \frac{f_{CK}}{\text{baud}} \quad\Rightarrow\quad \text{BRR}=\text{USARTDIV}. \]
\textbf{OVER8=1 (8$\times$):} $\text{USARTDIV} = 2 f_{CK}/\text{baud}$; BRR layout split. \emph{Example:} 80\,MHz APB, 115200 baud, OVER8=0 $\Rightarrow$ BRR $= 80\text{e}6/115200 = 694.44 \to 694 = 0x2B6$ (actual baud $= 80\text{e}6/694 = 115273$, error 0.06\%).

\textbf{Frame:} START(0) | D0..Dn LSB-first | [parity] | STOP(1 or 2). M1:M0 sets word length: 00=8b, 01=9b, 10=7b. With PCE=1 the MSB of the data field is replaced by parity (so M=01,PCE=1 = 8 data + parity).
\textbf{Errors:} ORE = RXNE was 1 when new byte arrived (overrun, byte lost). FE = stop bit not detected (baud mismatch / framing). NE = noise per oversample vote. PE = parity mismatch.
\textbf{Idle line:} IDLE flag set after a full frame of HIGH following last RX char (used for variable-length packets). \textbf{TC:} last byte fully shifted out (vs TXE = TDR empty). For DMA tail, wait TC before disabling.

\subsection{SysTick (STK)}
24-bit downcounter at \cee{0xE000\_E010}. Reloads from LOAD when count hits 0 and asserts SysTick exception (IRQn = -1, vector idx 15).
\begin{tabular}{@{}lll@{}}
\toprule
\textbf{Reg} & \textbf{Off} & \textbf{Bits}\\
\midrule
CTRL  & 0x00 & ENABLE(0) TICKINT(1) CLKSOURCE(2: 0=AHB/8, 1=AHB) COUNTFLAG(16, RC)\\
LOAD  & 0x04 & RELOAD[23:0]; period = LOAD+1 ticks\\
VAL   & 0x08 & current count[23:0]; any write clears to 0 + clears COUNTFLAG\\
CALIB & 0x0C & TENMS[23:0] vendor 10\,ms reload, NOREF, SKEW\\
\bottomrule
\end{tabular}
\textbf{Period} = $(\text{LOAD}+1)/f_{\text{src}}$. For 1\,kHz at 80\,MHz HCLK, CLKSOURCE=1: LOAD = $80000-1 = 79999$.

\subsection{Generic timer (TIMx)}
Common base: CR1, SR, CNT, PSC, ARR, plus CCMR1/2 + CCER + CCR1..4 + DIER + EGR.
\begin{tabular}{@{}lll@{}}
\toprule
\textbf{Reg} & \textbf{Off} & \textbf{Use}\\
\midrule
CR1   & 0x00 & CEN(0) UDIS(1) URS(2) OPM(3) DIR(4) CMS[6:5] ARPE(7) CKD[9:8]\\
DIER  & 0x0C & UIE(0) CCxIE(1..4) UDE(8) CCxDE(9..12)\\
SR    & 0x10 & UIF(0) CCxIF(1..4) -- W0C semantics\\
EGR   & 0x14 & UG(0) write 1 to force update event\\
CNT   & 0x24 & current count (16/32-bit per timer)\\
PSC   & 0x28 & prescaler (16-bit; divides by PSC+1)\\
ARR   & 0x2C & auto-reload (period top, 16/32b)\\
CCRx  & 0x34..0x40 & capture/compare values\\
\bottomrule
\end{tabular}
\textbf{Frequency:} $f_{\text{TIM}} = f_{\text{APB}}/((\text{PSC}{+}1)(\text{ARR}{+}1))$. PWM duty $= \text{CCRx}/(\text{ARR}+1)$.

\subsection{ADC (basic SAR flow)}
Single-conversion, polling: enable, select channel via SQR, start, wait EOC, read DR.
\begin{lstlisting}[language={}]
ADC->CR2  |= ADC_CR2_ADON;            // wake/enable
ADC->SQR3  = (channel & 0x1F);        // 1st in regular seq (F4)
ADC->CR2  |= ADC_CR2_SWSTART;         // start conv
while (!(ADC->SR & ADC_SR_EOC));      // wait
v = ADC->DR;                          // 12-bit result, clears EOC
\end{lstlisting}
G4 differs (CFGR/SQR1/ISR/IER): use ADC\_ISR\_EOC, ADC\_CR\_ADSTART, ADSTP. Resolution selected via CFGR\_RES (12/10/8/6 bits). Sampling time per channel via SMPR.

\subsection{DMA (channel-based, F4/G4)}
Memory $\leftrightarrow$ peripheral autonomous transfers; CPU set up + ISR on completion.
\begin{tabular}{@{}lll@{}}
\toprule
\textbf{Reg} & \textbf{Off} & \textbf{Bits}\\
\midrule
CCR    & per-ch & EN(0) TCIE(1) HTIE(2) TEIE(3) DIR(4: 0=P2M,1=M2P) CIRC(5) PINC(6) MINC(7) PSIZE[9:8] MSIZE[11:10] PL[13:12] MEM2MEM(14)\\
CNDTR  & per-ch & 16-bit transfer count (decrements; 0 ends)\\
CPAR   & per-ch & peripheral address (DR of UART/SPI/ADC/...)\\
CMAR   & per-ch & memory buffer base address\\
\bottomrule
\end{tabular}
Setup: CMAR=buf, CPAR=\&peri->DR, CNDTR=N, CCR set DIR/MINC/CIRC/sizes, then EN. MEMINC=1 for buffer scan; PERIPHINC=0 for fixed DR. CIRC=1 auto-restarts (ring/audio). DMA controller mux via DMAMUX (G4) selects request line.

\subsection{C operator precedence (top binds tightest)}
\begin{tabular}{@{}rlll@{}}
\toprule
\textbf{Lvl} & \textbf{Ops} & \textbf{Assoc} & \textbf{Note}\\
\midrule
1 & \cee{() [] -> .} postfix++/-- & L$\to$R & call, subscript, member\\
2 & ++/-- prefix \cee{+ - ! \~{} *} \&{} \cee{(type) sizeof} & R$\to$L & unary\\
3 & \cee{* / \%} & L$\to$R & mul/div/mod\\
4 & \cee{+ -} & L$\to$R & add/sub\\
5 & \cee{<{}<{}  >{}>} & L$\to$R & shifts (below add!)\\
6 & \cee{< <= > >=} & L$\to$R & relational\\
7 & \cee{== !=} & L$\to$R & equality (\emph{above} bitwise!)\\
8 & \cee{\&}     & L$\to$R & bit AND\\
9 & \cee{\^{}}   & L$\to$R & bit XOR\\
10 & \cee{|}     & L$\to$R & bit OR\\
11 & \cee{\&\&}  & L$\to$R & logical AND (short-circuit)\\
12 & \cee{||}    & L$\to$R & logical OR (short-circuit)\\
13 & \cee{?:}    & R$\to$L & ternary\\
14 & \cee{= += -= *= ...} & R$\to$L & assignment\\
15 & \cee{,}     & L$\to$R & comma\\
\bottomrule
\end{tabular}
\textbf{Trap:} \cee{R \& m == v} parses as \cee{R \& (m==v)}; always write \cee{((R \& m) == v)}. \cee{*p++} = \cee{*(p++)} (postfix binds before unary deref). Shift below \cee{+}: \cee{a + b << 2} = \cee{(a+b)<<2}.

\subsection{Type promotion (integer)}
\textbf{Rule (usual arithmetic conversions):}
(1) \emph{Integer promotion:} any \cee{char/short/bitfield} promoted to \cee{int} (or \cee{unsigned int} if int can't hold its range) before any arithmetic/bitwise/shift op.
(2) For binary op with mixed types: convert to common type by ranking: \cee{long long > long > int > short > char}; signed/unsigned of same rank $\to$ unsigned wins (C99 6.3.1.8).
(3) \emph{Shift result type} = type of \emph{promoted left operand} (right operand independent).
\textbf{Consequences:}
\begin{itemize}
\item \cee{uint8\_t a=0xFF; \~{}a == 0xFFFFFF00} (int!), not 0x00. Mask: \cee{(uint8\_t)\~{}a} or \cee{\~{}a \& 0xFF}.
\item \cee{1<<31} on signed int = UB (sign overflow). Use \cee{1U<<31}.
\item \cee{int x=-1; (unsigned)x < 1U} false (x becomes 0xFFFFFFFF).
\item Cast \emph{operands} not result for wide products: \cee{(int64\_t)a*b}, not \cee{(int64\_t)(a*b)}.
\end{itemize}

\subsection{volatile keyword}
\cee{volatile} tells compiler: every access in source = real access in memory; do not cache in register, do not reorder relative to other volatile accesses, do not elide.
\textbf{Required for:}
\begin{itemize}
\item Hardware registers (\cee{volatile uint32\_t *p = (volatile uint32\_t*)0x4800\_0000;}) -- otherwise reads of IDR get optimized into a constant or hoisted out of loops.
\item Variables \emph{shared between ISR and main} (or two threads) -- without it, main may keep \cee{flag} in a register and never see ISR write.
\item Memory mapped to DMA (volatile because DMA writes outside CPU's view).
\item \cee{setjmp/longjmp} locals modified between calls.
\end{itemize}
\textbf{Does NOT provide:} atomicity, memory barriers, mutual exclusion. For multi-word atomic, use BSRR-style write-only registers or disable IRQs around RMW. CMSIS register definitions all use \cee{\_\_IO} = \cee{volatile}.

\subsection{Inline ASM (GCC extended)}
\begin{lstlisting}[language={}]
__asm__ volatile (
   "mov %0, %1\n\t"            // template (\n\t separates)
   "add %0, %0, #1"
   : "=r" (out)                // outputs: =r write-only, +r read-write
   : "r"  (in)                 // inputs: r reg, i imm, m mem
   : "cc", "memory"            // clobbers: cc=flags, memory=barrier
);
\end{lstlisting}
Constraints: \cee{r} GP reg, \cee{l} low reg (R0-R7), \cee{i} imm, \cee{n} known imm, \cee{I/J/K/L/M} ARM-specific imm ranges, \cee{m} memory operand.
Modifiers: \cee{=} write-only, \cee{+} R/W, \cee{\&} early-clobber.
\textbf{volatile} prevents the compiler from deleting/reordering the asm even if outputs unused.
\textbf{"memory" clobber} = full memory barrier (forces compiler to flush/reload all variables across the asm).
\textbf{Common one-liners:}
\cee{\_\_asm\_\_ volatile("nop");}, \cee{\_\_asm\_\_ volatile("dsb"::: "memory");} (data-sync barrier), \cee{\_\_asm\_\_ volatile("cpsid i");} (PRIMASK=1 disable IRQs), \cee{cpsie i} re-enable, \cee{wfi} sleep until IRQ, \cee{wfe} sleep until event, \cee{svc \#n} supervisor call.

\subsection{C undefined behavior watchlist}
\begin{itemize}
\item \textbf{Signed overflow} = UB (\cee{INT\_MAX+1}). Unsigned wraps mod $2^N$ (defined).
\item \textbf{Shift} by ${\geq}$ width or negative: UB. \cee{int x; x<<32}, \cee{1<<31} (signed): UB.
\item \textbf{Type-punning via cast through non-char pointer:} strict-aliasing violation. Use \cee{union}, \cee{memcpy}, or \cee{char*}/\cee{uint8\_t*}.
\item \textbf{Misaligned access via cast:} \cee{*(uint32\_t*)(p+1)} on a byte-array p with \cee{p+1} not 4-aligned -- UB (Cortex-M4 mostly works for word but not LDRD/STRD/LDM).
\item \textbf{Sequence points / unsequenced modification:} \cee{i = i++ + ++i;} UB. Side-effects on same scalar between two seq points UB (pre-C11) / unsequenced (C11+) UB.
\item \textbf{Null deref}, \textbf{use-after-free}, \textbf{uninit auto var read} (other than via \cee{unsigned char}).
\item \textbf{Strict-aliasing safe pun:} \cee{uint32\_t v; uint8\_t b[4]; memcpy(b,\&v,4);} -- compiler optimizes to a single LDR.
\item \textbf{Comparing signed/unsigned:} signed becomes unsigned if same rank: \cee{-1 < 1U} false.
\end{itemize}

\subsection{Linker script structure}
\begin{lstlisting}[language={}]
MEMORY {
  FLASH (rx)  : ORIGIN = 0x08000000, LENGTH = 256K
  RAM   (xrw) : ORIGIN = 0x20000000, LENGTH =  32K
}
_estack = ORIGIN(RAM) + LENGTH(RAM);   // top of stack
SECTIONS {
  .isr_vector : { . = ALIGN(4); KEEP(*(.isr_vector)) } >FLASH
  .text       : { *(.text .text.*) *(.rodata .rodata.*)
                  _etext = .;          } >FLASH
  _sidata = LOADADDR(.data);            // init image addr
  .data : {                             // copied at startup
    _sdata = .;
    *(.data .data.*) ;
    _edata = .; } >RAM AT>FLASH
  .bss : {                              // zeroed at startup
    _sbss = .; *(.bss .bss.*) *(COMMON) ; _ebss = .;
  } >RAM
  ._user_heap_stack : { . += _Min_Heap_Size;
                        . += _Min_Stack_Size; } >RAM
}
\end{lstlisting}
Symbols \cee{\_sidata/\_sdata/\_edata/\_sbss/\_ebss/\_estack} consumed by Reset\_Handler.

\subsection{Startup code (Reset\_Handler)}
\begin{lstlisting}[language={}]
Reset_Handler:
  ldr r0, =_estack ; msr msp, r0       /* set MSP */
  ldr r0, =_sdata  ; ldr r1, =_edata
  ldr r2, =_sidata ; mov r3, #0        /* copy .data */
1: cmp r0, r1 ; ittt lt
   ldrlt r4, [r2], #4 ; strlt r4, [r0], #4 ; blt 1b
  ldr r0, =_sbss ; ldr r1, =_ebss      /* zero .bss */
2: cmp r0, r1 ; itt lt
   strlt r3, [r0], #4 ; blt 2b
  bl  SystemInit                       /* clocks/PLL */
  bl  __libc_init_array                /* C++ static ctors */
  bl  main
3: b 3b
\end{lstlisting}
\textbf{Vector table} (.isr\_vector section): word 0 = \cee{\_estack}, word 1 = \cee{Reset\_Handler}, then NMI, HardFault, MemManage, BusFault, UsageFault, (4 reserved), SVCall, DebugMon, (1 reserved), PendSV, SysTick, then 16+ peripheral handlers (vendor list).
\textbf{Default\_Handler} = infinite loop; declared \cee{.weak} alias for every IRQ name so user can override by defining the symbol.

\subsection{GCC ARM toolchain flags}
\begin{lstlisting}[language={}]
CC      = arm-none-eabi-gcc
CFLAGS  = -mcpu=cortex-m4 -mthumb \                # M4 + Thumb-2 ISA
          -mfloat-abi=hard -mfpu=fpv4-sp-d16 \     # use FPU regs (s0-s31)
          -Os -ffunction-sections -fdata-sections\ # per-fn sections (gc)
          -Wall -Wextra -Wno-unused-parameter \
          -ffreestanding -nostartfiles \           # we provide startup
          -DSTM32G431xx -DUSE_HAL_DRIVER \
          -g3 -gdwarf-2
LDFLAGS = -T linker.ld -Wl,--gc-sections \         # drop unused fns
          -Wl,-Map=app.map -specs=nano.specs \     # newlib-nano libc
          -lc -lm -lnosys
\end{lstlisting}
\textbf{Float ABI:} \cee{soft} = library calls (no FPU), \cee{softfp} = FPU compute but args in core regs (compat), \cee{hard} = FPU regs for args (fastest, requires matching libs). Cortex-M4F uses \cee{fpv4-sp-d16} (single-precision, 16 D-regs aliased over 32 S-regs).
\textbf{-Os} for code size; \cee{-O2} or \cee{-O3} for speed (watch flash budget).
\textbf{nano.specs} swaps in newlib-nano (smaller printf without float; add \cee{-u \_printf\_float} to relink float support).

\subsection{Quick numeric reference}
\begin{tabular}{@{}ll@{}}
\toprule
\textbf{Item} & \textbf{Value}\\
\midrule
PERIPH\_BASE & 0x4000\_0000\\
APB1PERIPH\_BASE & 0x4000\_0000\\
APB2PERIPH\_BASE & 0x4001\_0000\\
AHB1PERIPH\_BASE & 0x4002\_0000\\
AHB2PERIPH\_BASE & 0x4800\_0000 (G4) / +0x0800\_0000 (F4)\\
GPIOA\_BASE (G4) & 0x4800\_0000 (port spacing 0x0400)\\
SCS\_BASE & 0xE000\_E000\\
SysTick\_BASE & 0xE000\_E010\\
NVIC\_BASE & 0xE000\_E100\\
SCB\_BASE & 0xE000\_ED00\\
SCB->VTOR & 0xE000\_ED08\\
SCB->AIRCR & 0xE000\_ED0C (PRIGROUP[10:8], VECTKEY=0x05FA)\\
Flash base & 0x0800\_0000\\
SRAM base & 0x2000\_0000\\
SRAM bit-band alias & 0x2200\_0000\\
PERIPH bit-band alias & 0x4200\_0000\\
PSR ExcNum & 16 + IRQn\\
SysTick IRQn & $-1$ (ExcNum 15)\\
\bottomrule
\end{tabular}

\begin{ex}\textbf{Bit-band atomic toggle SRAM bit.} Word at 0x2000\_0040, bit 7. Alias = $0x2200\_0000 + 32 \cdot 0x40 + 4 \cdot 7 = 0x2200\_081C$. Read returns 0 or 1; write 0/1 stores into bit 7 atomically.\end{ex}

\begin{ex}\textbf{USART BRR for 9600 baud, $f_{CK}$=16\,MHz, OVER8=0.} USARTDIV = 16e6/9600 = 1666.67 $\to$ 1667 = 0x683. Actual baud = 16e6/1667 = 9598 (0.02\% err).\end{ex}

\begin{ex}\textbf{NVIC enable USART2\_IRQn=38.} \cee{NVIC->ISER[1] = 1U << 6;} (38>>5=1, 38\&31=6).\end{ex}

\begin{ex}\textbf{Volatile vs not.} Without \cee{volatile}, the compiler may hoist the load of \cee{*flag} out of the spin loop:
\begin{lstlisting}[language={}]
extern int flag;        // BAD: GCC -O2 caches flag in a reg
while (!flag) {}        //      loop becomes infinite if flag was 0
extern volatile int flag;  // GOOD: re-reads each iteration
while (!flag) {}
\end{lstlisting}\end{ex}

\begin{ex}\textbf{Operator precedence sanity.}
\cee{x \& 0xF == 0} parses as \cee{x \& (0xF==0)} = \cee{x \& 0} = 0 -- always 0!
Correct: \cee{(x \& 0xF) == 0}.
\cee{a + b >> 1} = \cee{(a+b)>>1} (shift below \cee{+}). To halve a sum is fine; \cee{a + (b>>1)} requires explicit parens.\end{ex}

\begin{ex}\textbf{Linker symbol lookup.} In Reset\_Handler, \cee{ldr r0, =\_estack} loads the literal-pool word that the linker filled with the symbol's address (= top of RAM). The C-side declaration is \cee{extern uint32\_t \_estack;} and you take its \emph{address}: \cee{\&\_estack}.\end{ex}

% ----- part2/09_arm_reference.tex -----
\section{ARM ASM Quick Reference}

\subsection{Instruction quick ref (immediate / reg-offset / MOV / stack)}
\textbf{Immediate-offset loads/stores} (basic, pre-, post-index):
\begin{tabular}{@{}l l@{}}
\asm{ldr Rd,[Ra,\#k]}     & \asm{Rd=*(Ra+k)}; Ra unchanged \\
\asm{ldr Rd,[Ra,\#k]!}    & Ra$\leftarrow$Ra+k \emph{first}, then load \\
\asm{ldr Rd,[Ra],\#k}     & load from Ra \emph{first}, then Ra+=k \\
\asm{ldrb/h Rd,[...]}     & 8/16b $\to$ 32b, \textbf{zero}-ext \\
\asm{ldrsb/sh Rd,[...]}   & 8/16b $\to$ 32b, \textbf{sign}-ext \\
\asm{ldrd Rl,Rh,[Ra,\#k]} & 8B; lo at lo addr; \textbf{no reg-offset} \\
\asm{str/strh/strb}       & no sign variant; STRB writes \emph{low byte} only \\
\asm{strd Rl,Rh,[Ra,\#k]} & 8B; same forms; \textbf{no reg-offset} \\
\end{tabular}

\textbf{Reg-offset} (\asm{[Ra,Ri,LSL\#n]}, $n=\log_2$elem):
\begin{tabular}{@{}l l@{}}
\asm{LSL\#2}        & word array (\cee{p[i]} on int32) \\
\asm{LSL\#1}        & halfword (int16) \\
\asm{LSL\#0} or none & byte (int8) \\
\asm{ldrsh/sb}      & sign-ext loads (same form) \\
\asm{str/strh/strb} & stores (same form) \\
\multicolumn{2}{@{}l@{}}{\textbf{No \asm{!} or post-idx for reg-offset.} Ra const.}\\
\end{tabular}

\textbf{MOV / pseudo-LDR (constants):}
\begin{tabular}{@{}l l@{}}
\asm{mov Rd,\#imm8r}  & 8b val rotated even (Thumb-2 modimm) \\
\asm{movw Rd,\#imm16} & low halfword; zero-clears top 16 \\
\asm{movt Rd,\#imm16} & top halfword; preserves low 16 \\
\asm{ldr Rd,=expr}    & assembler picks MOV/MOVW+T/literal pool \\
\asm{ldr Rd,=label}   & PC-relative load of address (any 32b) \\
\asm{adr Rd,label}    & near PC-relative addr (no pool) \\
\end{tabular}

\textbf{Stack (Full-Descending; SP$\to$last pushed):}
\begin{tabular}{@{}l l@{}}
\asm{push \{r4,r5,lr\}} & SP-=12; \textbf{lowest reg\# at lowest addr} \\
\asm{pop \{r4,r5,pc\}}  & undoes push; \asm{pop \{pc\}} returns \\
\multicolumn{2}{@{}l@{}}{Brace order is irrelevant; HW always sorts ascending.}\\
\multicolumn{2}{@{}l@{}}{Stem fn: \asm{push \{lr\}} on entry, \asm{pop \{pc\}} to return.}\\
\multicolumn{2}{@{}l@{}}{Leaf fn: just \asm{bx lr}; never trashes LR.}\\
\end{tabular}

\subsection{\asm{Bcc} decoder (full table)}
\begin{tabular}{@{}l l l@{}}
\toprule
mnem & branch if \dots & flags \\
\midrule
\asm{B}    & always (uncond)               & --- \\
\asm{BEQ}  & equal / zero                  & Z=1 \\
\asm{BNE}  & not equal / non-zero          & Z=0 \\
\asm{BHS/BCS} & uns $\ge$ (no borrow)      & C=1 \\
\asm{BLO/BCC} & uns $<$ (borrow)           & C=0 \\
\asm{BMI}  & negative result               & N=1 \\
\asm{BPL}  & non-negative                  & N=0 \\
\asm{BVS}  & signed overflow               & V=1 \\
\asm{BVC}  & no signed overflow            & V=0 \\
\asm{BHI}  & uns $>$                       & C=1, Z=0 \\
\asm{BLS}  & uns $\le$                     & C=0 or Z=1 \\
\asm{BGE}  & sgn $\ge$                     & N=V \\
\asm{BLT}  & sgn $<$                       & N$\ne$V \\
\asm{BGT}  & sgn $>$                       & Z=0, N=V \\
\asm{BLE}  & sgn $\le$                     & Z=1 or N$\ne$V \\
\bottomrule
\end{tabular}
\textbf{Pick by operand type:} signed $\Rightarrow$ EQ/NE/GT/GE/LT/LE; unsigned $\Rightarrow$ EQ/NE/HI/HS/LO/LS. Mixing (e.g.\ \asm{BGT} on uint compare) silently breaks at boundary cases.

\subsection{Pattern catalog: C $\to$ ASM}

\textbf{Integer arithmetic} (R0=a, R1=b, R2=c, R3=d):
\begin{tabular}{@{}l l@{}}
\cee{c=a+b}      & \asm{add r2,r0,r1} \\
\cee{c=a-b}      & \asm{sub r2,r0,r1} \\
\cee{c=-a}       & \asm{rsb r2,r0,\#0}      ($=0-$a) \\
\cee{c=b-a}      & \asm{rsb r2,r0,r1}       (reverse) \\
\cee{c=a*b}      & \asm{mul r2,r0,r1} \\
\cee{c=a*b+d}    & \asm{mla r2,r0,r1,r3}    (Ra last) \\
\cee{c=d-a*b}    & \asm{mls r2,r0,r1,r3} \\
\cee{c=a/b}      & \asm{sdiv r2,r0,r1}      / \asm{udiv} \\
\cee{c=a\%b}     & \asm{sdiv r2,r0,r1; mls r2,r2,r1,r0} \\
\cee{c=a+b+CF}   & \asm{adcs r2,r0,r1}      (after \asm{adds}) \\
\cee{int64=a+b}  & \asm{adds r0,r2; adc r1,r3} (lo,hi) \\
\end{tabular}

\textbf{Shifts / scaling} (Op2 barrel shift, free):
\begin{tabular}{@{}l l@{}}
\cee{2*a}         & \asm{lsl r2,r0,\#1}  \\
\cee{a/2 (uns)}   & \asm{lsr r2,r0,\#1}  \\
\cee{a/2 (sgn)}   & \asm{asr r2,r0,\#1}  \\
\cee{a*4 + a}=5a  & \asm{add r2,r0,r0,lsl \#2} \\
\cee{3*a}         & \asm{add r2,r0,r0,lsl \#1} \\
\cee{a - a/4}=3a/4 & \asm{sub r2,r0,r0,asr \#2} (sgn) \\
\cee{a + a/8}=9a/8 & \asm{add r2,r0,r0,asr \#3} (sgn) \\
\cee{a<<n | a>>(32-n)} & \asm{ror r2,r0,\#(32-n)} \\
\end{tabular}
\textbf{Op2 quirk:} when Op2 uses a shifted register, dest \emph{must} be written explicitly even if same as Rn (\asm{add r0,r0,lsl \#1} is invalid; need \asm{add r0,r0,r0,lsl \#1}).

\textbf{Bit manipulation:}
\begin{tabular}{@{}l l@{}}
\cee{c|=(1<<n)}      & \asm{mov r3,\#1; orr r2,r2,r3,lsl Rn} \\
\cee{c\&=$\sim$(1<<n)} & \asm{mov r3,\#1; bic r2,r2,r3,lsl Rn} \\
\cee{c\^{}=(1<<n)}   & \asm{eor r2,r2,r3,lsl Rn} \\
\cee{tst bit n}      & \asm{tst r2,\#(1<<n)} $\to$ Z=0 if set \\
\cee{c[s+w-1:s]=0}   & \asm{bfc r2,\#s,\#w} (clear field) \\
\cee{c[s+w-1:s]=v[w-1:0]} & \asm{bfi r2,r1,\#s,\#w} \\
\cee{ext uns field}  & \asm{ubfx r0,r1,\#s,\#w} \\
\cee{ext sgn field}  & \asm{sbfx r0,r1,\#s,\#w} \\
\cee{count lead 0}   & \asm{clz r0,r1} \\
\cee{$\sim$x}          & \asm{mvn r0,r1} \\
\cee{a \& $\sim$b}     & \asm{bic r0,r0,r1} \\
\cee{a | $\sim$b}      & \asm{orn r0,r0,r1} \\
\end{tabular}

\textbf{Memory pointer idioms} (R0=v, R1=p, R2=k):
\begin{tabular}{@{}l l@{}}
\cee{v=*p}        & \asm{ldr r0,[r1]} \\
\cee{*p=v}        & \asm{str r0,[r1]} \\
\cee{v=p[k]}      & \asm{ldr r0,[r1,r2,lsl \#2]} \\
\cee{v=*(p+k)}    & \asm{ldr r0,[r1,\#(4*k)]} (k const) \\
\cee{v=*p++}      & \asm{ldr r0,[r1],\#4}    (post-idx) \\
\cee{v=*++p}      & \asm{ldr r0,[r1,\#4]!}   (pre-idx) \\
\cee{*p++=v}      & \asm{str r0,[r1],\#4} \\
\cee{*++p=v}      & \asm{str r0,[r1,\#4]!} \\
\cee{*(p-1)=v}    & \asm{str r0,[r1,\#-4]}   (no wb) \\
\cee{++*p}        & \asm{ldr r0,[r1]; add r0,\#1; str r0,[r1]} \\
\cee{++*p++}      & \asm{ldrb r0,[r1],\#1; add r0,\#1; strb r0,[r1,\#-1]} \\
\cee{(int8) v=*p++}  & \asm{ldrsb r0,[r1],\#1} \\
\cee{(uint16) v=*p++}& \asm{ldrh  r0,[r1],\#2} \\
\end{tabular}

\textbf{Flow control:}
\begin{tabular}{@{}l l@{}}
\cee{if(a==b)$\dots$} & \asm{cmp r0,r1; bne \_else} (NL) \\
\cee{if(a==0)}        & \asm{cbz r0,\_then} (fwd $\le$126B) \\
\cee{if(a!=0)}        & \asm{cbnz r0,\_then} \\
\cee{if(a<0)}         & \asm{cmp r0,\#0; blt \_then} or \asm{tst}+\asm{bmi} \\
\cee{while(*p)}      & \asm{ldrb r2,[r1],\#1; cbz r2,\_end; ...; b lp} \\
\cee{return} (leaf)   & \asm{bx lr} \\
\cee{return} (stem)   & \asm{pop \{pc\}} (paired with entry \asm{push \{lr\}}) \\
\cee{int64 ret}       & lo$\to$R0, hi$\to$R1 \\
\end{tabular}

\subsection{Common recipes}

\textbf{Saturate signed to $[L,H]$:}
\begin{lstlisting}
@ R0 = clip(R0, R1=L, R2=H)
    cmp  r0, r1
    it   lt
    movlt r0, r1                   @ if (r0<L) r0=L
    cmp  r0, r2
    it   gt
    movgt r0, r2                   @ if (r0>H) r0=H
    bx   lr
\end{lstlisting}

\textbf{Array sum (int32, len R1):}
\begin{lstlisting}
    mov  r2, #0                    @ sum
    cbz  r1, sm_done
sm_lp:
    ldr  r3, [r0], #4
    add  r2, r2, r3
    subs r1, r1, #1
    bne  sm_lp
sm_done:
    mov  r0, r2
    bx   lr
\end{lstlisting}

\textbf{Conditional sum (negatives only):}
\begin{lstlisting}
    mov  r2, #0
cs_lp:
    cbz  r1, cs_done
    ldr  r3, [r0], #4
    cmp  r3, #0
    it   lt
    addlt r2, r2, r3               @ accumulate iff r3<0
    sub  r1, r1, #1
    b    cs_lp
cs_done:
    mov  r0, r2
    bx   lr
\end{lstlisting}

\textbf{Array max (signed):}
\begin{lstlisting}
    ldr  r2, [r0], #4              @ seed = a[0]
    sub  r1, r1, #1
mx_lp:
    cbz  r1, mx_done
    ldr  r3, [r0], #4
    cmp  r3, r2
    it   gt
    movgt r2, r3                   @ signed max
    sub  r1, r1, #1
    b    mx_lp
mx_done:
    mov  r0, r2
    bx   lr
\end{lstlisting}

\textbf{strlen:}
\begin{lstlisting}
    mov  r1, r0                    @ keep base
sl_lp:
    ldrb r2, [r0], #1
    cmp  r2, #0
    bne  sl_lp
    sub  r0, r0, r1                @ end - base
    sub  r0, r0, #1                @ exclude '\0'
    bx   lr
\end{lstlisting}

\textbf{strcpy (returns dst):}
\begin{lstlisting}
    mov  r2, r0                    @ save dst
sc_lp:
    ldrb r3, [r1], #1
    strb r3, [r0], #1              @ writes '\0' too
    cmp  r3, #0
    bne  sc_lp
    mov  r0, r2
    bx   lr
\end{lstlisting}

\textbf{memcpy byte-wise (R0=dst,R1=src,R2=n):}
\begin{lstlisting}
    cbz  r2, mc_done
mc_lp:
    ldrb r3, [r1], #1
    strb r3, [r0], #1
    subs r2, r2, #1
    bne  mc_lp
mc_done:
    bx   lr
\end{lstlisting}

\textbf{Q15$\times$Q15$\to$Q15} (renorm $\gg 15$ via 32b chunk):
\begin{lstlisting}
@ R0=a (Q15 in low 16, sign-ext), R1=b
    smull r2, r3, r0, r1           @ r3:r2 = a*b (Q30, 64b)
    lsr   r2, r2, #15              @ shift fractional
    orr   r0, r2, r3, lsl #17      @ stitch carry from hi
    sxth  r0, r0                   @ clamp/keep low 16 (Q15)
    bx    lr
\end{lstlisting}

\textbf{Pointer difference (typed):}
\begin{lstlisting}
@ int32_t *pa,*pb;  ret = pa - pb (in elements)
    sub  r0, r0, r1                @ raw byte diff
    asr  r0, r0, #2                @ /4 because sizeof(int32)=4
    bx   lr
@ for int16: asr #1; int8: no shift; int64: asr #3.
\end{lstlisting}

\textbf{Bit-field assign / extract:}
\begin{lstlisting}
@ struct { unsigned mode:3; unsigned val:5; ... } @bits 0..7
    bfi  r0, r1, #0, #3            @ mode = r1[2:0]
    bfi  r0, r2, #3, #5            @ val  = r2[4:0]
    ubfx r3, r0, #3, #5            @ uns extract val
    sbfx r4, r0, #0, #3            @ sgn extract mode
\end{lstlisting}

\subsection{Last-Minute Traps (extended)}
\begin{itemize}
\item \textbf{Sign-ext on byte/half MSbit=1.} \asm{ldrsb} of \asm{0x80}$\to$\asm{0xFFFFFF80}; \asm{ldrsh} of \asm{0x8000}$\to$\asm{0xFFFF8000}. \asm{LDRB/LDRH} always zero-ext. Pick suffix from C type, \emph{not} from data values.
\item \textbf{Sign-ext at the call site.} \cee{int8\_t a=-1} arg $\Rightarrow$ caller widens to \asm{R0=0xFFFFFFFF}. Callee that reads R0 as \asm{uint32\_t} sees a huge unsigned. EABI widens narrow types BEFORE the BL.
\item \textbf{\asm{STRD} has no reg-offset form.} Only \asm{[Ra,\#k]}, \asm{[Ra,\#k]!}, \asm{[Ra],\#k}. Same restriction on \asm{LDRD}.
\item \textbf{\asm{CBZ}/\asm{CBNZ} forbidden inside IT.} They don't read flags; using them in IT is a hard assemble error. \textbf{Conditional B inside IT only legal as the FINAL slot.}
\item \textbf{\asm{cmp r0,\#-5} re-encodes as \asm{CMN r0,\#5}.} Output is identical for signed CCS, but the disassembly will read \asm{CMN}; don't be confused.
\item \textbf{Prefix \asm{++} on \asm{*p} requires explicit STR writeback.} \cee{++*p} = LDR, ADD \#1, \emph{STR back}. The assembler/CPU does not auto-write-back the value side; only pre-/post-index handles \emph{pointer} writeback.
\item \textbf{Pointer arithmetic is scaled by sizeof.} \cee{p++} on \cee{int32\_t*} adds 4; on \cee{int16\_t*} adds 2; on \cee{int8\_t*} adds 1; on \cee{int64\_t*} adds 8. Diff between two typed ptrs returns \emph{elements}, not bytes.
\item \textbf{Flag update needs the \asm{S} suffix.} \asm{add}$\ne$\asm{adds}; \asm{sub}$\ne$\asm{subs}; \asm{mov}$\ne$\asm{movs}. Plain forms leave NZCV alone -- read flags immediately after \asm{cmp}/\asm{cmn}/\asm{tst}/\asm{teq} or an \asm{*S} op.
\item \textbf{\asm{UDIV}/\asm{SDIV} \emph{never} set flags.} No \asm{UDIVS} exists. If you need a flag after a divide, follow with \asm{cmp Rd,\#0}.
\item \textbf{ARM \asm{C}-on-subtract is the \emph{inverse} of x86.} \asm{C=1} after \asm{SUB/CMP} means \emph{no borrow} (i.e.\ Rn$\ge$Op2 unsigned). \asm{BHS/BCS}=$\ge$, \asm{BLO/BCC}=$<$. Don't reuse x86 borrow intuition.
\item \textbf{GNU shifted-Op2 quirk.} When Op2 is a register with a shift, the destination must be written explicitly even if it equals Rn. \asm{add r0,r0,lsl \#1} fails; use \asm{add r0,r0,r0,lsl \#1}.
\item \textbf{Q-format mul renorm.} Q$m.n$$\times$Q$m.n$ produces Q$2m.2n$ in 64b; renormalize back via \asm{ASR \#n} (signed) on the low/high stitch. For Q15$\to$Q15: \asm{smull;lsr \#15} after combining hi/lo halves.
\item \textbf{PUSH order is always ascending reg\#.} \asm{push \{r7,r5,r8,r6\}} stores R5 at the lowest address, R8 at the highest, regardless of brace order. POP is the inverse.
\item \textbf{\asm{STRB} silently drops upper 24 bits.} \asm{strb r0,[r1]} writes only \asm{r0[7:0]}. Hence no \asm{STRSB}: store-side sign info is meaningless.
\item \textbf{8-byte SP alignment required at every \asm{BL}.} If your prologue pushes an odd number of registers (incl.\ LR), pad with a dummy reg (e.g.\ \asm{push \{r4,lr\}} = 2 regs OK; \asm{push \{lr\}} alone violates -- pair with one more). AAPCS rule.
\item \textbf{Address bit0 in \asm{BX}/\asm{BLX} = Thumb-mode flag.} The PC always sees an even address; the CPU strips bit0 to set the T-bit. Don't manually clear it from a function pointer.
\item \textbf{Literal pool reach.} \asm{ldr Rd,=expr} works when there's a reachable pool ($\le 4095$B). For deep functions, drop \asm{.ltorg} or split.
\end{itemize}

\subsection{One-page summary block}
\textbf{EABI:} R0--R3 args/return, caller-saved. R4--R11 callee-saved. R12=IP scratch. R13=SP, R14=LR, R15=PC. 64b vals: low$\to$R0, high$\to$R1. Stack=Full-Descending; SP 8B-aligned at \asm{BL}. Leaf: \asm{bx lr}; stem: \asm{push \{lr\}}\dots\asm{pop \{pc\}}. \textbf{Loads:} basic \asm{[Ra,\#k]} no wb; pre \asm{[Ra,\#k]!}=\cee{*++p}; post \asm{[Ra],\#k}=\cee{*p++}. Suffixes: \asm{LDR}=32b, \asm{LDRH/B}=zero-ext, \asm{LDRSH/SB}=sign-ext. STRs: no sign variants; STRB drops upper 24b. \asm{LDRD/STRD} no reg-offset. Reg-offset: \asm{[Ra,Ri,LSL \#n]}, $n=\log_2$elem. \textbf{Constants:} \asm{ldr=expr} picks MOV/MOVW+T/pool. \textbf{Flags:} \asm{S}-suffix sets NZCV. \asm{cmp}=SUBS no-wb; \asm{cmn}=ADDS no-wb; \asm{tst}=ANDS no-wb; \asm{teq}=EORS no-wb. \asm{UDIV/SDIV} never set flags. C-on-SUB = no-borrow. \textbf{CCS pick:} signed$\to$EQ/NE/GT/GE/LT/LE; unsigned$\to$EQ/NE/HI/HS/LO/LS. \textbf{Cond exec:} IT/ITT/ITE/ITTT/ITTE/ITEE/ITTEE/ITTTT, max 4 slots, T=match cond, E=inverse, cond \asm{B} only as last slot, no \asm{CBZ} inside. \textbf{CBZ/CBNZ:} fwd $\le$126B, R0--R7, flag-neutral, ideal post-LDRB null test. \textbf{Loops:} top-test (\asm{while},\asm{for}) $\ge 0$ iters; bottom-test (\asm{do-while}) $\ge 1$. Priming-while and \asm{while(1)+break} have same dynamic LDR count, differ only in source duplication. \textbf{Switch:} CMP/BEQ ladder ($O(n)$) / stacked fall-through (shared body, no code between labels) / \asm{TBB [Rn,Rm]} byte table $O(1)$ / \asm{TBH [Rn,Rm,LSL \#1]} halfword table for far cases. \textbf{Q-fmt:} renorm by ASR $n$; Q15 mul = \asm{smull}+\asm{lsr \#15}. \textbf{Bit-fields:} \asm{BFC/BFI/UBFX/SBFX} eliminate hand-coded shift+mask sequences. \textbf{Pointer scaling}: \cee{p++} adds sizeof(*p). Diff returns elements. \textbf{Sign-ext at BL:} narrow signed types widen to 32b before placement in R0--R3. \textbf{Hard fails:} \asm{cmp Rn,\#-imm}$\to$CMN; CBZ in IT; reg-offset \asm{!}; STRD reg-offset; \asm{add r0,r0,lsl \#1} (need 3-arg form); flagless ops where flags expected. When in doubt: type drives suffix; suffix drives sign-ext; signed/unsigned drives CCS \emph{and} shift kind (ASR vs LSR).

% ----- part2/10_hw1_4_extended_drills.tex -----
% =============================================================================
% HW1-4 Extended Drills + Practice Problems (Part 2 add-on)
% =============================================================================
\section{HW1-4 Extended Drills + Practice Problems}

\subsection{HW1 Deep Walkthrough --- UART throughput (P1)}
\textbf{Setup:} 115200~bps, 1 start + 8 data + 1 parity + 1 stop = \textbf{11 bits/frame}, payload = 8 data bits = 1~byte.\\
\textbf{Step 1.} Frame rate $=$ 115200 / 11 $=$ 10472.727 frames/s.\\
\textbf{Step 2.} Each frame carries 1 payload byte $\Rightarrow$ throughput $=$ \textbf{10472.7~B/s}.\\
\textbf{Trap:} dividing by 8 (data bits only) gives 14400~B/s $-$ wrong; the line still ships start/parity/stop overhead.\\
\textbf{General formula:} $T_B = \dfrac{\text{baud}}{N_{\text{start}}+N_{\text{data}}+N_{\text{par}}+N_{\text{stop}}}\cdot\dfrac{N_{\text{data}}}{8}$ B/s when payload is whole byte. For 8N1: $\text{baud}/10$. For 8N2: $\text{baud}/11$. For 9600 8N1: 960~B/s. For 230400 8N1: 23040~B/s.

\subsection{HW1 Deep Walkthrough --- Pipeline loop (P3)}
3-stage pipe (F-D-E). Loop body = 5 task instr + 1 branch; branch FLUSHES (drops 2 in-flight stages).\\
\textbf{Single iter cycles:} pipe fill $=$ depth$-1$ $=$ 2 cyc; then 1 cyc/instr after fill. So 6 instr $\Rightarrow$ $2 + 6 = 8$ cyc. Throughput $=$ 6/8 $=$ \textbf{0.75 instr/cyc}.\\
\textbf{Unroll 4x} (5 task)$\times 4 + $1 branch $= 21$ instr/iter $\Rightarrow$ $2 + 21 = 23$ cyc. Throughput $=$ 21/23 $=$ \textbf{0.913 instr/cyc}.\\
\textbf{Why unrolling helps:} each branch costs depth$-1$ flush cycles, amortized over more useful work. Unroll by $K$: cycles $=(K\cdot 5 + 1)+2$, throughput $\to 1$ as $K\to\infty$.\\
\textbf{General:} cycles $= n + (\text{depth}-1)$ for straight code; throughput $\to 1$ as $n\to\infty$. 4-stage: $n+3$. 5-stage: $n+4$.
\begin{hwp}\textbf{P4:} 3-stage, 10 instr $\Rightarrow$ $10+2=12$ cyc, $10/12=$0.833.\end{hwp}

\subsection{HW1 Deep Walkthrough --- Shift conversions (P5/P6)}
\textbf{Decimal $\to$ binary} (divide-by-2, LSB first): $28: 28/2{=}14r0; 14/2{=}7r0; 7/2{=}3r1; 3/2{=}1r1; 1/2{=}0r1$ $\Rightarrow$ \cee{0b11100} (read remainders bottom$\to$top).\\
$28=$\cee{0b0001\_1100}; $14=$\cee{0b0000\_1110}. Shift relation: \cee{14<<1 = 28}. Confirm: \cee{0000\_1110 << 1 = 0001\_1100}.\\
$24=$\cee{0b0001\_1000}, $6=$\cee{0b0000\_0110}; \cee{24>>2=6}, so \textbf{b=a>>2}. Each \cee{>>1} divides by 2; \cee{<<1} multiplies by 2 (unsigned).
\begin{hwp}\textbf{P2 parity:} \cee{0b1101\_1100} has 5 ones (odd); even-par bit $=$ \textbf{1}; odd-par bit $=$ \textbf{0}.\end{hwp}

\subsection{HW2 5-bit Two's-Complement Drill}
Cardinality $2^5=32$. Range $[-16,+15]$. Encode neg: $v+32$. Decode: if MSB$=1$, subtract 32.
\begin{tabular}{@{}rlll@{}}
\toprule
dec & enc calc & binary (5b) & hex \\
\midrule
$-16$ & $-16+32=16$ & \cee{0b10000} & \cee{0x10} \\
$-15$ & $-15+32=17$ & \cee{0b10001} & \cee{0x11} \\
$-12$ & $-12+32=20$ & \cee{0b10100} & \cee{0x14} \\
$-8$  & $-8+32=24$  & \cee{0b11000} & \cee{0x18} \\
$-1$  & $-1+32=31$  & \cee{0b11111} & \cee{0x1F} \\
$\;0$ & $0$         & \cee{0b00000} & \cee{0x00} \\
$+1$  & $1$         & \cee{0b00001} & \cee{0x01} \\
$+7$  & $7$         & \cee{0b00111} & \cee{0x07} \\
$+15$ & $15$        & \cee{0b01111} & \cee{0x0F} \\
\bottomrule
\end{tabular}\\
Decode drill: \cee{0b11010}$=26\to 26{-}32=$\textbf{-6}; \cee{0b10110}$=22\to{-}10$; \cee{0b01101}$=13\to+13$.\\
\textbf{Sign-extend} $-6$ from 5b to 8b: repeat MSB $\Rightarrow$ \cee{0b1111\_1010 = 0xFA}. As 16b: \cee{0xFFFA}. As 32b: \cee{0xFFFF\_FFFA}.

\subsection{HW2 Endianness Drill (mem @\cee{0x2000\_0000})}
Bytes (low$\to$high addr): \cee{[0x01, 0x02, 0x03, 0x04, 0x05, 0x06, 0x07, 0x08]}.
\begin{tabular}{@{}llll@{}}
\toprule
addr off & type & bytes used & LE value \\
\midrule
+0 & uint8\_t  & \{0x01\}              & \cee{0x01} \\
+0 & int8\_t   & \{0x01\}              & $+1$ \\
+1 & uint8\_t  & \{0x02\}              & \cee{0x02} \\
+0 & uint16\_t & \{0x01,0x02\}         & \cee{0x0201} \\
+2 & uint16\_t & \{0x03,0x04\}         & \cee{0x0403} \\
+0 & int16\_t  & \{0x01,0x02\}         & $+513$ \\
+0 & uint32\_t & \{0x01..0x04\}        & \cee{0x04030201} \\
+4 & uint32\_t & \{0x05..0x08\}        & \cee{0x08070605} \\
+4 & int32\_t  & \{0x05..0x08\}        & $+134678021$ \\
\bottomrule
\end{tabular}\\
\textbf{Unaligned trap:} \cee{*(uint32\_t*)0x20000001} reads 4 bytes from offset 1 $\to$ \cee{0x05040302}. On Cortex-M3/4 unaligned word access is \emph{allowed} (LDR/STR) but slower; on M0/M0+ it raises UsageFault. Halfword reads need addr\%2==0; word reads addr\%4==0 for guaranteed single-cycle.\\
\textbf{Rule:} multi-byte read = (high addr byte)$\cdot 2^{8(N-1)}$ + ... + (low addr byte). LE places \emph{LSB at lowest address}, so the low-address byte is the lowest-order byte of the assembled value.

\subsection{HW2 Union Aliasing --- Concrete Hex}
\begin{lstlisting}[language=]
typedef union { uint8_t b[4]; uint16_t h[2]; uint32_t w; } u_t;
u_t u;
u.w = 0xDEADBEEF;
// LE bytes: b[0]=0xEF, b[1]=0xBE, b[2]=0xAD, b[3]=0xDE
// halfwords: h[0]=0xBEEF, h[1]=0xDEAD
\end{lstlisting}
\textbf{Write u8s, read u16:} after \cee{b[0]=0x12; b[1]=0x34;} $\Rightarrow$ \cee{h[0] = 0x3412} (LE assembles low addr as low byte).\\
\textbf{Write u16, read u8s:} after \cee{h[0]=0xCAFE;} $\Rightarrow$ \cee{b[0]=0xFE, b[1]=0xCA}.\\
\textbf{Partial overwrite:} starting from \cee{w=0x12345678}, then \cee{b[0]=0xAB} $\Rightarrow$ \cee{w=0x123456AB} (only LSB changed).\\
\textbf{Type punning safe via union} (C99+); reinterpret\_cast via pointer cast may break strict-aliasing rules but works in MCU code with \cee{volatile}/compiler flags.

\subsection{HW3 Register-Access Mapping --- Configure PA5 as Output PP, Low-Speed, No-Pull}
GPIOA base = \cee{0x4800\_0000} (AHB2). Pin 5 uses bits $[2\cdot 5{:}2\cdot 5{+}1] = [10{:}11]$ in 2-bit fields, bit 5 in 1-bit fields.\\
\textbf{Reset state} (G4 GPIOA): MODER=\cee{0xABFF\_FFFF} (most pins analog=0b11), OTYPER=\cee{0x0}, OSPEEDR=\cee{0x0}, PUPDR=\cee{0x6400\_0000}.
\begin{lstlisting}[language=]
// 1. Enable AHB2 GPIOA clock first (in RCC):
RCC->AHB2ENR |= RCC_AHB2ENR_GPIOAEN;
// 2. MODER: clear pin5 (bits 11:10), write 01 (output)
GPIOA->MODER  &= ~(3U << (5*2));   // clr bits 11:10
GPIOA->MODER  |=  (1U << (5*2));   // set 0b01
// 3. OTYPER: bit5 = 0 (push-pull)  -- already 0 at reset
GPIOA->OTYPER &= ~(1U << 5);
// 4. OSPEEDR: clear pin5 field, write 00 (low speed)
GPIOA->OSPEEDR &= ~(3U << (5*2));  // 0b00 = low
// 5. PUPDR: clear pin5 field, write 00 (no pull)
GPIOA->PUPDR  &= ~(3U << (5*2));
\end{lstlisting}
\textbf{Before/after} (assume reset values, change only pin 5):
\begin{tabular}{@{}lll@{}}
\toprule
reg & before & after \\
\midrule
MODER  (0x00) & \cee{0xABFF\_FFFF} & \cee{0xABFF\_F7FF} \\
OTYPER (0x04) & \cee{0x0000\_0000} & \cee{0x0000\_0000} \\
OSPEEDR(0x08) & \cee{0x0000\_0000} & \cee{0x0000\_0000} \\
PUPDR  (0x0C) & \cee{0x6400\_0000} & \cee{0x6400\_0000} \\
\bottomrule
\end{tabular}\\
\textbf{Drive PA5 high (atomic):} \cee{GPIOA->BSRR = (1U<<5);} (no read-modify-write).

\subsection{HW3 Bit-Field Replace Drill (5 examples)}
Pattern: \cee{R \&= \~{}(MASK<<POS); R |= ((VAL \& MASK)<<POS);} where MASK $=$ \cee{(1U<<W)-1}.
\begin{tabular}{@{}lllll@{}}
\toprule
\# & field & W & POS & MASK \\
\midrule
1 & MODER pin 7  & 2 & 14 & \cee{0x0000C000} \\
2 & EXTICR line3 & 4 & 12 & \cee{0x0000F000} \\
3 & PUPDR pin 12 & 2 & 24 & \cee{0x03000000} \\
4 & ADC SQR1 SQ1 & 5 & 6  & \cee{0x000007C0} \\
5 & RCC PLLM     & 4 & 4  & \cee{0x000000F0} \\
\bottomrule
\end{tabular}
\begin{lstlisting}[language=]
// 1. Set MODER pin7 to AF (0b10):
GPIOA->MODER   &= ~(3U  <<14);  GPIOA->MODER   |= (2U  <<14);
// 2. Route EXTI3 to port D (code 3):
SYSCFG->EXTICR[0] &= ~(0xFU<<12); SYSCFG->EXTICR[0] |= (3U<<12);
// 3. Pull-down on PA12 (0b10):
GPIOA->PUPDR   &= ~(3U  <<24);  GPIOA->PUPDR   |= (2U  <<24);
// 4. ADC SQR1 first conversion = ch 17:
ADC1->SQR1     &= ~(0x1FU<<6);  ADC1->SQR1     |= (17U<<6);
// 5. RCC PLLM = 8 (encoded as 7):
RCC->PLLCFGR   &= ~(0xFU <<4);  RCC->PLLCFGR   |= (7U  <<4);
\end{lstlisting}
\textbf{One-shot form:} \cee{R = (R \& \~{}(M<<P)) | ((V\&M)<<P);} avoids glitch window between two writes (matters for live timer/PWM regs).

\subsection{HW4 EXTI Configuration Trace --- PA0 Falling-Edge IRQ}
\textbf{IRQn:} EXTI0 $\Rightarrow$ IRQn $=$ \textbf{6}; PSR ExcNum = $16+6 = 22$.\\
\textbf{NVIC index/bit:} idx $=$ 6/32 $=$ 0; bit $=$ 6\%32 $=$ 6 $\Rightarrow$ \cee{ISER[0]} bit 6.
\begin{lstlisting}[language=]
// 0. clocks
RCC->AHB2ENR |= RCC_AHB2ENR_GPIOAEN;
RCC->APB2ENR |= RCC_APB2ENR_SYSCFGEN;
// 1. PA0 input, no-pull
GPIOA->MODER &= ~(3U<<0);          // 00 = input
GPIOA->PUPDR &= ~(3U<<0);          // 00 = no pull
// 2. Route EXTI0 to PA: SYSCFG_EXTICR1 bits[3:0] = 0b0000
SYSCFG->EXTICR[0] &= ~(0xFU<<0);   // clr -> port A code 0
// 3. Falling-edge trigger only
EXTI->RTSR &= ~(1U<<0);            // disable rising
EXTI->FTSR |=  (1U<<0);            // enable falling
// 4. Unmask IRQ
EXTI->IMR  |=  (1U<<0);
// 5. Clear stale pending in NVIC, then enable
NVIC->ICPR[0] = (1U<<6);
NVIC->ISER[0] = (1U<<6);
\end{lstlisting}
\textbf{Final reg state} (only EXTI0-related bits):
\begin{tabular}{@{}lll@{}}
\toprule
reg & contents (relevant) & note \\
\midrule
SYSCFG\_EXTICR1 & \cee{...0000} bits[3:0] & port A=0 \\
EXTI\_RTSR & bit0=0 & no rising \\
EXTI\_FTSR & bit0=1 & falling enabled \\
EXTI\_IMR  & bit0=1 & IRQ unmasked \\
NVIC\_ISER[0] & bit6=1 & EXTI0 enabled \\
\bottomrule
\end{tabular}
\textbf{Handler:}
\begin{lstlisting}[language=]
void EXTI0_IRQHandler(void){
  if(EXTI->PR & (1U<<0)){
    EXTI->PR = (1U<<0);   // W1C clear pending
    /* user logic */
  }
}
\end{lstlisting}

\subsection{HW4 Priority + Nested ISR Drill}
\textbf{Setup:} AIRCR PRIGROUP $=$ 4 (3 PPN bits, 1 SPN bit, N$=$4 implemented). Priorities (in IP[7:4] form, raw byte): IRQ\_A $=$ \cee{0x20}, IRQ\_B $=$ \cee{0x40}, IRQ\_C $=$ \cee{0x60}.\\
\textbf{Decode:} drop low 4 (unimplemented) $\Rightarrow$ PN = byte$>>$4. PRIGROUP$=$4 with N$=$4 means 3 PPN, 1 SPN: PPN $=$ PN$>>$1, SPN $=$ PN\&1.
\begin{tabular}{@{}lllll@{}}
\toprule
IRQ & raw IP & PN & PPN & SPN \\
\midrule
A & \cee{0x20} & 2 & 1 & 0 \\
B & \cee{0x40} & 4 & 2 & 0 \\
C & \cee{0x60} & 6 & 3 & 0 \\
\bottomrule
\end{tabular}\\
\textbf{(a) Preemption matrix:} A preempts B (1$<$2) yes; A preempts C (1$<$3) yes; B preempts C (2$<$3) yes; B preempts A no (2$\not<$1); C preempts neither.\\
\textbf{(b) Highest priority:} \textbf{IRQ\_A} (smallest PN=2).\\
\textbf{(c) Set IRQ\_X to PN=4:} write \cee{IP[X] = 4<<(8-4) = 4<<4 = 0x40 = 0b0100\_0000}. Equivalent CMSIS call: \cee{NVIC\_SetPriority(X, 4);} which masks to N bits and shifts.\\
\textbf{Nesting trace:} B running, A asserts $\to$ HW pushes 8-word frame (R0-R3, R12, LR, PC, xPSR), vectors to A. A returns $\to$ unstacks back to B context. C asserting during B does \emph{not} preempt (3$\not<$2) $\to$ stays pending until B returns.

\subsection{Memory Map Quick Reference}
\begin{tabular}{@{}lll@{}}
\toprule
addr & region & contents \\
\midrule
\cee{0x0800\_0000} & FLASH  & code (vector table aliased here) \\
\cee{0x2000\_0000} & SRAM   & .data, .bss, heap, stack(top) \\
\cee{0x4000\_0000} & APB1   & TIM2-7, USART2-3, I2C, etc. \\
\cee{0x4001\_0000} & APB2   & TIM1, USART1, SPI1, SYSCFG \\
\cee{0x4002\_0000} & AHB1   & DMA, RCC, FLASH ctrl, CRC \\
\cee{0x4800\_0000} & AHB2   & GPIOA..GPIOH (G4 family) \\
\cee{0xE000\_E000} & SCS    & (private peripheral bus) \\
\cee{0xE000\_E010} & SysTick & CTRL/LOAD/VAL/CALIB \\
\cee{0xE000\_E100} & NVIC   & ISER/ICER/ISPR/ICPR/IABR/IP \\
\cee{0xE000\_ED00} & SCB    & CPUID/ICSR/VTOR/AIRCR/SCR \\
\bottomrule
\end{tabular}

\subsection{Standalone Practice Problems}

\begin{ex}\textbf{Pipeline math.} 4-stage ideal pipe; loop body has 6 instr including 2 unconditional branches that flush. Throughput?\\
\textbf{Sol.} Each branch wastes (depth$-1$)$=$3 cyc. Useful instr/iter $=$ 6; cyc/iter $=$ 6 (1 cyc each in fill steady-state) $+$ 2$\cdot$3 flushes $=$ 12. Throughput $=$ 6/12 $=$ \textbf{0.5 instr/cyc}.\end{ex}

\begin{ex}\textbf{Two's complement encode.} Encode $-100$ in 8-bit TC.\\
\textbf{Sol.} $-100 + 256 = 156 =$ \cee{0b1001\_1100} $=$ \cee{0x9C}. Verify: MSB$=1$, $156-256 = -100$. \checkmark
\end{ex}

\begin{ex}\textbf{Two's complement decode.} Decode \cee{0xC8} as 8-bit signed.\\
\textbf{Sol.} \cee{0xC8} $=$ 200; MSB$=1$ $\Rightarrow$ $200-256 = $ \textbf{-56}. Bits: \cee{0b1100\_1000}; magnitude via invert+1: \cee{\~{}0xC8 = 0x37 = 55}; $+1=56$, sign neg $\Rightarrow -56$. \checkmark
\end{ex}

\begin{ex}\textbf{GPIO BSRR atomic.} Write the BSRR value to set PB7 and reset PB12 in one store.\\
\textbf{Sol.} BSRR low half (bits 15:0) = BS0..BS15 (set), high half (31:16) = BR0..BR15 (reset). Set PB7: bit 7 $=$ \cee{0x0080}. Reset PB12: bit $12+16=28$ $=$ \cee{0x1000\_0000}. OR them: \cee{GPIOB->BSRR = 0x1000\_0080;}\end{ex}

\begin{ex}\textbf{NVIC ISER for USART2 (IRQn=38).} Which ISER index/bit? What's its address?\\
\textbf{Sol.} idx $=$ 38/32 $=$ \textbf{1}; bit $=$ 38\%32 $=$ \textbf{6}. ISER[0] is at NVIC base 0xE000\_E100, each entry 4 bytes, so ISER[1] at \cee{0xE000\_E104}. C: \cee{NVIC->ISER[1] = (1U<<6);} or equivalently \cee{*(volatile uint32\_t*)0xE000E104 = 0x40;}\end{ex}

\begin{ex}\textbf{Memory map.} Identify peripherals at given addresses.\\
\textbf{Sol.} \cee{0x4002\_0000} = AHB1 base (RCC, DMA, FLASH ctrl region; specifically DMA1 on G4). \cee{0x4001\_0000} = APB2 base (SYSCFG, COMP, EXTI on G4 sit in this region). \cee{0xE000\_ED00} = SCB base (CPUID/ICSR/VTOR/AIRCR/SCR/CCR/SHPR/SHCSR).\end{ex}

\begin{ex}\textbf{Fixed-point Q1.7.} Encode 1.75 in Q1.7.\\
\textbf{Sol.} Q1.7 $=$ 1 sign + 1 int + 7 frac $=$ 9 bits typ stored in 16/8. Range $[-2,+2)$, resolution $2^{-7}$. Encode: $1.75 \cdot 2^7 = 224 =$ \cee{0b0\_1110\_0000} $=$ \cee{0x0E0}. As 16b stored: \cee{0x00E0}. Decode check: 224/128 $=$ 1.75. \checkmark\end{ex}

\begin{ex}\textbf{Pointer cast.} Memory at \cee{0x2000\_0004} = \cee{[0x01,0x02,0x03,0x04,0x05,0x06]}. Evaluate \cee{*(uint16\_t *)0x20000004}.\\
\textbf{Sol.} Cast addr to \cee{uint16\_t*}, dereference reads 2 LE bytes from \cee{0x20000004,0x20000005} $=$ \cee{0x01,0x02} $\Rightarrow$ value $=$ \cee{0x0201} = 513.\end{ex}

\subsection{Rapid-Fire Reference Card}
\begin{tabular}{@{}ll@{}}
\toprule
quantity & value \\
\midrule
8N1 byte time @115200 & 86.8~$\mu$s \\
8N1 throughput @115200 & 11520~B/s \\
3-stage pipe, $n$ instr & $n+2$ cyc \\
4-stage pipe, $n$ instr & $n+3$ cyc \\
ISR auto-stack cost & 8 cyc (push) + 8 cyc (pop) \\
EXTI0\_IRQn / PSR & 6 / 22 \\
USART2\_IRQn & 38 \\
SysTick\_IRQn & $-1$ \\
TC range, $N$ bits & $[-2^{N-1},\,2^{N-1}{-}1]$ \\
LE rule & LSB at lowest addr \\
GPIO MODER bits/pin & 2 (00 in, 01 out, 10 AF, 11 ana) \\
GPIO BSRR layout & \cee{[BR15..0\;|\;BS15..0]} \\
NVIC index & IRQn$\gg$5 \\
NVIC bit & IRQn\&31 \\
IP shift (N=4) & val$\ll$4 \\
\bottomrule
\end{tabular}

% ----- part2/11_reverse_eng_internals.tex -----
% =============================================================================
% Part 2 / 11 -- ASM-to-C reverse-engineering drills + System internals + Traps
% =============================================================================
\section{ASM Reverse-Engineering + System Internals + Exam-Day Traps}

\subsection{ASM $\to$ C drills (decode the loop, name the idiom)}

\textbf{R1: increment-and-store buffer fill.}
\begin{lstlisting}
@ R0 = p (uint32_t*), R1 = n
    cbz   r1, R1_end
R1_lp:
    str   r1, [r0], #4              @ *p++ = n
    subs  r1, r1, #1
    bne   R1_lp
R1_end:
    bx    lr
\end{lstlisting}
\cee{void f(uint32\_t *p,uint32\_t n)\{ while(n)\{ *p++=n; n--; \} \}} -- post-index store, decrement-and-branch loop. Final \cee{*p} value is 1, not 0.

\textbf{R2: conditional accumulator (IT block).}
\begin{lstlisting}
@ R0=arr, R1=n; ret R0 = sum of positive elems
    movs  r2, #0
R2_lp:
    cbz   r1, R2_done
    ldr   r3, [r0], #4
    cmp   r3, #0
    it    gt
    addgt r2, r2, r3                @ r2 += r3 iff r3 > 0
    subs  r1, r1, #1
    b     R2_lp
R2_done:
    mov   r0, r2
    bx    lr
\end{lstlisting}
\cee{int sum\_pos(int *a,int n)\{int s=0; while(n--)\{ if(*a>0) s+=*a; a++;\} return s;\}}. Single-slot IT; \cee{cmp} sets N/V/Z used by \cee{GT}. Note: writing \cee{addgt} \emph{outside} a guarding IT is illegal in Thumb-2.

\textbf{R3: 64-bit add (lo,hi pairs).}
\begin{lstlisting}
@ R1:R0 = a (lo,hi);  R3:R2 = b
    adds  r0, r0, r2                @ lo  + lo, sets C
    adc   r1, r1, r3                @ hi  + hi + C
    bx    lr
\end{lstlisting}
\cee{int64\_t add64(int64\_t a,int64\_t b)\{ return a+b; \}}. EABI: 64b in (R0=lo,R1=hi). \asm{adds} \emph{must} have S so that \asm{adc} sees the carry. For subtract: \asm{subs;sbc} (sbc subtracts the borrow = $\overline{C}$).

\textbf{R4: 64-bit subtract.}
\begin{lstlisting}
    subs  r0, r0, r2                @ lo - lo, C=1 if NO borrow
    sbc   r1, r1, r3                @ hi - hi - !C
    bx    lr
\end{lstlisting}
\cee{int64\_t sub64(int64\_t a,int64\_t b)\{ return a-b; \}}. ARM C-flag inverted vs x86 borrow.

\textbf{R5: strlen (null-terminator search).}
\begin{lstlisting}
@ R0 = char *s; ret R0 = length
    mov   r1, r0
R5_lp:
    ldrb  r2, [r0], #1
    cbnz  r2, R5_lp                 @ keep going while *p != 0
    sub   r0, r0, r1                @ delta = (one past '\0')
    sub   r0, r0, #1                @ exclude '\0'
    bx    lr
\end{lstlisting}
\cee{size\_t strlen(const char *s)\{const char *t=s; while(*t)t++; return t-s;\}}. Note CBNZ only sees R0--R7 and a forward $\le$126 B branch.

\textbf{R6: byte-wise memcpy (post-index).}
\begin{lstlisting}
@ R0=dst, R1=src, R2=n; ret R0 = dst
    mov   r3, r0
    cbz   r2, R6_end
R6_lp:
    ldrb  r12, [r1], #1
    strb  r12, [r0], #1
    subs  r2, r2, #1
    bne   R6_lp
R6_end:
    mov   r0, r3
    bx    lr
\end{lstlisting}
\cee{void *memcpy(void *d,const void *s,size\_t n)\{...\}}. R12 is scratch (caller-saved) so no PUSH needed.

\textbf{R7: signed saturate to $[-N,N]$.}
\begin{lstlisting}
@ R0 = x, R1 = N
    rsb   r2, r1, #0                @ r2 = -N
    cmp   r0, r2
    it    lt
    movlt r0, r2                    @ if x < -N, x = -N
    cmp   r0, r1
    it    gt
    movgt r0, r1                    @ if x >  N, x =  N
    bx    lr
\end{lstlisting}
\cee{int sat(int x,int N)\{ if(x<-N)x=-N; if(x>N)x=N; return x; \}}. Equivalent: \cee{SSAT R0,\#bits,R0} when N is a power-of-2 minus 1.

\textbf{R8: bit-field extract via UBFX/SBFX.}
\begin{lstlisting}
@ R0 = packed; extract bits[10:7] uns -> R1, bits[6:0] sgn -> R2
    ubfx  r1, r0, #7, #4            @ width 4 starting bit 7
    sbfx  r2, r0, #0, #7            @ signed 7-bit field
    bx    lr
\end{lstlisting}
\cee{r1 = (packed>>7) \& 0xF;} \cee{r2 = ((int32\_t)(packed<<25))>>25;} -- one instruction replaces 2 shifts + 1 AND.

\textbf{R9: switch via TBB jump table.}
\begin{lstlisting}
@ R0 = case index in [0..3]
    cmp   r0, #3
    bhi   def
    tbb   [pc, r0]                  @ byte table at NEXT instr
br_tbl:
    .byte (c0-br_tbl)/2
    .byte (c1-br_tbl)/2
    .byte (c2-br_tbl)/2
    .byte (c3-br_tbl)/2
    .align 1
c0: ...; b end
c1: ...; b end
c2: ...; b end
c3: ...; b end
def: ...
end: bx lr
\end{lstlisting}
\cee{switch(i)\{case 0:...; case 1:...; ...\}}. TBB looks up byte at \cee{table[r0]}, multiplies by 2, branches forward. TBH = halfword for far targets.

\subsection{C $\to$ ASM drills}

\textbf{C1: 64-bit return convention.}
\begin{lstlisting}[language={}]
int64_t add64(int64_t a, int64_t b) { return a + b; }
\end{lstlisting}
\begin{lstlisting}
@ a in R1:R0 (hi:lo), b in R3:R2; ret R1:R0
    adds r0, r0, r2
    adc  r1, r1, r3
    bx   lr
\end{lstlisting}
Both 64-bit args occupy 2 regs (lo,hi). On stack-spilled long longs, alignment is 8B; first stack arg may force a skip slot.

\textbf{C2: branchless abs (IT vs branch).}
\begin{lstlisting}[language={}]
int abs(int x) { return x<0 ? -x : x; }
\end{lstlisting}
\begin{lstlisting}
@ IT version (no branch, no flush)
    cmp   r0, #0
    it    lt
    rsblt r0, r0, #0                @ if x<0, x = 0 - x
    bx    lr
@ Branch version (older / GCC -O0)
    cmp   r0, #0
    bge   1f
    rsb   r0, r0, #0
1:  bx    lr
\end{lstlisting}
IT is faster on M4 (no pipeline flush); compiler picks it at -O2.

\textbf{C3: countdown delay using CBZ.}
\begin{lstlisting}[language={}]
void delay(uint32_t cycles){ while(cycles--) __asm__("nop"); }
\end{lstlisting}
\begin{lstlisting}
    cbz  r0, dly_end
dly_lp:
    nop
    subs r0, r0, #1
    bne  dly_lp
dly_end:
    bx   lr
\end{lstlisting}
CBZ saves the initial \cee{cmp} when zero is the early-out value.

\textbf{C4: register-offset array sum.}
\begin{lstlisting}[language={}]
int sum_array(int *a, int n) {
  int s=0; for(int i=0;i<n;i++) s+=a[i]; return s;
}
\end{lstlisting}
\begin{lstlisting}
    mov  r2, #0                     @ s
    mov  r3, #0                     @ i
sa_lp:
    cmp  r3, r1
    bge  sa_end
    ldr  r12, [r0, r3, lsl #2]      @ a[i]
    add  r2, r2, r12
    add  r3, r3, #1
    b    sa_lp
sa_end:
    mov  r0, r2
    bx   lr
\end{lstlisting}
Reg-offset \asm{[Rb,Ri,LSL\#2]} is one cycle, scales index by sizeof(int). Note: post-index pointer-bump variant from ARM ref is \emph{shorter} when \cee{i} isn't needed afterwards.

\textbf{C5: strcpy (returns dst).}
\begin{lstlisting}[language={}]
char *strcpy(char *d, const char *s){
  char *r=d; while((*d++=*s++)); return r;
}
\end{lstlisting}
\begin{lstlisting}
    mov  r2, r0
sc_lp:
    ldrb r3, [r1], #1
    strb r3, [r0], #1               @ writes terminator too
    cmp  r3, #0
    bne  sc_lp
    mov  r0, r2
    bx   lr
\end{lstlisting}

\textbf{C6: find\_max with conditional update.}
\begin{lstlisting}[language={}]
int find_max(int *a, int n){
  int m = a[0];
  for(int i=1;i<n;i++) if(a[i]>m) m=a[i];
  return m;
}
\end{lstlisting}
\begin{lstlisting}
    ldr  r2, [r0], #4               @ seed = a[0]
    sub  r1, r1, #1
fm_lp:
    cbz  r1, fm_end
    ldr  r3, [r0], #4
    cmp  r3, r2
    it   gt
    movgt r2, r3                    @ branchless max
    sub  r1, r1, #1
    b    fm_lp
fm_end:
    mov  r0, r2
    bx   lr
\end{lstlisting}

\subsection{Stack frame text-art (Full-Descending, 8B align at BL)}

\textbf{Leaf fn (no PUSH)} -- only LR \emph{in register}:
\begin{lstlisting}[language={}]
   higher addr
   ............
   |  caller's |
   |  frame    |
   +-----------+ <- SP at entry  (= SP at BX LR)
\end{lstlisting}
LR is live, never spilled. Cannot make a \asm{BL} (would clobber LR) without first saving it.

\textbf{Stem fn, \asm{push \{r4-r7,lr\}}} (5 regs $=$ 20\,B; padded to 24\,B for 8B align):
\begin{lstlisting}[language={}]
   +-----------+  caller frame
   |    LR     |  <- pushed last; HIGHEST addr of frame
   |    R7     |
   |    R6     |
   |    R5     |
   |    R4     |  <- LOWEST addr of saved set
   |  pad 4B   |  (inserted by GCC if needed for 8B align)
   +-----------+ <- SP after PUSH
\end{lstlisting}
\textbf{Rule:} ascending reg\# at ascending addr regardless of brace order. Total bytes pushed must keep SP $\equiv 0\pmod 8$ at any subsequent \asm{BL}.

\textbf{Stem with stack locals (\asm{sub sp,\#16}):}
\begin{lstlisting}[language={}]
   +-----------+
   |    LR     |
   |    R4     |
   +-----------+ <- SP after push {r4,lr}  (8B-aligned: 2 regs)
   |  loc[12]  |  local int a
   |  loc[8]   |  local int b
   |  loc[4]   |  local int c
   |  loc[0]   |  spill / temp     <- SP after sub sp,#16
   +-----------+
\end{lstlisting}
Access locals via \asm{[sp,\#0]}\dots\asm{[sp,\#12]}. Epilogue: \asm{add sp,\#16; pop \{r4,pc\}}.

\textbf{Inside ISR -- HW auto-stacks 8 words (R0-R3, R12, LR, retPC, xPSR):}
\begin{lstlisting}[language={}]
   +-----------+ <- pre-IRQ SP (8B aligned)
   |   xPSR    |  +0x1C
   |   PC      |  +0x18  (return address)
   |   LR      |  +0x14
   |   R12     |  +0x10
   |   R3      |  +0x0C
   |   R2      |  +0x08
   |   R1      |  +0x04
   |   R0      |  +0x00  <- SP at handler entry
   +-----------+
\end{lstlisting}
LR loaded with EXC\_RETURN (\cee{0xFFFFFFFD} thread/PSP, \cee{0xFFFFFFF9} thread/MSP, \cee{0xFFFFFFF1} handler). \asm{BX LR} pops the frame in HW. Tail-chain reuses this frame; late-arrival fixes up a higher-prio IRQ before unstacking. With FPU active, 18 extra words also stacked (s0--s15, FPSCR + pad).

\subsection{Cortex-M4 quick specs (extended)}
\begin{tabular}{@{}ll@{}}
\toprule
Item & Value \\
\midrule
Word size           & 32-bit (regs, ALU, addr bus) \\
ISA                 & Thumb-2 only (mix of 16/32-bit encodings) \\
Endianness          & Little (default; BE-8 optional, fixed at reset) \\
Pipeline            & 3-stage: Fetch / Decode / Execute \\
GP registers        & 16 (R0--R15); R13/14/15 = SP/LR/PC \\
SP banks            & MSP (default, handler), PSP (optional thread) \\
Multiplier          & 32$\times$32$\to$64 (1c MUL, 3--7c MULL) \\
Divider             & SDIV/UDIV, 2--12c, no flag update \\
DSP ext (M4)        & SIMD (UADD8, QADD16, SMLAD\dots), saturating arith \\
FPU (M4F)           & Single-precision IEEE 754, 32 S-regs / 16 D-regs \\
Bit-banding         & SRAM (0x2000\_0000) + Periph (0x4000\_0000), 1\,MB win \\
SysTick             & 24-bit downcount \\
NVIC IRQ slots      & up to 240 ext IRQs + 16 sys exceptions \\
NVIC priority bits  & 3--8 implemented; STM32 = 4 (16 levels) \\
Exc latency (no FP) & 12 cyc enter, 10 cyc exit (zero wait state) \\
Tail-chain latency  & 6 cyc (skips push+pop) \\
Late-arrival        & up to 6 cyc to swap higher-prio in mid-entry \\
Sleep modes         & WFI, WFE, sleep-on-exit \\
\bottomrule
\end{tabular}

\textbf{Special-purpose registers (CONTROL/PRIMASK/etc.):}
\begin{tabular}{@{}ll@{}}
\toprule
Reg & Effect \\
\midrule
\cee{PRIMASK}  & bit0=1 $\Rightarrow$ block all IRQs (NMI/HardFault still run); set/clr via \asm{cpsid i}/\asm{cpsie i} \\
\cee{FAULTMASK}& bit0=1 $\Rightarrow$ block all incl.\ HardFault (auto-cleared on return) \\
\cee{BASEPRI}  & 8-bit; if non-zero, mask IRQs whose priority $\ge$ BASEPRI value (selective gate) \\
\cee{CONTROL}  & b0=nPRIV (1=unpriv thread), b1=SPSEL (1=PSP), b2=FPCA (FP context active) \\
\cee{xPSR}     & APSR (NZCVQ) | IPSR (ExcNum 0..511) | EPSR (T-bit, IT-state, ICI) \\
\bottomrule
\end{tabular}
\textbf{IPSR = 16 + IRQn} for external IRQs; 0 in thread mode; 1=Reset, 2=NMI, 3=HardFault, \dots, 11=SVC, 14=PendSV, 15=SysTick.

\subsection{Vector table (first 16 system entries)}
Located in \cee{.isr\_vector} at start of FLASH (alias 0x0000\_0000 after reset). Each entry is a 32-bit address.
\begin{tabular}{@{}rlll@{}}
\toprule
Idx & Off & Name & Use \\
\midrule
0  & 0x00 & \cee{\_estack}      & Initial MSP value \\
1  & 0x04 & Reset\_Handler      & Entry after reset/POR \\
2  & 0x08 & NMI\_Handler        & Non-maskable interrupt (clock fail) \\
3  & 0x0C & HardFault\_Handler  & Catch-all fault if specific handlers off \\
4  & 0x10 & MemManage\_Handler  & MPU violation \\
5  & 0x14 & BusFault\_Handler   & Bus error (illegal access, abort) \\
6  & 0x18 & UsageFault\_Handler & Bad instr, unaligned, divide-by-0 \\
7  & 0x1C & ---                 & Reserved \\
8  & 0x20 & ---                 & Reserved \\
9  & 0x24 & ---                 & Reserved \\
10 & 0x28 & ---                 & Reserved \\
11 & 0x2C & SVC\_Handler        & Supervisor call (\asm{svc \#n}) \\
12 & 0x30 & DebugMon\_Handler   & Debug monitor \\
13 & 0x34 & ---                 & Reserved \\
14 & 0x38 & PendSV\_Handler     & SW-pendable, low-prio context switch \\
15 & 0x3C & SysTick\_Handler    & 24-bit tick timer \\
16+ & 0x40+ & IRQ0..IRQn        & Vendor peripheral handlers \\
\bottomrule
\end{tabular}
\textbf{Offset of vendor IRQ$n$:} \cee{0x40 + 4*n}. Vendor table is \cee{.weak} aliased to \cee{Default\_Handler} so user can override by defining the symbol. Relocate via \cee{SCB->VTOR = newbase;} (32-word align, $\ge$\,128 B aligned, point to RAM for runtime patching).

\subsection{Reset\_Handler ASM sketch}
\begin{lstlisting}
Reset_Handler:
  /* (0) MSP already loaded by HW from word 0 of vec table */
  ldr   r0, =_estack
  msr   msp, r0                     @ explicit, in case VTOR moved

  /* (1) copy .data init image FLASH -> SRAM */
  ldr   r0, =_sdata                 @ dest start  (RAM)
  ldr   r1, =_edata                 @ dest end
  ldr   r2, =_sidata                @ src  (FLASH image)
copy_loop:
  cmp   r0, r1
  ittt  lt
  ldrlt r3, [r2], #4
  strlt r3, [r0], #4
  blt   copy_loop

  /* (2) zero .bss */
  ldr   r0, =_sbss
  ldr   r1, =_ebss
  movs  r3, #0
zero_loop:
  cmp   r0, r1
  itt   lt
  strlt r3, [r0], #4
  blt   zero_loop

  /* (3) early SoC bring-up: clocks/PLL, FPU, VTOR */
  bl    SystemInit

  /* (4) C runtime init: static C++ ctors + libc */
  bl    __libc_init_array

  /* (5) jump to user main */
  bl    main

  /* (6) main() should not return; trap if it does */
hang:
  b     hang
\end{lstlisting}

\subsection{Linker script sketch}
\begin{lstlisting}[language={}]
MEMORY {
  FLASH (rx)  : ORIGIN = 0x08000000, LENGTH = 256K
  RAM   (xrw) : ORIGIN = 0x20000000, LENGTH =  32K
}
_estack = ORIGIN(RAM) + LENGTH(RAM);
SECTIONS {
  .isr_vector : { . = ALIGN(4); KEEP(*(.isr_vector)) } >FLASH
  .text       : { *(.text*) *(.rodata*) _etext = .; } >FLASH
  _sidata = LOADADDR(.data);
  .data : { _sdata = .; *(.data*) _edata = .; } >RAM AT>FLASH
  .bss  : { _sbss  = .; *(.bss*) *(COMMON) _ebss = .; } >RAM
}
\end{lstlisting}
\textbf{Key symbols consumed by Reset\_Handler:} \cee{\_estack, \_sdata, \_edata, \_sidata, \_sbss, \_ebss}. \cee{KEEP()} prevents \cee{--gc-sections} from dropping the vector table even though no C code references it directly.

\subsection{Exam-day error hit-list (extended)}
\begin{itemize}
\item \textbf{Forget to clear \cee{EXTI->PR} (W1C).} Handler returns, IRQ refires immediately $\to$ infinite loop. Cure: \cee{EXTI->PR = (1U<<line);} (write 1, NOT 0) at top of handler.
\item \textbf{Wrong AIRCR PRIGROUP.} Set via VECTKEY=\cee{0x05FA} + PRIGROUP[10:8]; mismatch with your \cee{NVIC\_SetPriority} expectations $\Rightarrow$ preempt vs sub split is wrong, ISRs preempt or fail to preempt unexpectedly.
\item \textbf{PUSH/POP register count mismatch.} \asm{push \{r4,r5,lr\}}\,/\,\asm{pop \{r4,pc\}} loses R5; subsequent reads of R5 see ISR junk or compiler-spilled value.
\item \textbf{B vs BL inside a function.} Plain \asm{B target} is a goto; doesn't update LR. If \cee{target} is meant to be called and return, you need \asm{BL}; otherwise you fall off the function.
\item \textbf{Leaf calling a non-leaf without saving LR.} If a "leaf" issues \asm{BL}, LR is overwritten, original return addr lost. Either turn it into a stem (\asm{push \{lr\}}) or inline the callee.
\item \textbf{Misaligned SP at \asm{BL}.} AAPCS requires 8B align at every external call. \asm{push \{lr\}} alone leaves SP off by 4; pair with one more reg (\asm{push \{r4,lr\}}) or pad with \asm{sub sp,\#4}.
\item \textbf{Reading \cee{IDR} for an output pin.} Works on M4 (the line is sampled in), but the \emph{intent} is the latched value; read \cee{ODR} for what was last written.
\item \textbf{Forgetting \cee{volatile} on a HW pointer.} \cee{uint32\_t *p = (uint32\_t*)0x40020000;} -- compiler caches/elides accesses. Always \cee{volatile uint32\_t *}.
\item \textbf{\cee{++} on \cee{*p} without write-back.} \cee{++*p} compiles to LDR; ADD; \emph{STR}. Skipping the STR (or writing inline asm that omits it) leaves the bumped value in a register only.
\item \textbf{\asm{LDRSH} on misaligned addr.} HW requires bit0=0 for \asm{LDRH/SH}. \asm{[Ra,\#1]} where \cee{Ra} is word-aligned $\to$ effective addr odd $\to$ UsageFault (UNALIGN\_TRP=1) or implementation-defined garbage.
\item \textbf{Inverting \asm{Bcc} family for signedness.} \asm{BGT/BLT} are signed; \asm{BHI/BLO} are unsigned. Mixing breaks at boundary (e.g.\ \asm{BGT} on \cee{uint32\_t} compares treats $0x80000000$ as negative).
\item \textbf{CBZ/CBNZ inside an IT block.} Hard assemble error -- they don't read flags so IT can't gate them. Also: only \emph{conditional B} is legal as the FINAL slot of an IT.
\item \textbf{\asm{cmp r0,\#imm} with imm not in modimm encoding.} Modified-imm = 8-bit value rotated by even amount, OR specific 32b patterns. \asm{cmp r0,\#0x12345} fails to assemble; use \asm{ldr r1,=0x12345; cmp r0,r1}.
\item \textbf{\asm{cmp Rn,\#-imm} silently re-encoded as CMN.} Same flag effect for signed CCS, but disassembly reads \asm{CMN} which can confuse hand-traced flag analysis.
\item \textbf{Reusing a clobbered LR.} After \asm{BL foo}, LR holds the return-from-foo address. Using \asm{BX LR} \emph{after} another \asm{BL} returns to the wrong place.
\item \textbf{Forgetting to enable peripheral clock.} \cee{RCC->AHB2ENR |= GPIOAEN;} required before any GPIOA register write -- silent failure (writes ignored, reads return 0).
\item \textbf{Endianness on union/cast.} \cee{union\{u32 w; u8 b[4];\} u; u.w=0x11223344;} $\Rightarrow$ \cee{u.b[0]==0x44} (LE). Don't assume MSB-first.
\item \textbf{Comparing signed and unsigned.} \cee{int x=-1; if(x < 1U)} is FALSE because $x$ promotes to \cee{0xFFFFFFFF}. C99 6.3.1.8.
\item \textbf{\asm{adds} omitted before \asm{adc}.} 64-bit chain breaks because C is stale. Same for \asm{subs;sbc}. Always pair the S-suffixed step.
\item \textbf{Trying TBB/TBH with an out-of-range index.} No bounds check -- guard with \asm{cmp Rn,\#max; bhi default} BEFORE the TBB.
\end{itemize}

\subsection{HAL vs direct register}
\begin{tabular}{@{}lll@{}}
\toprule
Aspect & HAL & Direct (CMSIS bare metal) \\
\midrule
Naming   & \cee{HAL\_GPIO\_WritePin}, \cee{HAL\_UART\_Transmit} & \cee{GPIOA->BSRR=...}, \cee{USART2->TDR=...} \\
Coverage & Drivers, error-checks, callbacks & Direct register access only \\
Speed    & 10--100$\times$ slower (struct deref + checks) & 1--3 instr per access \\
Setup    & CubeMX-generated init                     & Manual RCC + register writes \\
Best for & App code, complex periph (USB, SD), portability & ISRs, time-critical loops, exam Q's \\
\bottomrule
\end{tabular}
\textbf{Exam pattern:} questions almost always require \emph{direct register} answers -- e.g.\ "set PA5 high atomically" $\to$ \cee{GPIOA->BSRR = (1U<<5);} (NOT \cee{HAL\_GPIO\_WritePin(GPIOA,GPIO\_PIN\_5,GPIO\_PIN\_SET);}). Naming hint: HAL fns are \cee{HAL\_<MOD>\_<Verb>(...)}; CMSIS uses \cee{<PERI>-><REG>}.

\subsection{NVIC priority math (PRIGROUP $\to$ preempt vs sub)}
8-bit priority register; only top $N$ bits implemented (STM32 = 4, so low 4 read 0). PRIGROUP value $G\in\{0..7\}$ specifies the split point:
\[
  \begin{aligned}
    &\text{preempt bits} = 7-G \quad (\text{within top }N=4)\\
    &\text{sub bits}     = G   \quad (\text{within top }N=4)\\
    &\text{stored byte}  = (\text{PPN} \ll (G+1)\;|\;\text{SPN}) \ll (8-N)
  \end{aligned}
\]
For STM32 ($N=4$, low 4 bits don't exist):
\begin{tabular}{@{}cccl@{}}
\toprule
PRIGROUP $G$ & PP bits & SP bits & Layout in IP[IRQn] (top 4) \\
\midrule
3 & 4 & 0 & PPPP\,---- (16 levels preempt, 1 sub) \\
4 & 3 & 1 & PPPS\,---- (8 preempt, 2 sub) \\
5 & 2 & 2 & PPSS\,---- (4 preempt, 4 sub) \\
6 & 1 & 3 & PSSS\,---- (2 preempt, 8 sub) \\
7 & 0 & 4 & SSSS\,---- (no preempt, 16 sub) \\
\bottomrule
\end{tabular}

\textbf{Worked examples (decode \cee{IP[IRQn]} byte $\to$ PPN/SPN):}

\begin{ex}\textbf{$G$=4, IP byte \cee{0x60}.} Top 4 = \cee{0110}. Layout PPPS: PPN = \cee{011}b = 3, SPN = \cee{0}b = 0. Total PN $= (3{<<}1)|0 = 6$.\end{ex}

\begin{ex}\textbf{$G$=5, IP byte \cee{0xB0}.} Top 4 = \cee{1011}. Layout PPSS: PPN = \cee{10}b = 2, SPN = \cee{11}b = 3. PN $= (2{<<}2)|3 = 11$.\end{ex}

\begin{ex}\textbf{$G$=3, IP byte \cee{0xA0}.} Top 4 = \cee{1010}. All 4 bits are preempt: PPN = 10, SPN = 0. Same-PPN ISRs cannot preempt each other.\end{ex}

\begin{ex}\textbf{Encode: $G$=4, want PPN=2, SPN=1.} Field = \cee{010\,1} = \cee{0101}b. Place in top 4: \cee{IP[IRQn] = 0x50;}. Equivalent: \cee{NVIC\_SetPriority(IRQn, NVIC\_EncodePriority(4,2,1));}.\end{ex}

\textbf{Preemption rule:} an IRQ \emph{can} preempt the running ISR iff its \emph{PPN} is strictly less. Equal PPN $\to$ pending; tie broken by lower IRQn first. Sub-priority only orders \emph{simultaneously pending} same-PPN IRQs at dispatch -- never preempts.

\subsection{Useful macros + intrinsics (CMSIS / HAL)}
\begin{tabular}{@{}ll@{}}
\toprule
Symbol & Effect \\
\midrule
\cee{\_\_HAL\_RCC\_GPIOA\_CLK\_ENABLE()} & set RCC AHB2ENR.GPIOAEN bit \\
\cee{\_\_HAL\_RCC\_USART2\_CLK\_ENABLE()} & set RCC APB1ENR1.USART2EN \\
\cee{\_\_HAL\_RCC\_SYSCFG\_CLK\_ENABLE()} & needed before EXTICR write \\
\cee{NVIC\_EnableIRQ(IRQn)}            & ISER[IRQn/32] $|=$ 1$<<$(IRQn\%32) \\
\cee{NVIC\_DisableIRQ(IRQn)}           & ICER set bit \\
\cee{NVIC\_SetPriority(IRQn,p)}        & writes IP[IRQn] = p$<<$(8$-$N) \\
\cee{NVIC\_SetPriorityGrouping(g)}     & AIRCR PRIGROUP=g + VECTKEY \\
\cee{NVIC\_ClearPendingIRQ(IRQn)}      & ICPR set bit \\
\cee{\_\_disable\_irq() / \_\_enable\_irq()} & \asm{cpsid i} / \asm{cpsie i} (PRIMASK) \\
\cee{\_\_set\_BASEPRI(p)}              & selective IRQ mask above level $p$ \\
\cee{\_\_DSB() / \_\_ISB() / \_\_DMB()}& barriers (data sync / instr sync / mem) \\
\cee{\_\_NOP() / \_\_WFI() / \_\_WFE()}& single-instr intrinsics \\
\cee{\_\_REV(x) / \_\_REV16 / \_\_RBIT}& byte-swap / hw-swap / bit-reverse \\
\cee{\_\_CLZ(x)}                       & count leading zeros (1 cyc) \\
\cee{\_\_LDREXW / \_\_STREXW}          & exclusive monitor for atomics \\
\cee{SET\_BIT / CLEAR\_BIT / READ\_BIT}& \cee{REG |= MASK} / \cee{\&= \~{}MASK} / \cee{REG \& MASK} \\
\cee{MODIFY\_REG(R,M,V)}               & \cee{R = (R \& \~{}M) | (V \& M)} \\
\bottomrule
\end{tabular}

\begin{ex}\textbf{Atomic NVIC enable USART2 (IRQn=38).} \cee{NVIC->ISER[1] = 1U << 6;} or \cee{NVIC\_EnableIRQ(USART2\_IRQn);}. Both compile to a single 32-bit STR -- ISER is W1S so no RMW.\end{ex}

\begin{ex}\textbf{Disable IRQs around 64-bit RMW.} \cee{uint32\_t pri = \_\_get\_PRIMASK(); \_\_disable\_irq(); /*update*/ \_\_set\_PRIMASK(pri);} -- preserves nesting (don't blindly \cee{\_\_enable\_irq()}).\end{ex}

\begin{ex}\textbf{Read\,/modify\,/write that should be atomic.} \cee{REG |= mask;} compiles to LDR; ORR; STR. ISR firing between LDR and STR loses ISR's update. Use BSRR-style W1S/W1C registers when available, otherwise mask IRQs.\end{ex}

\subsection{One-page summary block (System internals)}
\textbf{Vector table:} 32-bit entries; index 0=\cee{\_estack}, 1=Reset, 2--15 system, 16+ vendor. Offset = $4\cdot$idx. Reset\_Handler: load MSP $\to$ copy .data $\to$ zero .bss $\to$ SystemInit $\to$ \_\_libc\_init\_array $\to$ main $\to$ trap loop. \textbf{Stack frame:} HW auto-stacks 8 words on IRQ entry (R0--R3, R12, LR, PC, xPSR), pushed in ascending-addr order; SW prologue may push R4--R11+LR for stem fns. Always 8B-aligned at \asm{BL}. \textbf{NVIC priority:} byte per IRQ, top $N=4$ bits used on STM32; PRIGROUP $G$ splits top 4 into PPN (high) + SPN (low) where PP bits = $7-G$ within those 4. Preempt iff PPN strictly less. \textbf{System regs:} PRIMASK gates all (cpsid i), FAULTMASK gates incl.\ HardFault, BASEPRI selective gate, CONTROL holds nPRIV/SPSEL/FPCA. \textbf{Exam tells:} HAL questions ask for register names, not function calls; clear PR (W1C) in EXTI handler; pair \asm{adds}+\asm{adc} for 64b; ARM C-flag inverted on subtract; CBZ no-IT and forward-only; PUSH order always ascending; 8B align at every BL. \textbf{Reverse engineer:} look for post-index (\asm{[Rb],\#k}) $\to$ \cee{*p++}; pre-index (\asm{[Rb,\#k]!}) $\to$ \cee{*++p}; reg-offset \asm{[Rb,Ri,LSL\#n]} $\to$ \cee{a[i]} where $n$=log2 sizeof; CMP+IT+MOVcc $\to$ ternary; CBZ on a freshly LDRB'd byte $\to$ string null check; \asm{adds;adc} $\to$ 64-bit; \asm{ldr;add\#1;str} on same addr $\to$ \cee{(*p)++}. \textbf{Translate forward:} type chooses LDR suffix (sgn $\Rightarrow$ S-variant); array stride chooses LSL\#n; loop shape chooses CBZ vs cmp+B; conditional update chooses IT block; 64-bit return uses (R0,R1)=(lo,hi). \textbf{Final sanity:} flag-setting needs S-suffix; CMP/TST/TEQ/CMN already do; SDIV/UDIV never do; modimm has 8b+rotate constraint; literal pool reachable within 4\,KB.

% ----- part2/12_quiz6a_practice.tex -----
\section{Quiz 6a Practice (qz-25b, ptr + flow ctrl)}

\textbf{Note:} this is the \emph{Fall-2025} (qz-25b) qz6a worksheet --- the
current-semester qz6a is not in the repo, so use this version as the
canonical practice. Worksheet only, no instructor solution PDF; answers
below derived from first principles via HW9--12.

\subsection{P1 (45) --- LDRD/STRD + LDRSB pre-idx copy}
\begin{qp}
\cee{uint8\_t ipt[20]} with \cee{ipt[i]=0x70+(i<<2)};
\cee{int32\_t opt[4]}; \asm{task\_a(ipt,opt)}.
\begin{lstlisting}
task_a:
  ldrd  r2,r3,[r0,#8]!   @ pre-idx, then load 2 words
  strd  r2,r3,[r1,#0]    @ opt[0], opt[1]
  ldrsb r2,[r0,#4]!      @ pre-idx byte, sign-ext
  ldrsb r3,[r0,#4]!
  strd  r2,r3,[r1,#8]    @ opt[2], opt[3]
  bx    lr
\end{lstlisting}
\end{qp}

\textbf{(a) Memory map of \cee{ipt[20]}.}
$0x70+(i\!\ll\!2) = 0x70+4i$ \;$\Rightarrow$ stride$=4$, byte values:
\begin{tabular}{@{}ll|ll|ll@{}}
\toprule
i & ipt & i & ipt & i & ipt \\
\midrule
0  & 0x70 & 7  & 0x8C & 14 & 0xA8 \\
1  & 0x74 & 8  & 0x90 & 15 & 0xAC \\
2  & 0x78 & 9  & 0x94 & 16 & 0xB0 \\
3  & 0x7C & 10 & 0x98 & 17 & 0xB4 \\
4  & 0x80 & 11 & 0x9C & 18 & 0xB8 \\
5  & 0x84 & 12 & 0xA0 & 19 & 0xBC \\
6  & 0x88 & 13 & 0xA4 &    &      \\
\bottomrule
\end{tabular}
\textbf{Trap:} stride is 4 \emph{bytes}, not 4 indices --- \cee{ipt[i]} is
still a single byte, neighbors differ by $0x04$.

\textbf{(b) Trace.} R0 in $=$ \cee{\&ipt[0]}, R1 in $=$ \cee{\&opt[0]}.
\begin{tabular}{@{}lll@{}}
\toprule
step & R0 after & loaded \\
\midrule
\asm{ldrd r2,r3,[r0,\#8]!} & \cee{\&ipt[8]}  & R2,R3 \\
\asm{ldrsb r2,[r0,\#4]!}   & \cee{\&ipt[12]} & R2 \\
\asm{ldrsb r3,[r0,\#4]!}   & \cee{\&ipt[16]} & R3 \\
\bottomrule
\end{tabular}
\textbf{LDRD words} (LE pack of 4 bytes at \cee{\&ipt[8]}, \cee{\&ipt[12]}):
R2 $=$ bytes \cee{\{90,94,98,9C\}} $\to$ \asm{0x9C989490};\;
R3 $=$ bytes \cee{\{A0,A4,A8,AC\}} $\to$ \asm{0xACA8A4A0}.\\
\textbf{LDRSB:} byte at \cee{\&ipt[12]}$=0xA0$ (MSB$=1$) $\to$
R2 $=$ \asm{0xFFFFFFA0}.\;
Byte at \cee{\&ipt[16]}$=0xB0$ $\to$ R3 $=$ \asm{0xFFFFFFB0}.

\textbf{Printout:}\\
\asm{opt[0]=0x9C989490, opt[1]=0xACA8A4A0,}\\
\asm{opt[2]=0xFFFFFFA0, opt[3]=0xFFFFFFB0}

\textbf{(c) Exit registers.}
R0 $=$ \cee{\&ipt[16]} (last pre-idx update; R0 advanced
$8+4+4=16$ bytes).\;
R1 $=$ \cee{\&opt[0]} unchanged --- both \asm{strd}'s used basic
offset (no \asm{!}).

\textbf{(d) C translation.}
\begin{lstlisting}[language=C]
void task_a(uint8_t *p, int32_t *q) {
  int32_t *pw = (int32_t*)(p + 8);   // ldrd from ipt+8
  q[0] = pw[0];                       // word @ ipt[8..11]
  q[1] = pw[1];                       // word @ ipt[12..15]
  int8_t *pb = (int8_t*)(p + 16);     // ldrsb from ipt+16
  q[2] = (int32_t)pb[0];              // sign-ext ipt[12]? NO --
  q[3] = (int32_t)pb[1];              // see below
}
\end{lstlisting}
\textbf{Careful:} \asm{ldrsb r2,[r0,\#4]!} \emph{after} ldrd-pre means
R0 was already at \cee{ipt+8}, so first ldrsb reads
\cee{ipt[12]}, second reads \cee{ipt[16]}.
\begin{lstlisting}[language=C]
// Cleanest form, follows ASM byte-by-byte:
q[0] = *(int32_t*)(p+8);
q[1] = *(int32_t*)(p+12);
q[2] = (int32_t)(int8_t)p[12];
q[3] = (int32_t)(int8_t)p[16];
\end{lstlisting}

\subsection{P2 (20) --- Saturate to $[-5,5]$}
\begin{qp}
\cee{int math\_fun\_c(int x)\{ if(x<-5) x=-5; else if(x>5) x=5; return x; \}}
Translate to ASM \asm{int math\_fun\_s(int x)}.
\end{qp}

\textbf{Branch form (cleanest):}
\begin{lstlisting}
math_fun_s:                @ x->R0, ret R0
  cmp   r0, #-5
  bge   .Lchk_hi           @ x>=-5 -> check high
  mov   r0, #-5
  bx    lr
.Lchk_hi:
  cmp   r0, #5
  ble   .Ldone             @ x<=5 -> already ok
  mov   r0, #5
.Ldone:
  bx    lr
\end{lstlisting}

\textbf{IT-block form (denser, no labels):}
\begin{lstlisting}
math_fun_s:
  cmp   r0, #-5
  it    lt
  movlt r0, #-5            @ if x<-5 then x=-5
  cmp   r0, #5
  it    gt
  movgt r0, #5             @ else if x>5 then x=5
  bx    lr
\end{lstlisting}
\textbf{Why \asm{blt}/\asm{bgt}/\asm{lt}/\asm{gt}?} \cee{int} is signed
$\Rightarrow$ signed cond codes (\asm{lt/gt/le/ge}). Using
\asm{blo/bhi} would \emph{wrongly} treat $-5$ as a huge unsigned value
$0xFFFFFFFB$.\;
\textbf{2nd cmp safe even after clamp-low:} after \asm{movlt r0,\#-5},
$x=-5<5$ so the \asm{movgt} is correctly skipped --- no need for
explicit \asm{else if} branching.

\subsection{P3 (35) --- Sum of negative elements}
\begin{qp}
\cee{int mp\_neg\_sum\_array(int *arr, int N)} returning sum of
\cee{arr[i]} for which \cee{arr[i]<0}. Write C (15) and ASM (20).
\end{qp}

\textbf{C (15 pts):}
\begin{lstlisting}[language=C]
int mp_neg_sum_array_c(int *arr, int N) {
  int sum = 0;
  for (int i = 0; i < N; i++) {
    if (arr[i] < 0) sum += arr[i];
  }
  return sum;
}
\end{lstlisting}

\textbf{ASM (20 pts) --- index + scaled load:}
\begin{lstlisting}
mp_neg_sum_array_s:        @ R0=arr, R1=N
  mov   r2, #0             @ sum = 0
  mov   r3, #0             @ i   = 0
.Lloop:
  cmp   r3, r1
  bge   .Ldone             @ i>=N -> exit
  ldr   r12,[r0,r3,lsl #2] @ tmp = arr[i]
  cmp   r12,#0
  bge   .Lskip             @ tmp>=0 -> skip
  add   r2, r2, r12        @ sum += tmp
.Lskip:
  add   r3, r3, #1         @ i++
  b     .Lloop
.Ldone:
  mov   r0, r2             @ return sum
  bx    lr
\end{lstlisting}

\textbf{IT-block variant} (saves the \asm{.Lskip} label):
\begin{lstlisting}
.Lloop:
  cmp   r3, r1
  bge   .Ldone
  ldr   r12,[r0,r3,lsl #2]
  cmp   r12, #0
  it    lt
  addlt r2, r2, r12        @ sum += arr[i] if <0
  add   r3, r3, #1
  b     .Lloop
\end{lstlisting}

\textbf{Pointer-walk variant} (no scaled index):
\begin{lstlisting}
  mov   r2, #0             @ sum
  add   r1, r0, r1, lsl #2 @ end = arr + N
.L:
  cmp   r0, r1
  bge   .D
  ldr   r12,[r0],#4        @ tmp=*p++; post-idx
  cmp   r12, #0
  it    lt
  addlt r2, r2, r12
  b     .L
.D:
  mov   r0, r2
  bx    lr
\end{lstlisting}

\textbf{Reg map / pattern.}
\begin{tabular}{@{}lll@{}}
\toprule
role     & reg  & note \\
\midrule
arg1 \cee{arr} & R0   & also caller-save scratch \\
arg2 \cee{N}   & R1   & loop bound        \\
sum            & R2   & final $\to$ R0    \\
i              & R3   & loop counter      \\
tmp \cee{arr[i]}& R12 & \asm{ip}, scratch \\
\bottomrule
\end{tabular}
R0--R3, R12 are caller-save (AAPCS) $\Rightarrow$ no push/pop needed.\;
\textbf{Index vs ptr-walk:} index form = 1 ldr/iter w/ \asm{lsl \#2};
ptr-walk = 1 ldr post-idx --- both 4 cycles in inner loop, choice is
style.

\subsection{Cross-cutting traps (qz6a checklist)}
\begin{itemize}\setlength{\itemsep}{0pt}
\item \textbf{LDRSB on $\geq$0x80} $\to$ top 24 bits become \asm{1};
      \cee{ipt[12]=0xA0}, \cee{ipt[16]=0xB0} both sign-extend.
\item \textbf{LDRD/STRD pair} $=$ consecutive registers (R2,R3 here);
      no register-offset form, but pre-/post-idx OK.
\item \textbf{Pre-idx \asm{!}} updates base \emph{before} access $=$
      \cee{*++p}; cumulative across multiple \asm{!} ops $\Rightarrow$
      total advance $= \sum$ offsets ($8+4+4=16$).
\item \textbf{STR/STRD with no \asm{!}} = base unchanged;
      R1 stays at \cee{\&opt[0]} throughout P1.
\item \textbf{Signed cmp} for \cee{int}: use
      \asm{blt/bge/bgt/ble} (\emph{not} \asm{blo/bhs/bhi/bls}).
\item \textbf{Array-loop template} (drill until automatic):
      init \asm{sum=0,i=0}; top: \asm{cmp i,N / bge done};
      \asm{ldr tmp,[base,i,LSL\#2]}; predicate; conditional accumulate;
      \asm{add i,\#1 / b loop}; done: \asm{mov r0,sum / bx lr}.
\item \textbf{IT block} useful when body is 1--4 instructions all
      sharing one cond code; replaces label+branch.
\item \textbf{C $\to$ ASM} for \cee{a[i]} on \cee{int*}: scaled
      \asm{[base, i, LSL \#2]} (no writeback) $=$ \cee{p[i]} idiom.
\end{itemize}

% ==============================================================================
% Common footer for both parts. Closes multicols + document.
% ==============================================================================
\end{multicols*}
\end{document}
